\documentclass[11pt,a4paper]{article}
\usepackage[T1]{fontenc}
\usepackage[utf8]{inputenc}
\usepackage{lmodern}
\usepackage[margin=26mm,headheight=14pt]{geometry}
\usepackage{amsmath,amssymb,amsthm,mathtools}
\usepackage{microtype}
\usepackage{booktabs,longtable,array,enumitem}
\usepackage{etoolbox}
% Finite-shrink bottom glue avoids longtable page-splitting diagnostics
% with recent LaTeX output routines; stretch still fills the same space.
\makeatletter
\patchcmd{\LT@output}{\vss}{\vfil}{}{}
\patchcmd{\LT@output}{\vss}{\vfil}{}{}
\makeatother
\usepackage{xurl}
\usepackage[dvipsnames]{xcolor}
\usepackage[hypertexnames=false,colorlinks=true,linkcolor=MidnightBlue,citecolor=MidnightBlue,urlcolor=MidnightBlue]{hyperref}
\usepackage{fancyhdr}
\pagestyle{fancy}
\fancyhf{}
\fancyhead[L]{\small Ginibre and Poisson Boolean percolation}
\fancyhead[R]{\small Research report}
\fancyfoot[C]{\thepage}
\setlength{\parindent}{0pt}
\setlength{\parskip}{5pt plus 1pt minus 1pt}
\setlength{\emergencystretch}{2em}
\setlength{\LTpre}{8pt}
\setlength{\LTpost}{8pt}
\newcolumntype{L}[1]{>{\raggedright\arraybackslash}p{#1}}
\setcounter{tocdepth}{2}
\allowdisplaybreaks[2]
\newtheorem{theorem}{Theorem}
\newtheorem{lemma}[theorem]{Lemma}
\newtheorem{proposition}[theorem]{Proposition}
\newtheorem{corollary}[theorem]{Corollary}
\theoremstyle{definition}
\newtheorem{definition}[theorem]{Definition}
\theoremstyle{remark}
\newtheorem{remark}[theorem]{Remark}
\numberwithin{equation}{section}
\providecommand{\tightlist}{\setlength{\itemsep}{0pt}\setlength{\parskip}{0pt}}
\DeclareMathOperator{\dist}{dist}
\DeclareMathOperator{\Tr}{Tr}
\DeclareMathOperator{\Var}{Var}
\DeclareMathOperator{\Cov}{Cov}
\DeclareMathOperator{\DPP}{DPP}
\hypersetup{pdftitle={Ginibre and Poisson Boolean percolation: investigation report},
 pdfauthor={Research record prepared with Codex},
 pdfsubject={Verified results, recovered arguments, and obstructions; original comparison unresolved}}
\begin{document}

% ===== BEGIN front =====
\begin{titlepage}
\centering
\vspace*{20mm}
{\Large Mathematical research report\par}
\vspace{9mm}
{\huge\bfseries Ginibre and Poisson Boolean percolation\par}
\vspace{6mm}
{\Large Proofs, recovered arguments, and obstructions\par}
\vspace{15mm}
{\large Investigation of 15 September 2026\\
Amended 21 September 2026\par}
\vfill
\begin{minipage}{0.88\textwidth}
\textbf{Status of the original target.}
The strict inequality
\[
             r_c(G)<r_c(\Pi_1)
\]
for equal unit intensity and deterministic disk grains remains unproved
and undisproved in this investigation.

\medskip
The critical-radius limit under area-preserving stretching is established
by the proof included here. The recovered small-\(\beta\) expansion is
preserved in detail, but its central uniform geometric estimate has not
received a complete fresh verification. Neither result is silently
identified with the original target.
\end{minipage}
\vfill
{\small A theoretical investigation; no numerical experiments.\par}
\end{titlepage}

\begin{abstract}
This report records the investigation of the critical-radius comparison
between planar Ginibre and Poisson Boolean models. It contains complete
proofs of an optimal homogeneous Poisson domination for thinned and
rescaled Ginibre, exact Palm interpolation and Fock variational identities,
buffered conditional Bergman domination, an obstruction to uniform exterior
Palm comparison, and finite-ensemble deletion calculations. It restores
the ellipse route: fixed-eccentricity large-area comparison, quantitative
local Poisson convergence under area-preserving stretching, and a
negative-association block argument proving convergence of the actual
critical radii. It also preserves the small-\(\beta\) asymptotic argument,
with its principal geometric estimate explicitly separated from its
verified algebraic and limiting consequences. The remaining routes are
organized by their mathematical outcomes and exact obstructions.
The amended account also restores the quantitative stretching rate,
shape-variation and switching calculations, coupling barriers, the
positive-bath mixed-model comparison, conditional clearing laws, and
further certificate, transport, and finite-rank refinements.
The full original notes accompany the report; historical claims of audit
success in those notes are not used as substitutes for the proofs here.
No claim of novelty or resolution of the original strict inequality is made.
\end{abstract}

\tableofcontents
\clearpage

\section*{Scope, conventions, and verification status}
\addcontentsline{toc}{section}{Scope, conventions, and verification status}

\subsection*{The mathematical target}
Identify \(\mathbb R^2\) with \(\mathbb C\), and write \(dA\) for Lebesgue
area. For a locally finite configuration \(\eta\), define
\[
 \mathcal O_r(\eta)=\bigcup_{x\in\eta}\overline B(x,r).
\]
Its critical radius is defined by
\[
 r_c(\mathcal L(\eta))=
 \inf\{r>0:\mathbb P(\mathcal O_r(\eta)
           \text{ has an unbounded component})>0\}.
\]
Every grain is a full closed disk. Thus the center connection range is
\(2r\), not \(r\). A restricted crossing means that its occupied path is
restricted to the observation rectangle; centers outside that rectangle
are still allowed when their disks enter it. Such an event is determined
by the radius-\(r\) expansion of the rectangle.

Let \(\Pi_\lambda\) be homogeneous Poisson of intensity \(\lambda\),
write \(\Pi=\Pi_1\), and put \(r_P=r_c(\Pi_1)\).
The unit-intensity Ginibre process \(G\) has kernel
\[
 H(z,w)=\exp\!\left(\pi z\overline w
                -\frac\pi2(|z|^2+|w|^2)\right),
 \qquad |H(z,w)|^2=e^{-\pi|z-w|^2}.
\]
Kernels are relative to \(dA\), unless a change of gauge is stated.
The \(n\)-point factorial density is \(\det[H(z_i,z_j)]\).
The kernel is an orthogonal projection in the physical Gaussian gauge.

For a finite distinct tuple \(A\), \(G^{!A}\) denotes its canonical
reduced Palm law. The prescribed points are omitted. This is defined
through factorial densities, not ordinary conditioning on a probability-zero
event. Whenever those points play the role of occupied parents or arms,
their restoration is explicitly part of the configuration. The corresponding
reduced kernel is \(Q_A=H-P_A\), where \(P_A\) projects onto the span of
the coherent evaluation vectors at \(A\).

For any functional \(F\), write
\[
 D_qF(\eta)=F(\eta\cup\{q\})-F(\eta).
\]
For increasing Boolean \(F\), the mixed difference \(D_qD_zF\) can have
either sign. It equals \(+1\) when two insertions cooperate and \(-1\)
when they are individually successful alternatives. This sign distinction
is central to the investigation.

Two deformations must be kept separate:
\begin{itemize}
\item \(G_\beta=\operatorname{DPP}(\beta H_\beta)\) is unit intensity
after thinning by \(\beta\) and spatial contraction by \(\sqrt\beta\).
The small-\(\beta\) endpoint is Poisson; the full Ginibre endpoint is
\(\beta=1\).
\item \(\Phi_t=S_t^{-1}G\), with
\(S_t(x,y)=(e^t x,e^{-t}y)\), has unit intensity because
\(\det S_t=1\). The parameter \(t\) changes anisotropy without thinning.
\end{itemize}
The notation \(G_s=\operatorname{DPP}(sH)\) used in the signed
interpolation section instead has intensity \(s\); a Poisson bath restores
total intensity one on the path \(\nu=1-s\).

\subsection*{How to read the status statements}
\begin{enumerate}[leftmargin=*]
\item \textbf{Verified in this report.} The displayed argument has been
reconstructed and checked, using the cited published inputs. This wording
is a statement about the supplied proof, not external certification or novelty.
\item \textbf{Recovered argument.} A substantive derivation survives in
the original notes, sometimes with older positive audit records. Its
complete dependency chain has not been freshly verified. It is preserved
with its exact unresolved verification dependency; it is not declared false.
\item \textbf{Disproved auxiliary assertion.} A counterexample or exact
calculation contradicts the assertion. The scope of that counterexample
is stated; failure of a stronger auxiliary assertion does not disprove
the original critical-radius comparison.
\end{enumerate}

\begin{longtable}{@{}L{0.36\textwidth}L{0.20\textwidth}L{0.37\textwidth}@{}}
\toprule
Result or route & Present status & Consequence and limitation\\
\midrule
\endhead
Optimal upper Poisson domination of \(G_\beta\)
& Verified & Gives a lower bound on \(r_c(G_\beta)\), for \(\beta<1\).\\[3pt]
Palm interpolation and Fock variational formulas
& Verified & Exact identities; their pivotal average has no established favorable sign.\\[3pt]
Buffered Bergman bound and exterior Palm likelihood obstruction
& Verified & Control conditioned remainders and rule out uniform arbitrary-exterior comparisons.\\[3pt]
Fixed-shape large-area ellipse comparison
& Verified & Strict comparison of large-grain percolation probabilities, away from an identified critical comparison.\\[3pt]
\(r_c(S_t^{-1}G)\to r_P\)
& Verified & Actual infinite-volume thresholds converge; no ordering of the value at \(t=0\) follows.\\[3pt]
Small-\(\beta\) expansion with coefficient \(-r_P/4\)
& Recovered argument & Depends on a uniform geometric pivotal estimate, including the zero-bath regime.\\[3pt]
Further certificate, motion, and multiple-mark constructions
& Status specified individually & Their source proofs are preserved; none supplies the missing full-Ginibre pivotal average.\\[3pt]
Quantitative shape rate, mixed-model gap, and additional local refinements
& Recovered arguments & Restored with source hypotheses, quantitative constants, and remaining dependencies.\\[3pt]
\(r_c(G)<r_P\)
& Unresolved & No positive radius gap at \(\beta=1\) is established.\\
\bottomrule
\end{longtable}

\subsection*{Completeness and provenance}
The supplied research-history PDF and ellipse note are retained unchanged
in the companion source archive. The archive also includes every surviving
individual mathematical Markdown note and the earlier consolidated checkpoint.
The report provides a coherent mathematical account rather than treating
successive checkpoint descriptions as independent proofs. Its source index
maps back to the original files. Earlier labels such as ``audited'' and
``proved'' are historical assertions when quoted in the supplement;
the status statements in this report govern its own conclusions.

The amendment of 21 September 2026 incorporates the twelve main coverage
findings and six additional quantitative refinements identified in the
companion coverage audit. Their result-level locations are listed in
Section~\ref{br:coverage-map}. New summaries retain the recovered-argument
status and their explicit conditional or finite-volume scope; restoration
does not certify every historical proof afresh. The pre-amendment report
and PDF are preserved in \path{history/report_2026-09-15/}.

In particular, the two arguments shown in the user's screenshots were
not lost. The small-\(\beta\) note contains both the original
\(O(\beta^{6/5})\) remainder and its later optimized remainder.
The stretching-limit proof survives a fresh check and is restored here
to the verified results. The previous selective write-up's omission of
that limit was an omission, not a mathematical retraction.
\clearpage

% ===== END front =====

% ===== BEGIN core =====
\hypertarget{sharp-homogeneous-poisson-upper-domination}{%
\section{Sharp homogeneous Poisson upper domination}\label{sharp-homogeneous-poisson-upper-domination}}

For \(0<\beta<1\), let \[
H_\beta(z,w)=\frac1\beta
\exp\!\left(\frac{\pi z\bar w}{\beta}
 -\frac{\pi(|z|^2+|w|^2)}{2\beta}\right),\qquad
G_\beta=\operatorname{DPP}(\beta H_\beta).
\] Here \(H_\beta\) is an orthogonal projection of diagonal \(1/\beta\), and \(G_\beta\) has intensity one. Set \[
L_\beta=-\frac{\log(1-\beta)}{\beta}.
\]

\textbf{Theorem 1.} \[
G_\beta\preceq\Pi_{L_\beta}.
\tag{1.1}
\] Moreover, if \(G_\beta\preceq\Pi_\lambda\) for a homogeneous Poisson process, then \(\lambda\ge L_\beta\). Consequently, \[
\boxed{\quad
r_c(G_\beta)\ge\frac{r_c(\Pi_1)}{\sqrt{L_\beta}}.
\quad}
\tag{1.2}
\] The stochastic order in (1.1) is inclusion order on locally finite configurations. The critical-radius bound is a lower bound.

\hypertarget{proof-of-domination}{%
\subsection{Proof of domination}\label{proof-of-domination}}

Write \(\mu_s=\mathcal L(\operatorname{DPP}(sH_\beta))\), for \(0<s<1\), and \(\mu_s^{!q}\) for its reduced Palm law.

Goldman's one-point Ginibre--Palm coupling, after spatial scaling, couples the full projection process with its reduced Palm process by adding exactly one point. Independently retain every point with probability \(s\). The added point is not retained with probability \(1-s\), independently of the thinned Palm remainder. Therefore, as measures, \[
\mu_s\ge(1-s)\mu_s^{!q}.
\tag{1.3}
\] This use of Goldman's theorem does not assert independence of the added point's location and the Palm remainder.

Let \(F\) be a bounded increasing function determined by a bounded set \(W\). For an independent Poisson process \(\Pi_\nu\), put \[
C(s,\nu)=\mathbb EF(\operatorname{DPP}(sH_\beta)\cup\Pi_\nu),
\qquad D_qF(\eta)=F(\eta\cup\{q\})-F(\eta).
\] Independent thinning, followed by the Campbell formula of the full projection process, gives \[
\partial_s C(s,\nu)
 =\frac1\beta\int_W
      \mathbb E_{\mu_s^{!q},\Pi_\nu}D_qF\,dA(q).
\tag{1.4}
\] Poisson insertion gives \[
\partial_\nu C(s,\nu)
 =\int_W\mathbb E_{\mu_s,\Pi_\nu}D_qF\,dA(q).
\tag{1.5}
\] These differentiations are justified directly: conditional on the full projection process in \(W\), thinning involves finitely many Bernoulli variables. The derivative is bounded by \(\operatorname{osc}(F)N(W)\), whose expectation is finite. The same couplings give continuity of both derivatives, with one-sided derivatives at the endpoints.

Since \(D_qF\ge0\), (1.3)--(1.5) imply \[
\partial_s C(s,\nu)
\le\frac1{\beta(1-s)}\partial_\nu C(s,\nu).
\tag{1.6}
\] Integrate this inequality along \[
\nu(s)=\frac1\beta\log\frac{1-s}{1-\beta},
\qquad 0\le s\le\beta.
\] The chain rule shows that \(C(s,\nu(s))\) is nonincreasing. Its initial law is \(\Pi_{L_\beta}\), and its final law is \(G_\beta\). Thus (1.1) holds for every bounded increasing local function.

For completeness, local stochastic order extends to whole-plane inclusion order. On each bounded open disk, the order on locally finite counting measures is closed, being determined by integration against nonnegative continuous compactly supported functions. Strassen's closed-order theorem gives a coupling of the restrictions. Extend each to a coupling of the full laws using regular conditional exterior laws. These couplings form a tight family because their marginals are fixed. Any subsequential weak limit preserves all compactly supported order inequalities, and hence is a whole-plane inclusion coupling.

The external input here is \href{https://arxiv.org/pdf/math/0610243}{Goldman's Theorem 1}, whose printed normalization has intensity \(1/\pi\). Spatial scaling gives the projection normalization used above.

\hypertarget{proof-that-l_beta-is-optimal}{%
\subsection{\texorpdfstring{Proof that \(L_\beta\) is optimal}{Proof that L\_\textbackslash beta is optimal}}\label{proof-that-l_beta-is-optimal}}

Let \(D_R=B(0,R)\) and \(Q_R=1_{D_R}H_\beta1_{D_R}\). The Fredholm hole formula gives \[
-\log\mathbb P(G_\beta(D_R)=0)
 =\sum_{j\ge1}\frac{\beta^j}{j}\operatorname{Tr}(Q_R^j).
\tag{1.7}
\] Also \[
\operatorname{Tr}Q_R=\frac{|D_R|}{\beta},\qquad
\operatorname{Tr}(Q_R-Q_R^2)
 =\int_{D_R}\int_{D_R^c}|H_\beta(x,y)|^2\,dA(y)dA(x)
 =O_\beta(R).
\tag{1.8}
\] To verify the last estimate, use \(|H_\beta(x,y)|^2=\beta^{-2}e^{-\pi|x-y|^2/\beta}\). For displacement \(h\), the area of \(\{x\in D_R:x+h\notin D_R\}\) is at most \(C(R|h|+|h|^2)\). Integrating this bound against the Gaussian proves (1.8).

The eigenvalues of \(Q_R\) lie in \([0,1]\). Hence, for fixed \(j\), \[
0\le\operatorname{Tr}Q_R-\operatorname{Tr}(Q_R^j)
\le(j-1)\operatorname{Tr}(Q_R-Q_R^2).
\] Therefore \(\operatorname{Tr}(Q_R^j)/|D_R|\to1/\beta\). The bound \(\operatorname{Tr}(Q_R^j)\le|D_R|/\beta\), together with \(\sum_{j\ge1}\beta^j/j<\infty\), justifies dominated convergence in (1.7): \[
\lim_{R\to\infty}
\frac{-\log\mathbb P(G_\beta(D_R)=0)}{|D_R|}
=L_\beta.
\tag{1.9}
\] If \(G_\beta\preceq\Pi_\lambda\), its hole probability is at least \(e^{-\lambda|D_R|}\). Equation (1.9) forces \(L_\beta\le\lambda\).

\hypertarget{infinite-volume-implication}{%
\subsection{Infinite-volume implication}\label{infinite-volume-implication}}

In an inclusion coupling, every occupied component of \(G_\beta\) is contained in an occupied component of \(\Pi_{L_\beta}\), using the same disk radius. Poisson scaling gives \(r_c(\Pi_\lambda)=r_c(\Pi_1)/\sqrt\lambda\), proving (1.2). Alternatively, apply the local comparison to connection from a fixed bounded seed to distance \(L\), then let \(L\uparrow\infty\). Only finitely many grains touch the seed, so the decreasing intersection is exactly connection of that seed to an unbounded component. No claim about percolation at the Poisson critical radius is needed.

\hypertarget{exact-signed-interpolation-formula}{%
\section{Exact signed interpolation formula}\label{exact-signed-interpolation-formula}}

Let \(H=H_1\), \(G_s=\operatorname{DPP}(sH)\), and \[
C_F(s,\nu)=\mathbb P(F(G_s\cup\Pi_\nu)=1)
\] for a local increasing Boolean event \(F\). For two absent points \(q,z\), define \[
\begin{aligned}
\mathcal R_{q,z}(\eta)
 &=\mathbf1\{F(\eta)=0,\ F(\eta+q)=F(\eta+z)=1\},\\
\mathcal C_{q,z}(\eta)
 &=\mathbf1\{F(\eta)=F(\eta+q)=F(\eta+z)=0,\
                         F(\eta+q+z)=1\}.
\end{aligned}
\] Thus \(\mathcal R\) records two individually successful alternatives, whereas \(\mathcal C\) records success requiring both insertions.

Translate one fixed Goldman coupling at the origin to obtain a jointly measurable family of couplings at \(q\): \[
X^q\sim G^{!q},\qquad X^q+\{Z_q\}\sim G.
\] Its added point has density \(e^{-\pi|z-q|^2}\), but is generally dependent on \(X^q\). Let \(X_s^q\) be an independent \(s\)-thinning of \(X^q\), and set \(Y_s^q=X_s^q\cup\Pi_\nu\).

\textbf{Theorem 2.} \[
\boxed{\quad
(\partial_s-\partial_\nu)C_F(s,\nu)
=s\int
\mathbb E\!\left[
 \mathcal R_{q,Z_q}(Y_s^q)-\mathcal C_{q,Z_q}(Y_s^q)
\right]\,dA(q).
\quad}
\tag{2.1}
\] The integrals are over a bounded active region of \(F\). Along the equal-intensity path \(\nu=1-s\), (2.1) is its ordinary derivative.

\hypertarget{proof}{%
\subsection{Proof}\label{proof}}

For every increasing Boolean \(F\), checking its four possible input configurations gives \[
D_qF(\eta)-D_qF(\eta+z)
=\mathcal R_{q,z}(\eta)-\mathcal C_{q,z}(\eta).
\tag{2.2}
\] Thin the full coupled process with common independent marks. The extra point survives with an independent Bernoulli variable \(B_s\) of parameter \(s\). The ordinary background has law \(Y_s^q+B_s\{Z_q\}\), while \(Y_s^q\) is its DPP reduced-Palm background with the same independent Poisson bath. Consequently, writing \[
J_q=\mathbb ED_qF(Y_s^q),\qquad
I_q=\mathbb ED_qF(G_s\cup\Pi_\nu),
\] equation (2.2) gives \[
J_q-I_q
=s\,\mathbb E[
\mathcal R_{q,Z_q}(Y_s^q)-\mathcal C_{q,Z_q}(Y_s^q)].
\] The derivative formulas are \(\partial_s C_F=\int J_q\,dA(q)\) and \(\partial_\nu C_F=\int I_q\,dA(q)\): the full parent has intensity one. This proves (2.1).

For each fixed bounded active window, the common thinning marks eventually retain all its points as \(s\uparrow1\). Bounded convergence therefore justifies this endpoint, also after integration over active roots. This does not give convergence uniform in growing windows.

There is also an exact formulation under the ordinary current law. Disintegrate the coupling given \(G_s\), and let \(w_s(q,G_s,z)\) be the conditional probability that \(z\) is its surviving deleted point. For \(0<s<1\), its no-deletion probability is \[
(1-s)\frac{d\mu_s^{!q}}{d\mu_s}(G_s)=1-T_q^s(G_s),
\qquad
\sum_{z\in G_s}w_s(q,G_s,z)=T_q^s(G_s).
\] Then the right side of (2.1) is equivalently \[
\int\mathbb E\sum_{z\in G_s}w_s(q,G_s,z)
 [\mathcal R_{q,z}-\mathcal C_{q,z}]
 ((G_s\setminus\{z\})\cup\Pi_\nu)\,dA(q).
\tag{2.3}
\] There is no additional \(s\) in (2.3); its weights already include the survival probability. Neither formula supplies the sign.

\hypertarget{exact-fock-variational-formula-and-a-negative-local-example}{%
\section{Exact Fock variational formula, and a negative local example}\label{exact-fock-variational-formula-and-a-negative-local-example}}

Let \[
\mathcal F=\left\{f\text{ entire}:
\|f\|_{\mathcal F}^2=\int_{\mathbb C}|f(z)|^2e^{-\pi|z|^2}dA(z)
<\infty\right\}.
\] Fix a finite distinct set \(A\), let \(Q=H-P_A\) be the projection onto the corresponding Gaussian-gauged functions vanishing on \(A\), and fix \(q\notin A\). For bounded Borel \(B\), let \(Q_B=1_BQ1_B\). For \(0\le s\le1\), define \[
M_{A,B,s}(q)=Q(q,q)+s\left\langle
 1_BQ(\cdot,q),(I-sQ_B)^{-1}1_BQ(\cdot,q)
\right\rangle.
\]

\textbf{Theorem 3.} \[
\boxed{\quad
\frac1{M_{A,B,s}(q)}
=\min_{\substack{f\in\mathcal F,\ f|_A=0\\
                   f(q)=e^{\pi|q|^2/2}}}
\int |f(z)|^2e^{-\pi|z|^2}
                 (1-s\mathbf1_B(z))\,dA(z).
\quad}
\tag{3.1}
\] In particular \(s=1\) gives the norm outside \(B\).

\hypertarget{proof-1}{%
\subsection{Proof}\label{proof-1}}

In the physical \(L^2\) gauge, set \(T=I-sQ1_BQ\) on \(\operatorname{ran}Q\). The operator \(Q1_BQ\) is compact: its trace is \(\int_B Q(z,z)dA(z)\le|B|\). Its norm is strictly less than one. Otherwise compactness gives a nonzero eigenvector supported in bounded \(B\), whereas its analytic representative vanishes on the open complement of a sufficiently large disk. Analyticity makes that impossible.

Thus \(T\) is positive and invertible, including at \(s=1\). Evaluation at \(q\) in the \(T\)-inner product has representer \(T^{-1}Q(\cdot,q)\). Its squared norm is \(\langle Q(\cdot,q),T^{-1}Q(\cdot,q)\rangle=M_{A,B,s}(q)\); the last equality is the resolvent identity. Cauchy--Schwarz gives (3.1), with equality at the normalized representer.

For every \(q\notin A\), whether or not \(q\in B\), the determinant lemma identifies \(M_{A,B,s}(q)\) as \[
Q(q,q)\,
\frac{\mathbb P(\operatorname{DPP}(sQ^{!q})(B)=0)}
     {\mathbb P(\operatorname{DPP}(sQ)(B)=0)}.
\tag{3.2}
\] For the process \(\operatorname{DPP}(sQ)\), the actual root intensity times this void ratio is \(sM_{A,B,s}(q)\).

\hypertarget{a-local-multiplier-can-be-below-one}{%
\subsection{A local multiplier can be below one}\label{a-local-multiplier-can-be-below-one}}

Fix a connection range \(R>0\). Let \(A=\{-a,a\}\) on the real axis, where \(R/2<a<R\), and let \[
B_a=B(-a,R)\cap B(a,R).
\] The parents are disconnected, and inserting the origin joins them. Direct inversion of their \(2\times2\) coherent Gram matrix gives \[
P_A(0,0)=\frac{2e^{-\pi a^2}}{1+e^{-2\pi a^2}}
=\operatorname{sech}(\pi a^2).
\] Also \(|B_a|\to0\) as \(a\uparrow R\). Since \(Q(z,z)\le1\), \[
M_{A,B_a,s}(0)
\le1-\operatorname{sech}(\pi a^2)
       +\frac{|B_a|}{1-|B_a|},\qquad |B_a|<1.
\] For \(a\) sufficiently close to \(R\), this is at most \[
1-\tfrac12\operatorname{sech}(\pi R^2)<1
\] uniformly for \(0\le s\le1\). Thus an empty connecting lens does not give a multiplier at least one for every two-parent geometry. This counterexample concerns a proposed pointwise estimate, not the critical radius.

\hypertarget{buffered-conditional-domination-with-prescribed-zeros}{%
\section{Buffered conditional domination with prescribed zeros}\label{buffered-conditional-domination-with-prescribed-zeros}}

Let \(D=B(p,R)\), let \(P\subset D\) be a fixed finite distinct set, and let \(C\) have compact closure in \(D\). Write \(Q_P=H-P_P\). Let \(B_D^P\) be the physical reproducing kernel of the Gaussian weighted analytic Bergman space on \(D\) whose functions vanish on \(P\).

\textbf{Theorem 4.} For any DPP with kernel \(0\le K\le Q_P\), conditional on its entire configuration outside \(D\), its restriction to \(C\) is stochastically dominated by \(\operatorname{DPP}(B_D^P|_C)\). In particular, almost surely in the observed exterior, \[
\mathbb P(N(C)>0\mid X_{D^c})
\le\int_C B_D^P(z,z)dA(z).
\tag{4.1}
\] The buffer \(D\setminus C\) is unobserved. This is a statement about the reduced remainder process. To infer a hole in \(C\) after restoring the prescribed Palm marks, those marks must also lie outside \(C\).

\hypertarget{proof-2}{%
\subsection{Proof}\label{proof-2}}

Use the usual unitary Ginibre translation to center the Gaussian weight at \(p\). The accompanying conjugate phase factors preserve operator order and determinantal laws; all kernels below use this common gauge.

First observe the exact configuration in the finite annulus \(A_L=B(p,L)\setminus D\), \(L>R\). On the bounded disk \(W=B(p,L)\), all compressed kernels have norm strictly below one. Their Janossy \(L\)-operators satisfy \[
L_K=K_W(I-K_W)^{-1}
\le Q_{P,W}(I-Q_{P,W})^{-1}=L_Q.
\] After conditioning on the finite annular configuration, the conditional interior \(L\)-operator is a Schur complement, so is no larger than the \(D\)-block of \(L_Q\). Applying \(L\mapsto L(I+L)^{-1}\) shows that the conditional correlation kernel is no larger than the kernel of \(Q_P\) conditioned to have no points in \(A_L\). Singular tuples of zero Janossy density may be omitted.

The latter kernel is the reproducing kernel of Fock functions vanishing on \(P\), with norm \[
\int_D |f|^2e^{-\pi|z-p|^2}dA
+\int_{B(p,L)^c}|f|^2e^{-\pi|z-p|^2}dA.
\] Restriction to \(D\) is contractive into the weighted Bergman space with the same zeros. Taking dual norms of finite linear combinations of evaluations therefore bounds this reproducing kernel by \(B_D^P\).

Its restriction to \(C\) is a positive trace-class contraction. Indeed choose \(M<R\) with \(C\subset B(p,M)\). The unconstrained weighted disk Bergman compression has eigenvalues \(p_n(M)/p_n(R)\), where \[
p_n(t)=\frac1{n!}\int_0^{\pi t^2}e^{-u}u^n\,du.
\] These ratios decrease in \(n\), since increasing \(n\) tilts the radial density toward larger radii. Its norm is \((1-e^{-\pi M^2})/(1-e^{-\pi R^2})<1\). The constrained kernel is smaller. Kernel-order stochastic domination follows from \href{https://arxiv.org/html/1406.2707v1}{Lyons, Theorem 3.8}: the kernels are locally trace-class positive contractions on the same area-measure space. Extend them by zero to the plane if the Borel set \(C\) is not itself locally compact.

For each bounded increasing local test, let \(L\uparrow\infty\). The conditional expectations converge by the bounded martingale theorem. A countable determining class gives the asserted conditional-law domination. Truncating the count and then using monotone convergence proves (4.1). No inverse on the full exterior is taken.

\hypertarget{explicit-clustered-zero-bound}{%
\subsection{Explicit clustered-zero bound}\label{explicit-clustered-zero-bound}}

For a \(k\)-fold zero at \(p\), interpreted as deletion of the first \(k\) derivative-evaluation modes, let \(C\subset B(p,M)\) and \(\rho=M/R<1\). Orthogonality of the centered monomials gives \[
\int_D|z-p|^{2n}e^{-\pi|z-p|^2}dA
\ge e^{-\pi R^2}\frac{\pi R^{2n+2}}{n+1}.
\] Summing the remaining modes \(n\ge k\) yields \[
\boxed{\quad
\int_C B_D^{(k,p)}(z,z)dA
\le
\frac{|C|e^{\pi R^2}}{\pi R^2}
\frac{(k+1)\rho^{2k}}{(1-\rho^2)^2}.
\quad}
\tag{4.2}
\] For fixed \(k\), distinct zeros tending to \(p\) converge to this derivative constraint: use divided differences of the evaluation vectors and convergence of their nonsingular Gram matrices. Thus sufficiently close distinct marks give (4.2) with any prescribed positive additive tolerance. No coincident actual points are asserted. These marks must still be supplied, with their probabilities, in a percolation argument.

\hypertarget{unbounded-exterior-palm-likelihood-ratios}{%
\section{Unbounded exterior Palm likelihood ratios}\label{unbounded-exterior-palm-likelihood-ratios}}

\textbf{Theorem 5.} Let \(D=B(0,T)\), let \(A\subset D\) be a fixed finite distinct set, and take distinct \(q,z\in D\setminus A\) with \(q^2\ne z^2\). Let \(\nu_q,\nu_z\) be the restrictions to \(D^c\) of the canonical reduced Ginibre Palm laws at \(A\cup\{q\}\) and \(A\cup\{z\}\). Then \[
\nu_z\sim\nu_q,\qquad
\operatorname*{ess\,sup}_{\nu_q}\frac{d\nu_z}{d\nu_q}=\infty,
\qquad
\operatorname*{ess\,inf}_{\nu_q}\frac{d\nu_z}{d\nu_q}=0.
\tag{5.1}
\] Hence no finite constant bounds either exterior law by the other. The restriction \(q^2\ne z^2\) suffices for this theorem; no assertion about the omitted symmetric case is needed here.

The published inputs are the radial conditional-density and same-order Palm-equivalence results of \href{https://arxiv.org/html/1608.03736v2}{Bufetov--Qiu, Proposition 1.1 and Corollary 1.6}, and Ginibre number rigidity and fixed-count conditional tolerance from \href{https://arxiv.org/abs/1211.2381}{Ghosh--Peres}. The proof below explains the support construction and fixed-tuple passage.

\hypertarget{control-of-the-full-distant-product}{%
\subsection{5.1 Control of the full distant product}\label{control-of-the-full-distant-product}}

For \(L>2T\), take centered radial cutoffs in \[
g_{>L}(v)=\prod_{|x|>L}(1-v/x),\qquad |v|\le T.
\] The analytic logarithm of each factor is unambiguous, and \[
\log g_{>L}(v)=-vS_1(L)-\tfrac12v^2S_2(L)
-\sum_{|x|>L}\sum_{j\ge3}\frac{v^j}{jx^j}.
\] Here \(S_1(L)=\lim_{U\to\infty}\sum_{L<|x|\le U}x^{-1}\) requires care. It is not justified by an exterior \(L^2(dA)\) bound on \(1/x\), which diverges.

For a smooth compactly supported complex function \(f\), the projection variance identity and the fundamental theorem of calculus give \[
\operatorname{Var}\sum_{x\in G}f(x)
=\frac12\iint|f(x)-f(y)|^2e^{-\pi|x-y|^2}dA(x)dA(y)
\le\frac1{2\pi}\int|\nabla f|^2dA.
\tag{5.2}
\] The displayed constant is a sufficient bound; no optimal constant is needed. Approximate the radial cutoff of \(1/x\) by cutoffs changing over strips of width one at radii \(L,U\). Their Dirichlet integrals are \(O(L^{-1}+U^{-1})\). The difference from the sharp cutoff has \(L^2(dA)\) norm squared of the same order; the elementary DPP variance bound applies to that difference. All means vanish by angular integration. Thus, for an absolute \(C\), \[
\mathbb E|S_1(L,U)|^2\le C/L,\qquad U>L\ge2.
\] Applying this to differences of outer cutoffs proves the radial \(L^2\) limit. The elementary variance estimate also gives \[
\mathbb E|S_2(L)|^2\le\pi/L^2,\qquad
\mathbb E\sum_{|x|>L}|x|^{-3}=2\pi/L.
\] As \(\sum_{j\ge3}|v/x|^j/j\le(2/3)|v/x|^3\), it follows that \[
\mathbb E\sup_{|v|\le T}|\log g_{>L}(v)|
\le \frac{C^{1/2}T}{\sqrt L}
 +\frac{\sqrt\pi T^2}{2L}+\frac{4\pi T^3}{3L}
\longrightarrow0.
\tag{5.3}
\] The normalized radial product in the BQ conditional-density formula has the same limit, by uniqueness of convergence in probability. There is no free unaccounted-for linear analytic factor.

\hypertarget{realizing-a-large-quadratic-exterior-field}{%
\subsection{5.2 Realizing a large quadratic exterior field}\label{realizing-a-large-quadratic-exterior-field}}

Fix an inside count \(n\), an error \(\varepsilon>0\), and \(M\ge16\pi\). Choose \(L\) large. With positive probability both the product in (5.3) is within \(\varepsilon/3\) of zero in logarithm and \[
\pi L^2/2\le N(B_L)\le2\pi L^2.
\tag{5.4}
\] The product condition tends to probability one. The count condition does too, by its mean \(\pi L^2\) and variance at most \(\pi L^2\). Both are measurable from the exterior of \(B_L\); this uses number rigidity for the count.

Condition on that exterior first. Its inside count is a fixed \(N\). Set \(J=\lfloor(N-n)/2\rfloor\) and \(d_J=\sqrt{J/M}\). Place \(n\) points in any positive-volume open set inside \(D\), \(J\) points in a small patch near \(+id_J\), and \(J\) in a patch near \(-id_J\). An odd remaining point goes near \(3iL/4\). For large \(L\), these patches lie outside \(D\) and inside \(B_L\). Their template product is \[
(1+Mv^2/J)^J\longrightarrow e^{Mv^2}
\] uniformly on \(\overline D\), because the logarithmic error is at most \(M^2T^4/J\). Taking patches small gives the same approximation for distinct actual points. The possible odd factor tends uniformly to one.

Fixed-count conditional tolerance assigns positive probability to each such open configuration set. There are only finitely many counts in (5.4), so the templates may depend measurably on \(N\). Integrating proves the existence of a positive-probability set \(E_M\) of actual \(D\)-exteriors with rigid count \(n\) and \[
g_O(v)=\exp(M\omega v^2+\eta_O(v)),\qquad
\sup_{\overline D}|\eta_O|<\varepsilon,
\tag{5.5}
\] where any unit \(\omega\) is obtained by rotating the patches. This ordering of conditioning is essential: the tail estimate has not been conditioned on a previously selected rare patch event.

\hypertarget{fixed-palm-tuples-not-just-almost-every-tuple}{%
\subsection{5.3 Fixed Palm tuples, not just almost every tuple}\label{fixed-palm-tuples-not-just-almost-every-tuple}}

For ordinary exterior \(O\), BQ identifies the conditional interior process as a polynomial ensemble with kernel \(K_O\), rank \(N(O)\), and weight \(e^{-\pi|v|^2}|g_O(v)|^2\). Its analytic range is contained in the weighted disk Bergman space, so \(K_O\le B_D\). Campbell disintegration initially gives, for almost every distinct \(m\)-tuple \(A\), \[
\frac{d\mu^A_{\rm out}}{d\mu_{\rm out}}(O)
=\frac{\det K_O(A,A)}{\det H(A,A)}.
\tag{5.6}
\] This identity holds for every fixed distinct tuple. Indeed, its right side is continuous in the tuple and dominated locally by products of disk Bergman diagonals. The canonical reduced Palm kernels depend locally uniformly on their distinct roots. Their correlation densities are bounded by one, so the absolutely convergent local Janossy series gives total-variation continuity on every bounded observation window. Equality for almost every tuple therefore extends to every fixed tuple on a countable determining class of exterior cylinder tests, and hence for the full exterior laws.

The same argument with a joint inside/exterior test identifies the conditional reduced interior law. The determinant-weighted finite polynomial-ensemble Palm integral is continuous, including the case of an empty remainder, and has the same deterministic domination. This supplies the needed disintegration without an uncountable intersection of unrelated null sets.

Apply the support construction with \(n=m+1\). Polynomial interpolation gives \(\det K_O(A,A)>0\) on \(E_M\) for every fixed distinct \(A\). Thus \(\mu^A_{\rm out}(E_M)>0\). Under that law, the residual conditional interior process has one point with density \[
t_O(v)=\frac1{Z_O}
 |p_A(v)|^2e^{-\pi|v|^2}|g_O(v)|^2,\qquad
p_A(v)=\prod_{a\in A}(v-a).
\tag{5.7}
\]

\hypertarget{the-likelihood-ratio}{%
\subsection{5.4 The likelihood ratio}\label{the-likelihood-ratio}}

Choose \(\omega\) so \(\Delta=\Re[\omega(z^2-q^2)]>0\). On \(E_M\), (5.5)--(5.7) imply \[
\frac{t_O(z)}{t_O(q)}
\ge
\left|\frac{p_A(z)}{p_A(q)}\right|^2
 e^{-\pi(|z|^2-|q|^2)}e^{2M\Delta-4\varepsilon}.
\] Divide the \((m+1)\)-root version of (5.6) by the \(m\)-root version. The Schur determinant formula gives \[
\frac{d\nu_z}{d\nu_q}(O)
=\frac{Q_A(q,q)}{Q_A(z,z)}\frac{t_O(z)}{t_O(q)}
\quad\hbox{on }E_M.
\] All fixed factors are positive. The event \(E_M\) has positive probability under both further Palm laws, since both conditional intensities are positive there. Increasing \(M\) proves unbounded essential supremum. Reversing \(\omega\) proves essential infimum zero. Same-order Palm equivalence supplies the asserted mutual absolute continuity.

This theorem rules out a uniform bound over arbitrary exteriors. It provides no lower bound on the probability of the adverse exterior sets and no conclusion about their weight in a pivotal event.

\hypertarget{exact-finite-ensemble-deletion-calculation-and-counterexample}{%
\section{Exact finite-ensemble deletion calculation and counterexample}\label{exact-finite-ensemble-deletion-calculation-and-counterexample}}

Use the rank-\(n\) Ginibre projection generated by the first \(n\) Gaussian-weighted monomials, and the reduced Palm root \(q=0\). If \(X\) has \(n-1\) points, a deletion rule \(w(G,z)\) reconstructs this Palm source from the rank-\(n\) ensemble precisely when its pointwise weights sum to one and, almost everywhere in \(X\), \[
\int e^{-\pi|z|^2}
\prod_{x\in X}\frac{|z-x|^2}{|x|^2}\,
w(X+z,z)\,dA(z)=1.
\tag{6.1}
\] This is the Campbell formula with the Vandermonde density ratio. The root kernel diagonal is one.

\textbf{Proposition 6.} In rank two the rule \[
w(\{x,z\},z)=\frac{|x|^2}{|x|^2+|z|^2}
\] is valid. In rank three, the normalized Lagrange-square rule is not valid.

For rank two, its conditional insertion density is \[
k(x,z)=e^{-\pi|z|^2}
\frac{|z-x|^2}{|x|^2+|z|^2}.
\] Angular averaging replaces \(|z-x|^2\) by \(|z|^2+|x|^2\), so its integral is one. The two deletion weights also sum to one. Thus both marginals, not just an averaged normalization, are correct.

For the failed rule, put \[
L_i(0)=\prod_{j\ne i}\frac{-z_j}{z_i-z_j},
\qquad S(G)=\sum_i|L_i(0)|^2,\qquad
w_L(G,z_i)=\frac{|L_i(0)|^2}{S(G)}.
\] The Vandermonde terms in (6.1) cancel, leaving the necessary condition \[
\int e^{-\pi|z|^2}\frac1{S(X+z)}\,dA(z)=1.
\tag{6.2}
\] In rank three, take \(X=\{a,-a\}\), \(a>0\). Direct interpolation gives \[
S(\{a,-a,z\})
=\frac{a^4+\frac{|z|^2}{2}(a^2+|z|^2)}
 {|a^2-z^2|^2}.
\] The angular mean of its reciprocal is \[
\frac{a^4+t^4}{a^4+\frac12a^2t^2+\frac12t^4},
\qquad t=|z|.
\] It tends to two for every \(t>0\) as \(a\downarrow0\). The interpolation identity \(\sum_iL_i(0)=1\) gives \(1/S\le3\). Dominated convergence therefore makes the left side of (6.2) tend to two. Continuity in the distinct remainder points, with the same Gaussian majorant, makes the failure persist on an open set of positive Palm probability. This is not merely a coincidence example.

% ===== END core =====

\setcounter{theorem}{6}

% ===== BEGIN ellipse =====
\section{The ellipse route: large grains, the Poisson limit, and its scope}
\label{ell:section}

This section records the results obtained from the ellipse route, with
their proofs.  In particular, the critical-radius convergence theorem
below is an established conclusion of the argument.  It was omitted
from one of the shorter summaries of this investigation; that omission
was not a mathematical refutation.  The theorem does not, however,
order the undeformed Ginibre critical radius and the limiting Poisson
critical radius.

Throughout this section, $G$ has intensity one with respect to planar
Lebesgue measure $dA$.  Its kernel is
\begin{equation}\label{ell:ginibre-kernel}
 H(z,w)=\exp\left(\pi z\overline w
       -\frac\pi2|z|^2-\frac\pi2|w|^2\right).
\end{equation}
For a locally finite configuration $\xi$, write
\[
 \mathcal O_r(\xi)=\bigcup_{x\in\xi}\overline B(x,r),
 \qquad
 \theta_\xi(r)=\mathbb P\bigl(0\text{ belongs to an unbounded
 component of }\mathcal O_r(\xi)\bigr)
\]
when $\xi$ denotes a random configuration, and set
\[
 r_c(\xi)=\inf\{r>0:\theta_\xi(r)>0\}.
\]
Thus $r$ is the grain radius and two grains intersect when their
centers are at distance at most $2r$.  All probabilities in this
section are ordinary probabilities at a fixed spatial point, not Palm
probabilities.  Open disks will be used for hole events and closed
disks for grains.  The intensity measures of the processes to which
we apply the results are absolutely continuous, so a fixed circle
contains no point almost surely.

Define the area-preserving deformation and its transformed process by
\begin{equation}\label{ell:deformation}
 S_t(x,y)=(e^t x,e^{-t}y),\qquad \Phi_t=S_t^{-1}G.
\end{equation}
For area $A>0$, put $r=\sqrt{A/\pi}$ and
\begin{equation}\label{ell:ellipse}
 E_{A,t}=S_t\overline B(0,r)
 =\{(x,y):e^{-2t}x^2+e^{2t}y^2\le A/\pi\}.
\end{equation}
If $\theta_G(E)$ is the probability that the origin belongs to an
unbounded component of $\bigcup_{x\in G}(x+E)$, then the deterministic
homeomorphism $S_t^{-1}$ gives
\begin{equation}\label{ell:affine-equivalence}
 \theta_G(E_{A,t})=\theta_{\Phi_t}(\sqrt{A/\pi}).
\end{equation}
It preserves the origin and boundedness of components.  Since
$\det S_t=1$, every $\Phi_t$ has intensity one.

\subsection{A contour bound for bounded occupied components}

The contour argument needs only a particular consequence of negative
association.  Say that a stationary point process $\Phi$ has the
disjoint-hole inequality if, for pairwise disjoint bounded Borel sets
$D_1,\ldots,D_k$,
\begin{equation}\label{ell:hole-na}
 \mathbb P\left(\Phi\left(\bigcup_{j=1}^kD_j\right)=0\right)
 \le\prod_{j=1}^k\mathbb P(\Phi(D_j)=0).
\end{equation}
For such a process, define
\[
 q_\Phi(r)=\mathbb P(\Phi(B(0,r))=0),\qquad
 f_\Phi(r)=\mathbb P\bigl(0\in\mathcal O_r(\Phi)
       \text{ and its component is bounded}\bigr).
\]

\begin{lemma}[Contour bound]\label{ell:contour}
For every $\varepsilon\in(0,1)$ there are constants
$C_\varepsilon<\infty$ and $\eta_\varepsilon>0$, independent of
$\Phi$ and $r>0$, such that a stationary locally finite process
satisfying \eqref{ell:hole-na} obeys
\begin{equation}\label{ell:contour-bound}
 q_\Phi((1-\varepsilon)r)\le\eta_\varepsilon
 \quad\Longrightarrow\quad
 f_\Phi(r)\le C_\varepsilon
                  q_\Phi((1-\varepsilon)r)^2.
\end{equation}
\end{lemma}

\begin{proof}
Fix $r>0$ and use a translated square grid of mesh
$h=\varepsilon r/\sqrt2$, with the origin at a square center.
Suppose the origin is covered and its occupied component $C$ is
bounded.  Every center of a grain belonging to $C$ lies in $C$;
local finiteness therefore makes $C$ a finite union of closed disks.
In particular, $C$ is compact.

Let $Q$ be the union of all closed grid squares meeting $C$.
There are finitely many.  At every point of $C$, all incident grid
squares are included, so $C\subset\operatorname{int}Q$.
The squares of $Q$ are connected through shared edges.  One way to
see this is to follow a continuous path in the finite disk union
$C$: transitions through square edges preserve edge connectivity,
and transitions through a grid vertex include all four incident
squares.  The external boundary of $Q$ can consequently be traced
by a closed grid-edge walk surrounding the origin, with distinct
edges, although vertices can be revisited.  Let $n$ be its length.

For every $p\in\partial C$,
\begin{equation}\label{ell:boundary-clearance}
 \operatorname{dist}(p,\Phi)\ge r.
\end{equation}
Indeed, a grain whose interior contains $p$ meets $C$ and belongs
to the same component; it would make $p$ an interior point of $C$.
Every vertex $v$ of the external grid boundary is outside $C$ and
belongs to a square meeting $C$.  Hence it is within distance
$\sqrt2h$ of $\partial C$.  The distance function to $\Phi$ is
$1$-Lipschitz, and \eqref{ell:boundary-clearance} implies
\begin{equation}\label{ell:vertex-void}
 \Phi(B(v,(1-\varepsilon)r))=0.
\end{equation}

The component contains a complete grain, so its enclosing contour
has horizontal span at least $2r$.  Choose the leftmost and rightmost
contour vertices first.  Their distance exceeds
$2(1-\varepsilon)r$.  Extend this pair to a maximal collection of
contour vertices separated by at least $2(1-\varepsilon)r$.
For a fixed contour this selection is made deterministically.
There is an integer $K_\varepsilon$ bounding the number of grid
vertices in any disk of radius $2(1-\varepsilon)r$, because the
ratio of this radius to $h$ depends only on $\varepsilon$.
The contour has at least $n/2$ distinct vertices: its $n$ distinct
edges contribute $2n$ incidences and each grid vertex has degree
at most four.  Maximality of the selection thus gives at least
\[
 \max\{2,n/M_\varepsilon\},\qquad M_\varepsilon=2K_\varepsilon,
\]
selected vertices.  Their open disks in \eqref{ell:vertex-void}
are pairwise disjoint.  If $u=q_\Phi((1-\varepsilon)r)$,
stationarity and \eqref{ell:hole-na} bound the probability of all
the empty disks required by this fixed contour by
$u^{\max\{2,n/M_\varepsilon\}}$.

There are at most $C(n+1)^2 4^n$ length-$n$ grid walks enclosing
the origin, for an absolute constant $C$.  In fact, their convex
hull contains the origin, and their diameter is at most $nh$.
Every possible starting vertex is therefore within distance $nh$
of the origin; there are $O((n+1)^2)$ such vertices and at most
$4^n$ step sequences from each.  A union bound yields
\begin{equation}\label{ell:contour-series}
 f_\Phi(r)\le C\sum_{n\ge4}(n+1)^2 4^n
                  u^{\max\{2,n/M_\varepsilon\}}.
\end{equation}
For $n\le2M_\varepsilon$ this is bounded by a constant depending
only on $\varepsilon$ times $u^2$.  Writing
$n=2M_\varepsilon+k$ in the remaining terms and factoring out
$u^2$ leaves a constant times
\[
 \sum_{k\ge1}(2M_\varepsilon+k+1)^2
                   (4u^{1/M_\varepsilon})^k.
\]
This is uniformly bounded when
$4u^{1/M_\varepsilon}\le1/2$.
The choice $\eta_\varepsilon=8^{-M_\varepsilon}$ is sufficient;
the case $u=0$ follows directly from the contour union bound.
\end{proof}

\begin{corollary}[Failure is asymptotically a hole]\label{ell:hole-failure}
Assume the hypotheses of Lemma~\ref{ell:contour} and
\begin{equation}\label{ell:quartic-hole}
 \log q_\Phi(r)=-c r^4+o(r^4),\qquad c>0.
\end{equation}
Then, with $\log0=-\infty$,
\[
 \limsup_{r\to\infty}\frac{\log f_\Phi(r)}{r^4}\le-2c,
 \qquad f_\Phi(r)=o(q_\Phi(r)).
\]
If $\Phi$ has no point on a fixed circle almost surely, then
\begin{equation}\label{ell:failure-equivalence}
 1-\theta_\Phi(r)\sim q_\Phi(r).
\end{equation}
\end{corollary}

\begin{proof}
For each fixed $\varepsilon\in(0,1)$,
Lemma~\ref{ell:contour} and \eqref{ell:quartic-hole} give
\[
 \limsup_{r\to\infty}\frac{\log f_\Phi(r)}{r^4}
 \le-2c(1-\varepsilon)^4.
\]
Taking the infimum of these valid upper bounds as
$\varepsilon\downarrow0$ proves the first assertion.  More directly,
fix any $\varepsilon$ with $2(1-\varepsilon)^4>1$.  Then
\[
 \log\frac{f_\Phi(r)}{q_\Phi(r)}
 \le -c\bigl(2(1-\varepsilon)^4-1\bigr)r^4+o(r^4)
 \longrightarrow-\infty.
\]
No estimate uniform in $\varepsilon$ is needed.  The failure event
is the disjoint union of the uncovered event and the event defining
$f_\Phi(r)$.  The circle-null condition makes the first probability
equal to $q_\Phi(r)$ with the stated boundary conventions.
\end{proof}

\subsection{The fixed-eccentricity large-area comparison}

The external determinantal input used here and below is
\href{https://arxiv.org/html/1406.2707v1}{Lyons, Theorem 3.7}:
a locally trace-class positive contraction on a Radon measure space
defines a negatively associated point process.  The theorem applies
to arbitrary bounded increasing functions supported on disjoint
Borel sets.  It also applies to decreasing functions by replacing
both functions by their negatives.  Iteration gives
\eqref{ell:hole-na}.  The Ginibre projection and its images under
invertible area-preserving linear transformations satisfy the
operator hypotheses.  Such transformations also preserve
disjointness of the supporting sets.

The hole asymptotic is a published input, rather than a new result
of this investigation.  In the normalization of intensity $1/\pi$,
\href{https://arxiv.org/html/1604.08363v3}{Adhikari--Reddy,
Theorem 6 and Example 17} give, for fixed semiaxes $a,b>0$,
\begin{equation}\label{ell:ar-hole}
 \lim_{R\to\infty}\frac1{R^4}
 \log\mathbb P\left(G_{1/\pi}\left(
 \left\{(x,y):\frac{x^2}{R^2a^2}+\frac{y^2}{R^2b^2}<1\right\}
 \right)=0\right)
 =-\frac{(ab)^3}{2(a^2+b^2)}.
\end{equation}
Their domain theorem assumes an open domain in the unit disk and,
in the case used here, an exterior-ball condition.  An ellipse is
convex and satisfies that condition.  To remove the unit-disk
restriction, choose $d>0$ with $d\max(a,b)<1$, use semiaxes
$da,db$ and scale $R/d$.  The coefficient is homogeneous of
degree four, so the factors $d^4$ cancel.  The ellipse coefficient
in \eqref{ell:ar-hole} has been checked against their displayed
Example 17; the normalization is not inferred from a disk formula
alone.

\begin{theorem}[Equal-area disks versus fixed noncircular ellipses]
\label{ell:large-area}
For every fixed $t\in\mathbb R$, at unit intensity,
\begin{align}
 1-\theta_G(E_{A,t})
  &\sim \mathbb P(G(E_{A,t})=0),\label{ell:large-equivalence}\\
 \lim_{A\to\infty}\frac{-\log(1-\theta_G(E_{A,t}))}{A^2}
  &=\frac1{4\cosh(2t)}.\label{ell:large-rate}
\end{align}
For every fixed $t\ne0$, there is $A_0(t)<\infty$ such that
\begin{equation}\label{ell:large-strict}
 A>A_0(t)\quad\Longrightarrow\quad
 \theta_G(E_{A,0})>\theta_G(E_{A,t}).
\end{equation}
More precisely,
\begin{equation}\label{ell:failure-ratio}
 \lim_{A\to\infty}\frac1{A^2}
  \log\frac{1-\theta_G(E_{A,t})}{1-\theta_G(E_{A,0})}
 =\frac14\left(1-\frac1{\cosh(2t)}\right).
\end{equation}
\end{theorem}

\begin{proof}
The dilation $x\mapsto\sqrt\pi x$ maps unit-intensity Ginibre to
the intensity-$1/\pi$ normalization.  Apply
\eqref{ell:ar-hole} with $a=e^t$, $b=e^{-t}$, and
$R=\sqrt\pi r$.  Since $ab=1$ and
$a^2+b^2=2\cosh(2t)$, this gives
\begin{equation}\label{ell:unit-hole}
 \log q_{\Phi_t}(r)
 =-\frac{\pi^2 r^4}{4\cosh(2t)}+o_t(r^4).
\end{equation}
The process $\Phi_t$ is stationary and satisfies the disjoint-hole
inequality, and its intensity measure charges no fixed circle.
Corollary~\ref{ell:hole-failure} applies with
$c=\pi^2/(4\cosh(2t))$.  The affine equivalence
\eqref{ell:affine-equivalence}, and $A=\pi r^2$, prove
\eqref{ell:large-equivalence} and \eqref{ell:large-rate}.
Subtracting the two logarithmic asymptotics gives
\eqref{ell:failure-ratio}.  For fixed $t\ne0$, its right side
is strictly positive, so the failure-probability ratio eventually
exceeds one.  This proves \eqref{ell:large-strict} and identifies
exactly where strictness enters.
\end{proof}

\begin{remark}[Intensity scaling and quantifiers]\label{ell:large-scope}
For Ginibre intensity $\lambda>0$, replace $A^2$ in the denominators
of \eqref{ell:large-rate} and \eqref{ell:failure-ratio} by
$\lambda^2A^2$: dilation by $\sqrt{\pi\lambda}$ reduces the
claim to \eqref{ell:ar-hole}.  The eccentricity is fixed before
$A\to\infty$.  Neither the error term nor $A_0(t)$ has been made
uniform as $t\to0$ or $|t|\to\infty$.  A comparison at sufficiently
large area does not compare the areas where percolation first begins.
\end{remark}

\subsection{Exact kernel and local convergence in total variation}

We now let $t\to+\infty$.  Write $z=x+iy$ and $w=x'+iy'$.
Because \eqref{ell:deformation} preserves area, the kernel of
$\Phi_t$ with respect to $dA$ is $K_t(z,w)=H(S_tz,S_tw)$.
Expanding the exponent gives
\begin{equation}\label{ell:anisotropic-kernel}
 K_t(z,w)=\exp\left\{
 -\frac\pi2\left[e^{2t}(x-x')^2+e^{-2t}(y-y')^2\right]
 +i\pi(yx'-xy')\right\}.
\end{equation}
In particular, the phase is unchanged.  Let $W$ be a bounded Borel
window of area $m$, contained in a horizontal strip of height $H_W$.
On $L^2(W,dA)$ let $A_t$ be the integral operator with kernel
$|K_t(z,w)|$, and set
\begin{equation}\label{ell:at-st}
 a_t=\sqrt2 H_W e^{-t},\qquad
 s_t=\int_{W^2}|K_t(z,w)|^2\,dA(z)dA(w).
\end{equation}
The absolute kernel is a positive definite Gaussian convolution
kernel on the plane, compressed to $W$.  Thus $A_t$ is positive
and trace class, with trace $m$.  Direct integration and Schur's
bound give
\begin{equation}\label{ell:operator-bounds}
 \|A_t\|\le a_t,\qquad
 \operatorname{Tr}(A_t^2)=s_t\le mH_W e^{-t}.
\end{equation}
For example, the first estimate follows by integrating the
horizontal Gaussian over the whole line and bounding the vertical
integral by $H_W$.  Positivity and the spectral theorem then yield
\begin{equation}\label{ell:cycle-trace}
 \operatorname{Tr}(A_t^k)\le a_t^{k-2}s_t\qquad(k\ge2).
\end{equation}

\begin{proposition}[Quantitative local Poisson approximation]
\label{ell:tv}
For $2a_t<1$,
\begin{equation}\label{ell:tv-bound}
 d_{\mathrm{TV}}\bigl(\mathcal L(\Phi_t|_W),
                    \mathcal L(\Pi_1|_W)\bigr)
 \le\frac{e^{2m}}2
  \left[\exp\left(\frac{2s_t}{1-2a_t}\right)-1\right].
\end{equation}
Consequently the restrictions converge in total variation for
every fixed bounded $W$.
\end{proposition}

\begin{proof}
Write
$\rho_n(z_1,\ldots,z_n)=\det[K_t(z_i,z_j)]_{i,j=1}^n$,
and put $\rho_0=1$.  The diagonal is one, so Hadamard's
inequality gives $0\le\rho_n\le1$.  In the determinant
expansion the identity permutation contributes one.  Bound all
other terms in absolute value, integrate over $W^n$, and group
permutations by cycle type.  A fixed point contributes $m$ and
a cycle of length $k\ge2$ contributes
$\operatorname{Tr}(A_t^k)$.  The exponential generating function
for permutations with these cycle weights gives, for $u a_t<1$,
\begin{align}
 &\sum_{n\ge0}\frac{u^n}{n!}
       \int_{W^n}|\rho_n-1|\,dA^n\notag\\
 &\hspace{8mm}\le
 e^{um}\left[
   \exp\left(\sum_{k\ge2}\frac{u^k}{k}
                           \operatorname{Tr}(A_t^k)\right)-1
             \right].\label{ell:cycle-majorant}
\end{align}
For clarity, the exponential arises because $j_k$ cycles of
length $k$ have combinatorial weight $1/(k^{j_k}j_k!)$.
Every term in this majorant is nonnegative, so all regroupings
are justified by Tonelli's theorem.  By \eqref{ell:cycle-trace},
\[
 \sum_{k\ge2}\frac{2^k}{k}\operatorname{Tr}(A_t^k)
 \le 2s_t\sum_{j\ge0}(2a_t)^j
 =\frac{2s_t}{1-2a_t}.
\]

The Janossy density for exactly $n$ points in $W$ is
\begin{equation}\label{ell:janossy}
 j_n(z_1,\ldots,z_n)
 =\sum_{k\ge0}\frac{(-1)^k}{k!}
       \int_{W^k}\rho_{n+k}(z_1,\ldots,z_n,w_1,\ldots,w_k)
                 \,dA^k(w).
\end{equation}
The bound $0\le\rho_j\le1$ makes this inclusion--exclusion
series absolutely convergent, also after integration over the
first $n$ coordinates.  The Poisson Janossy density is $e^{-m}$,
obtained by replacing every $\rho_j$ by one.  Hence
\begin{align*}
 2d_{\mathrm{TV}}(\Phi_t|_W,\Pi_1|_W)
 &=\sum_{n\ge0}\frac1{n!}\int_{W^n}|j_n-e^{-m}|\,dA^n\\
 &\le\sum_{n,k\ge0}\frac1{n!k!}
                         \int_{W^{n+k}}|\rho_{n+k}-1|\,dA^{n+k}\\
 &=\sum_{j\ge0}\frac{2^j}{j!}
                         \int_{W^j}|\rho_j-1|\,dA^j.
\end{align*}
Apply \eqref{ell:cycle-majorant} at $u=2$.  Finally
$a_t,s_t\to0$ for fixed $W$, proving convergence.
\end{proof}

\begin{remark}\label{ell:tv-scope}
This approximation holds for every measurable event in a fixed
window, including closed-grain crossing events.  No continuity-set
argument or tangency exception is needed for the transfer of such
event probabilities.  The factor $e^{2m}$ is kept explicitly:
the proposition is not a uniform approximation on windows growing
arbitrarily with $t$.
\end{remark}

\subsection{Two finite-size criteria using negative association}

The next arguments supply the passage from local convergence to
infinite-volume percolation.  Neither a bound on a mixing length
nor isotropy of the approximating processes is used.

Fix $r>0$ and $L>2r$, and tile the plane by closed squares
\[
 S_v=L\bigl(v+[-1/2,1/2]^2\bigr),\qquad v\in\mathbb Z^2.
\]
Overlapping square boundaries cause no ambiguity in the event
definitions.  Supports of events within each color class below
are separated by a positive distance.

\begin{lemma}[Upper finite-size criterion]\label{ell:upper-criterion}
There is an absolute $b_0>0$ with the following property.
Let $\Phi$ be stationary, locally finite and negatively associated.
For each nearest-neighbor edge $e=\{v,w\}$ of $\mathbb Z^2$,
let $R_e=S_v\cup S_w$ and declare $e$ good if the occupied set
has a long-direction crossing of $R_e$ and a transverse crossing
of each of $S_v,S_w$.  Crossings are paths restricted to the
indicated rectangles.  If both orientations of edges have bad
probability at most $b<b_0$, then
$\mathcal O_r(\Phi)$ has an unbounded component with positive
probability.
\end{lemma}

\begin{proof}
All three crossing events are increasing.  Since the grain radius
is the deterministic number $r$, the good-edge event depends only
on centers in $R_e+\overline B(0,r)$; centers outside this set
cannot contribute even to a path on the rectangle boundary.

First check deterministic gluing.  A long crossing of $R_e$
contains a crossing of each endpoint square in the same direction:
take the relevant subpath between the exterior side and the middle
side.  For two collinear good edges, the transverse crossing of
their common square intersects the square crossings supplied by
both long witnesses.  For two good edges making a turn, the long
witnesses themselves supply horizontal and vertical crossings of
the common square, and these intersect.  The fact that horizontal
and vertical paths in a closed rectangle intersect follows from
planar separation; it does not require independence or uniqueness
of crossings.  Thus witnesses for any connected family of good
edges belong to one occupied component.  If the edge family is
infinite, these witnesses are spatially unbounded.

Color the lattice edges by their orientation and the coordinates
modulo $5$ of their canonical initial vertex (the endpoint with
smaller coordinate in the edge direction).  There are $50$
colors.  Two different edges of a color are separated by at least
$5L$ in one coordinate before the sizes of their rectangles are
taken into account.  A rectangle has side lengths $2L,L$; after
expansion by $r$, both side lengths remain below $3L$, since
$2r<L$.  Thus expanded supports are disjoint within each color.

A prescribed set of $n$ bad edges contains at least $n/50$
edges of one color.  Their events are decreasing and supported
on disjoint sets.  Negative association, iterated over these
sets, gives
\begin{equation}\label{ell:bad-edge-product}
 \mathbb P(\text{all prescribed edges are bad})\le b^{n/50}.
\end{equation}
If the good-edge cluster of the lattice origin is finite, its
outer edge boundary contains a simple dual circuit of bad edges
surrounding the origin.  The number of length-$n$ possible circuits
is at most $C(n+1)^2 4^n$: choose a starting vertex within distance
$O(n)$ of the origin, then a sequence of nearest-neighbor steps.
Consequently
\begin{equation}\label{ell:peierls}
 \mathbb P(\text{the origin's good-edge cluster is finite})
 \le C\sum_{n\ge4}(n+1)^2\bigl(4b^{1/50}\bigr)^n.
\end{equation}
Choose an absolute $b_0>0$ for which the series is less than one.
The probability of an infinite good-edge cluster is then positive;
the gluing argument proves the stated occupied percolation.
\end{proof}

\begin{lemma}[Lower finite-size criterion]\label{ell:lower-criterion}
Under the same stationarity, local finiteness and negative
association assumptions, let $H_v$ be the event that
$\mathcal O_r(\Phi)$ has a path from $S_v$ to the boundary of
the concentric square
\[
 T_v=L\bigl(v+[-3/2,3/2]^2\bigr),
\]
with its path contained in $T_v$.  If
\begin{equation}\label{ell:small-annulus}
 p=\mathbb P(H_v)<8^{-25},
\end{equation}
then almost surely $\mathcal O_r(\Phi)$ has no unbounded component.
\end{lemma}

\begin{proof}
The event $H_v$ is increasing and depends only on centers in
$T_v+\overline B(0,r)$.  Color vertices $v$ by their two
coordinates modulo $5$.  These supports have side lengths
$3L+2r<4L$, and hence are disjoint within each of the $25$
colors.  For any prescribed $n$ distinct vertices,
negative association implies that the probability that all their
$H_v$ events hold is at most $p^{n/25}$.

If an unbounded occupied component exists, the tiles it meets
form an infinite connected set for star adjacency (squares may
meet at an edge or a vertex).  The finite-degree lattice graph
then contains an infinite self-avoiding star path in that set.
Every tile on that path has $H_v$: follow an occupied path from
the component's intersection with $S_v$ until its first meeting
with $\partial T_v$.  Such finite occupied paths exist because
a locally finite union of fixed-radius disks has path-connected
components.  There are at most $8^n$ candidate $n$-vertex star
paths from a specified tile.  Thus the probability of an infinite
such path from that tile is bounded, for every $n$, by
\[
 8^n p^{n/25}=(8p^{1/25})^n\longrightarrow0.
\]
Take a countable union over starting tiles.
\end{proof}

\begin{remark}[Existence and origin membership]\label{ell:existence}
For a stationary process with fixed radius $r>0$, positive
probability of some unbounded occupied component is equivalent
to positive probability of origin membership in such a component.
Indeed, an unbounded disk component contains the interior of a
disk and hence a rational point.  The existence event is contained
in the countable union of the events that a specified rational
point belongs to an unbounded component.  If the union has
positive probability, one event does; stationarity makes its
probability equal to the origin probability.  The converse is
immediate.  No ergodicity assertion is needed for this equivalence.
\end{remark}

\subsection{Convergence of the actual critical radii}

The Poisson input is
\href{https://arxiv.org/html/1605.05926v1}{Ahlberg--Tassion--Teixeira,
Theorem 1.1(i),(iii)}.  For a radius distribution with finite
second moment, fixed-aspect-ratio occupied rectangle crossings
tend to one above the Poisson critical intensity and to zero
below it.  The deterministic unit-radius distribution meets that
hypothesis.  Dilation $x\mapsto x/r$ changes intensity-one
radius-$r$ disks into intensity-$r^2$ unit disks.  Thus their
theorem gives precisely the off-critical crossing limits needed
at a fixed radius $r$ on either side of $r_c(\Pi_1)$.

\begin{theorem}[A local-convergence criterion for critical radii]
\label{ell:general-limit}
Let $(\Xi_n)$ be stationary, locally finite, negatively associated
planar point processes.  Suppose that for every bounded Borel set
$W$,
\[
 d_{\mathrm{TV}}\bigl(\mathcal L(\Xi_n|_W),
                     \mathcal L(\Pi_1|_W)\bigr)\longrightarrow0.
\]
For deterministic disk grains,
\begin{equation}\label{ell:general-rc-limit}
 r_c(\Xi_n)\longrightarrow r_c(\Pi_1).
\end{equation}
\end{theorem}

\begin{proof}
Write $r_P=r_c(\Pi_1)$.  First fix $r>r_P$.
The cited Poisson crossing limit, and a union bound for the three
crossings defining a good edge, allow one to choose a finite
$L>2r$ for which the Poisson bad-edge probability is below
$b_0/2$.  Rotation invariance gives this for both orientations.
Now keep $r$ and $L$ fixed.  Each good-edge event depends on one
fixed bounded expanded rectangle.  Local total-variation
convergence makes both bad-edge probabilities below $b_0$
for all sufficiently large $n$.  Lemma~\ref{ell:upper-criterion}
and Remark~\ref{ell:existence} give percolation at radius $r$.
Thus $r_c(\Xi_n)\le r$ eventually.  Since this holds for every
$r>r_P$,
\[
 \limsup_{n\to\infty}r_c(\Xi_n)\le r_P.
\]

Next fix $0<r<r_P$.  An $H_v$ path implies a short-direction
crossing of one of four strips of dimensions $3L$ by $L$.
For example, translate $v$ to zero and suppose the path first
reaches the top side of $T_0$ at height $3L/2$.  Its subpath
after its last visit to height $L/2$ stays in
$[-3L/2,3L/2]\times[L/2,3L/2]$ and crosses that strip vertically.
The other first-exit sides give the other three strips.
Consequently the subcritical Poisson crossing limit and a union
bound give $\mathbb P_{\Pi_1}(H_0)\to0$ as $L\to\infty$.
Choose one finite $L>2r$ with
\[
 \mathbb P_{\Pi_1}(H_0)<\tfrac12 8^{-25}.
\]
For this fixed choice, local total-variation convergence gives
$\mathbb P_{\Xi_n}(H_0)<8^{-25}$ for sufficiently large $n$.
Lemma~\ref{ell:lower-criterion} proves nonpercolation at $r$.
Monotonicity in the grain radius then gives
$r_c(\Xi_n)\ge r$.  Letting $r\uparrow r_P$ yields
\[
 \liminf_{n\to\infty}r_c(\Xi_n)\ge r_P.
\]
This proves \eqref{ell:general-rc-limit}.
\end{proof}

\begin{corollary}[The anisotropic Ginibre threshold limit]
\label{ell:ginibre-limit}
For \eqref{ell:deformation},
\begin{equation}\label{ell:main-rc-limit}
 \boxed{\displaystyle\lim_{t\to+\infty}r_c(S_t^{-1}G)
                         =r_c(\Pi_1).}
\end{equation}
Equivalently, the critical areas of equally oriented Ginibre
ellipse grains of shape $t$ converge to the Poisson disk critical
area $\pi r_c(\Pi_1)^2$.
\end{corollary}

\begin{proof}
Each $\Phi_t$ is stationary, locally finite and negatively
associated by the determinantal input stated above.
Proposition~\ref{ell:tv} supplies local total-variation convergence.
Apply Theorem~\ref{ell:general-limit} along every sequence
$t_n\to+\infty$; the sequential characterization of a real limit
proves \eqref{ell:main-rc-limit}.  Finally use
\eqref{ell:affine-equivalence} and $A=\pi r^2$.
\end{proof}

\begin{remark}[The order of choices]\label{ell:limit-scope}
The proof fixes a radius strictly away from the Poisson threshold,
then one finite block size, and only then sends the deformation
parameter to infinity.  Infinite-volume percolation is established
by the finite-size criteria, not by exchanging an infinite-volume
limit with local convergence. This qualitative proof alone supplies neither a rate
for \eqref{ell:main-rc-limit} nor a sign for
$r_c(\Phi_t)-r_c(\Pi_1)$. The separate recovered rate in
Section~\ref{ell:recovered-rate} uses an additional near-critical input.
\end{remark}

\subsection{A joint large-area and large-deformation regime}

The fixed-$t$ quartic hole exponent is not uniform in the
deformation parameter.  The following direct Fredholm estimate
describes a different joint regime and makes that distinction
quantitative.

\begin{proposition}[Thin-ellipse hole regime]\label{ell:joint-hole}
Let $q_t(r)=\mathbb P(\Phi_t(B(0,r))=0)$.  If
$2\sqrt2 r e^{-t}<1$, then
\begin{equation}\label{ell:joint-hole-bound}
 0\le-\log q_t(r)-\pi r^2
 \le\frac{\pi r^3e^{-t}}{1-2\sqrt2 r e^{-t}}.
\end{equation}
Along every sequence with $r\to\infty$ and $r e^{-t}\to0$,
\begin{equation}\label{ell:joint-hole-limit}
 \frac{-\log q_t(r)}{\pi r^2}\longrightarrow1,
 \qquad
 1-\theta_{\Phi_t}(r)\sim q_t(r).
\end{equation}
\end{proposition}

\begin{proof}
Let $T$ be the compression of $K_t$ to $B(0,r)$.
It is a positive contraction with trace $\pi r^2$.
Use \eqref{ell:operator-bounds} with $m=\pi r^2$ and
$H_W=2r$.  Taking the absolute kernel in the Schur estimate
gives
\[
 \|T\|\le2\sqrt2 r e^{-t},\qquad
 \operatorname{Tr}(T^2)\le2\pi r^3e^{-t}.
\]
The Fredholm determinant hole formula and the spectral theorem
give, under the displayed norm restriction,
\begin{align*}
 -\log q_t(r)-\pi r^2
 &=\sum_{k\ge2}\frac{\operatorname{Tr}(T^k)}k\\
 &\le\frac{\operatorname{Tr}(T^2)}{2(1-\|T\|)}.
\end{align*}
The series is absolutely convergent; all traces are nonnegative.
This proves \eqref{ell:joint-hole-bound} and the first limit.

Fix $\varepsilon\in(0,1)$ with $2(1-\varepsilon)^2>1$.
The same estimate applies at $(1-\varepsilon)r$ along the given
sequence.  The constants in Lemma~\ref{ell:contour} are uniform
over the processes $\Phi_t$, so for all sufficiently large
indices it gives
\[
 \log\frac{f_{\Phi_t}(r)}{q_t(r)}
 \le\log C_\varepsilon
    -\pi\bigl(2(1-\varepsilon)^2-1\bigr)r^2+o(r^2)
 \longrightarrow-\infty.
\]
The hole and finite-component decomposition proves the second
assertion in \eqref{ell:joint-hole-limit}.
\end{proof}

\begin{remark}\label{ell:regime-scope}
At fixed $t$, \eqref{ell:unit-hole} has an exponent proportional
to $r^4$.  In Proposition~\ref{ell:joint-hole}, the exponent is
asymptotic to $\pi r^2$.  The latter regime makes the minor
semiaxis $r e^{-t}$ tend to zero while the area tends to infinity.
Both assertions concern large grains; neither by itself is a
near-critical comparison.
\end{remark}

\subsection{The first finite-window perturbation at the Poisson end}

This calculation retains the geometry of the leading interaction.
Unlike the isotropically concentrated interaction in a thinned
Ginibre limit, the anisotropic interaction concentrates only the
horizontal separation of two insertion locations.

For a bounded measurable functional $F$ determined by the points
in a fixed bounded window $W$, define
\[
 D_zF(\xi)=F(\xi\cup\{z\})-F(\xi),\qquad
 D_{z_1,\ldots,z_n}F=D_{z_1}\cdots D_{z_n}F.
\]
Insertion differences are zero outside $W$.  Coincident insertion
locations form a null set in the integrals below and cause no
ambiguity.  Keep $m,A_t,a_t,s_t$ from
\eqref{ell:at-st}--\eqref{ell:operator-bounds}, and put
\begin{equation}\label{ell:Bt}
 B_t=\sum_{k\ge2}\frac{2^k}{k}\operatorname{Tr}(A_t^k).
\end{equation}
For fixed $W$, $B_t=2s_t+O_W(e^{-2t})$ and
$B_t\le2s_t/(1-2a_t)$ when $2a_t<1$.

\begin{proposition}[Exact leading local correction]
\label{ell:local-perturbation}
For $2a_t<1$,
\begin{align}
 \mathbb EF(\Phi_t)
 &=\mathbb EF(\Pi_1)
 -\frac12\int_{W^2}|K_t(z,w)|^2
             \mathbb E[D_zD_wF(\Pi_1)]\,dA(z)dA(w)
 +R_t,\label{ell:perturbation}\\
 |R_t|&\le \|F\|_\infty e^{2m}
                  (e^{B_t}-1-2s_t)
       =O_W(\|F\|_\infty e^{-2t}).\label{ell:remainder}
\end{align}
This is a fixed-window statement, with no asserted uniformity
over growing percolation boxes.
\end{proposition}

\begin{proof}
For every finite configuration $\eta\subset W$, the elementary
subset expansion is
\[
 F(\eta)=\sum_{X\subseteq\eta}D_XF(\varnothing).
\]
Since $|D_XF|\le2^{|X|}\|F\|_\infty$ and the factorial
densities are bounded by one, Campbell's formula and absolute
summability give the Newton expansion
\begin{equation}\label{ell:newton}
 \mathbb EF(\Phi_t)=\sum_{n\ge0}\frac1{n!}
      \int_{W^n}\rho_n(Z)D_ZF(\varnothing)\,dA^n(Z).
\end{equation}
The sum of the absolute integrals is at most
$\|F\|_\infty e^{2m}$.  The same formula holds for Poisson
with $\rho_n=1$.

Expand each determinant in \eqref{ell:newton}.  The identity
permutation gives $\mathbb EF(\Pi_1)$.  A permutation with
exactly one transposition $(i,j)$ and all remaining coordinates
fixed contributes $-|K_t(z_i,z_j)|^2$.  By symmetry and
$\binom n2/n!=1/(2(n-2)!)$, summing these terms gives
\begin{align*}
 -\frac12\int_{W^2}|K_t(z,w)|^2
  \sum_{k\ge0}\frac1{k!}\int_{W^k}
      D_{z,w,u_1,\ldots,u_k}F(\varnothing)\,dA^k(u)
                  \,dA(z)dA(w).
\end{align*}
The inner sum is the Poisson Newton expansion of
$\mathbb E D_zD_wF(\Pi_1)$.

All remaining permutations contain a cycle of length at least
three or at least two transpositions.  The same absolute cycle
majorant as in \eqref{ell:cycle-majorant}, now multiplied by
$\|F\|_\infty$, bounds their total contribution by
$\|F\|_\infty e^{2m}(e^{B_t}-1-2s_t)$.
The subtraction $2s_t$ is exactly the majorant weight of a single
transposition at $u=2$.  Finally,
\[
 0\le B_t-2s_t
 \le 2s_t\frac{2a_t}{1-2a_t}=O_W(e^{-2t}),
 \qquad s_t=O_W(e^{-t}),
\]
so $e^{B_t}-1-2s_t=O_W(e^{-2t})$.  Every interchange above
is covered by these absolute majorants.
\end{proof}

\begin{proposition}[The vertical-pair coefficient for a rectangle crossing]
\label{ell:vertical-coefficient}
Fix $r>0$ and a nondegenerate closed rectangle
$R=I_R\times J_R$.  Let $F$ indicate a left-to-right crossing
of $R$ by radius-$r$ occupied disks, the crossing path being
restricted to $R$.  Choose a closed rectangle
$W=I_W\times J_W$ containing
$R+\overline B(0,r)$.  Then
\begin{align}
 &\lim_{t\to\infty}e^t
       \bigl(\mathbb EF(\Phi_t)-\mathbb EF(\Pi_1)\bigr)
 \notag\\
 &\qquad=-\frac12\int_{I_W}\int_{J_W}\int_{J_W}
       \mathbb E\left[D_{(a,b)}D_{(a,c)}F(\Pi_1)\right]
                  \,dc\,db\,da.\label{ell:vertical-limit}
\end{align}
The coefficient is a finite signed integral.  No assertion about
its sign is made.
\end{proposition}

\begin{proof}
Write
$h(z,w)=\mathbb E[D_zD_wF(\Pi_1)]$, taking $h=0$ if an
insertion is outside the determining window.  The only point
requiring care beyond Proposition~\ref{ell:local-perturbation}
is the existence of the horizontal trace of $h$ at equal
horizontal coordinates.  We prove that for Lebesgue-almost every
$(a,b,c)$,
\begin{equation}\label{ell:trace}
 h((a,b),(a+\delta,c))\longrightarrow
 h((a,b),(a,c))\quad(\delta\to0).
\end{equation}

Conditional on the finite Poisson configuration in $W$, represent
the crossing by the finite intersection graph of clipped convex
grains $C_u=\overline B(u,r)\cap R$, with additional marks
indicating intersection with the left and right sides of $R$.
The event holds exactly when this graph connects the two marks.
Its incidences are locally stable in the insertion coordinates
away from the following contact equalities.

For one fixed center $u$ with $C_u\ne\varnothing$, contact
between its clipped grain and a varying disk centered at $v$
is governed by membership in
$C_u+\overline B(0,r)$.  Its boundary is the boundary of a
convex body with nonempty interior and has planar Lebesgue
measure zero.  The same observation applies to meeting a marked
side segment, and to appearance or disappearance of a clipped
grain, using respectively the segment or $R$ in place of $C_u$.
These facts, with Fubini's theorem, exclude all relevant contacts
between inserted and Poisson grains for almost every
$(a,b,c)$ and almost every Poisson configuration.  Existing
grains do not move; any contact among them is in any case a
null event under their joint densities.

It remains to check contact between the two inserted clipped
grains when their horizontal coordinates coincide.  Put
$d=|b-c|/2$ and $h_0=(b+c)/2$.  Their intersection is nonempty
exactly when
\begin{equation}\label{ell:covertical-contact}
 \operatorname{dist}(a,I_R)^2+
       \bigl(d+\operatorname{dist}(h_0,J_R)\bigr)^2\le r^2.
\end{equation}
Indeed, for $(x,y)\in R$, the maximum of its squared distances
to $(a,b)$ and $(a,c)$ is
$(x-a)^2+(d+|y-h_0|)^2$, whose minimum over the rectangle
is the left side of \eqref{ell:covertical-contact}.
On each of the finitely many regions specifying the nearest
points in the two intervals and the sign of $b-c$, equality
in \eqref{ell:covertical-contact} is a nontrivial polynomial
equation.  Its solution set has three-dimensional Lebesgue
measure zero.  Coincident centers $b=c$ also form a null set.

Outside these exceptional sets every relevant graph incidence
is stable under a sufficiently small horizontal displacement
of the second inserted center.  Therefore all four Boolean
terms defining $D_zD_wF$ stabilize for almost every fixed
background and almost every triple.  The difference is bounded,
so conditional Poisson integration and dominated convergence
prove \eqref{ell:trace}.

Now use \eqref{ell:anisotropic-kernel} in the integral of
\eqref{ell:perturbation} and substitute $x=a$,
$x'=a+e^{-t}u$, $y=b$, $y'=c$.  After multiplication by
$e^t$, its integrand becomes
\[
 e^{-\pi u^2}e^{-\pi e^{-2t}(b-c)^2}
       h((a,b),(a+e^{-t}u,c)),
\]
with the window indicator also imposed.  For almost every
$(a,b,c,u)$, \eqref{ell:trace} supplies the limit.  A constant
multiple of $e^{-\pi u^2}$, with $(a,b,c)$ restricted to a
fixed bounded box, is an integrable dominating function.
The window boundary contributes a null set.  Since
$\int_{\mathbb R}e^{-\pi u^2}\,du=1$, dominated convergence
gives the right side of \eqref{ell:vertical-limit}.  The
remainder vanishes after multiplication by $e^t$ by
\eqref{ell:remainder}.
\end{proof}

\begin{remark}[The remaining sign is a geometric question]
\label{ell:vertical-scope}
For an increasing Boolean event, $D_zD_wF$ equals $-1$ when
the background fails but either insertion separately succeeds;
it equals $+1$ when the background and both single insertions
fail but the pair succeeds; in all other cases it is zero.
Thus the minus sign in \eqref{ell:vertical-limit} gives a
positive contribution from redundant individually sufficient
insertions and a negative contribution from cooperative pairs.
The two vertical coordinates $b,c$ need not coincide.
Consequently the repeated-insertion identity $D_z^2F=-D_zF$
does not determine this coefficient.  A sign estimate uniform
over the boxes needed near criticality has not been proved.
\end{remark}

\subsection{A recovered quantitative rate for the stretching limit}
\label{ell:recovered-rate}

The source \path{ellipse_critical_variation.md}, with
\path{ellipse_critical_variation_audit.md}, contains a stronger conclusion
than the qualitative limit above. We restore it as a recovered argument,
with its additional published input explicit. Write
$c(t)=r_c(S_t^{-1}G)$. The recorded estimate is
\begin{equation}\label{ell:polynomial-rate}
 |c(t)-r_P|\le C t^{-\eta/2}\qquad(t\ge t_0)
\end{equation}
for some $C,t_0<\infty$ and $\eta>0$.

For unit-radius Poisson disks,
\href{https://arxiv.org/html/2211.11661v1}{Manolescu--Santoro,
Theorem 2.4 and equation (2.16)} give a transition window of order
$[n^2\pi_4(n)]^{-1}$ and a polynomial lower bound
$\pi_4(n)\ge c n^{-2+\eta}$ on the critical four-arm probability.
The critical specialization of the latter follows by fixing $n$,
approaching critical intensity, and using finite-window continuity and
critical RSW. The authors give a proof sketch of this near-critical
input; this report uses their stated result rather than supplying a
replacement proof.

For windows of diameter $O(n)$, \eqref{ell:tv-bound} is at most
$Cn^3\exp(C_0n^2-t)$ once its norm condition holds. Choose
$n=\lfloor a_0\sqrt t\rfloor$, with $C_0a_0^2<1/2$.
The local error then tends to zero exponentially. Choose the Poisson
crossing tolerance below the two fixed tolerances in the upper and
lower negative-association criteria. Intensity deviations of order
$n^{-\eta}$ give the required Poisson block tests; the scaling
$\lambda=r^2$ converts them to radius deviations of the same order
near $r_P>0$. Transfer these tests to $\Phi_t$ in total variation and
apply the criteria. This gives \eqref{ell:polynomial-rate}.
The scale is selected before the local comparison; no unrestricted
exchange of a growing-window and deformation limit is used.
The rate has no asserted sign.

\subsection{Recovered orientation and critical-radius variation identities}
\label{ell:orientation}

Use the center-terminal convention of \path{ellipse_critical_variation.md}:
in the fixed original configuration $G$, connect a center in $B_m$ to
a center outside the open $B_L$ using edges
$|S_{-t}(p-q)|\le2r$. The terminal labels do not move with $t$ or $r$.
For $|t|\le T$ and $r\le R$, the event is local in $B_{L+2Re^T}$,
by stopping a witness at its first center outside $B_L$. Let
$U_{m,L}(t)$ be its minimax threshold capped at $R$, assigning $R$
when no path exists below that cap.

For an edge of weight $w(t)=|S_{-t}v|/2$ and transformed angle
$\theta_t$,
\[
 w'=-w\cos(2\theta_t),\qquad
 w''=w[2-\cos^2(2\theta_t)].
\]
At each fixed $t$, distinct pair weights tie only on Lebesgue-null
hypersurfaces. Thus the active edge below the cap is almost surely
unique, with angle $\Theta_{m,L}(t)$, and
\begin{equation}\label{ell:threshold-orientation}
 U'=-U\cos(2\Theta)\mathbf1_{\{U<R\}}
\end{equation}
almost everywhere. The capped threshold is $R$-Lipschitz.
Absolute continuity and Fubini justify integration of
\eqref{ell:threshold-orientation} before taking expectations.

If $\mu_t(dr,d\theta)$ is the joint law of the uncapped threshold
and active angle, the finite-window crossing probability satisfies
the distributional contact identity
\[
 \partial_t C(t,r)=rM(t,r)\partial_r C(t,r),\qquad
 M(t,r)=\mathbb E_{\mu_t}[\cos(2\Theta)\mid U=r],\quad |M|\le1,
\]
where the conditional mean is understood relative to the radius
contact measure. This identity supplies no favorable sign.

Deterministic grain inclusion gives
$e^{-|t-s|}c(s)\le c(t)\le e^{|t-s|}c(s)$.
Choose $R>\sup_{|t|\le T}c(t)$ and a slow diagonal of seeds and
outer radii valid for rational $t$ and noncritical rational radii.
The seed exhaustion gives convergence of capped thresholds in
probability and mean at rational $t$. The common Lipschitz bound
upgrades it to
$\sup_{|t|\le T}\mathbb E|U_n(t)-c(t)|\to0$.
Passing only the endpoint expectations through the integrated identity
therefore gives the recovered exact representation
\begin{equation}\label{ell:critical-orientation}
 c(T)-c(0)=-\lim_n\int_0^T
 \mathbb E[U_n(t)\mathbf1_{\{U_n(t)<R\}}
                         \cos(2\Theta_n(t))]\,dt.
\end{equation}
For each fixed $\delta>0$, contributions with
$|U_n(t)-c(t)|>\delta$ vanish in absolute value after integration.
No pointwise limit of angular contact measures is assumed.
Evenness of $c$ does not exclude a cusp at zero.

Combining \eqref{ell:critical-orientation} with the recovered rate,
an orientation integral with limit at most
$-CT^{-\eta/2}-\delta$ would give $c(0)\le r_P-\delta$.
This states a sufficient quantitative sign estimate; that estimate
remains unproved.

\subsection{Switching destroys the proposed convexity argument}
\label{ell:switching}

The finite calculation in \path{ellipse_minmax_switching.md} concerns
$T(t)=\min_\gamma\max_{e\in\gamma}|S_tv_e|/2$ on a finite Euclidean
graph with fixed terminal vertices. On intervals with one active
weight, $T=w_e$ and $w_e''>0$. At a switch, however,
\[
 T''=w_{\rm active}''(t)\,dt+
       \sum_j[T'(t_j+)-T'(t_j-)]\delta_{t_j}.
\]
Contract edges strictly below a tied threshold and discard those
strictly above it. Giving each remaining critical edge slope $v_e$,
the two directional derivatives are
\[
 T'(t_0+)=\min_\gamma\max_{e\in\gamma}v_e,\qquad
 T'(t_0-)=\max_\gamma\min_{e\in\gamma}v_e.
\]
These follow by expanding finitely many edge weights to first order.
Series alternatives give positive switching atoms; parallel
alternatives give negative ones.

Both occur with all Euclidean edges present. The points
$(0,0),(1,0),(1,1)$, with the first and last as terminals, give
$T(t)=e^{|t|}/2$ near zero and atom $+1$. For source $(0,0)$,
target $(1,0)$, and additional points $(0,1/10),(3/5,9/10)$,
the two short within-cluster edges leave two competing unit-length
bridges. Their slopes are $1/2$ and $-7/50$, giving atom $-16/25$.
After a common uniform rotation, the right derivatives of the mean
thresholds are respectively $1/\pi$ and $-4/(5\pi)$.

The negative example can be smoothed: replace $(1,0)$ by $(e^u,0)$,
sample $u$ from a smooth density $h$ supported in a sufficiently small
interval about zero, and then rotate uniformly. Splitting the
$u$-integral at the moving equality of the two bridge weights yields
\begin{equation}\label{ell:negative-curvature}
 M'(0)=0,\qquad
 M''(0)=\tfrac32M(0)-\tfrac{16}{25}h(0).
\end{equation}
Since $M(0)\le1/2$, taking $h(0)>75/64$ makes the curvature negative.
The moving-boundary term in the second derivative is essential;
ties at a fixed time have probability zero. The analogous smoothed
series example has $M''(0)=\tfrac32M(0)+h(0)$.
Thus neither isotropy nor positive edge curvature implies convexity
of an averaged connecting threshold. These finite-law examples do
not determine the Ginibre critical-radius curvature.

\subsection{Riesz coefficients and the two-point Palm angular response}
\label{ell:riesz-palm}

Further recovered calculations in \path{ellipse_critical_variation.md},
Sections 6--7, distinguish two finite approximations. For a fixed
rotation-invariant local event in the transformed process, isotropy
rewrites the covertical Poisson coefficient as
\begin{equation}\label{ell:riesz}
 \lim_{t\to\infty}e^t[\mathbb EF(\Phi_t)-\mathbb EF(\Pi_1)]
 =-\frac1{2\pi}\iint
       \frac{\mathbb E_\Pi D_xD_yF}{|x-y|}\,dx\,dy.
\end{equation}
The kernel is locally integrable and insertion support is bounded.
This is not repeated insertion at one point.

For two disjoint disks of radius $a$, separation $D>2a$, put
$u=\pi a^2$, $J_A=16\pi a^3/3$, and
$J_{AB}=\int_{A\times B}|x-y|^{-1}dx\,dy$.
For the event that both disks contain a point, averaged over common
rotations, its Riesz integral is
$2e^{-2u}[J_{AB}-(e^u-1)J_A]$.
It is positive if $u\le1$ and $a(D+2a)<3/32$.
The one-disk occupancy event instead has integral $-e^{-u}J_A$.
These examples give opposite leading signs without a claim about
the near-critical annulus.

Return now to fixed original centers and moving ellipse edges.
At a contact pair $x,y=x+2rS_t(\cos\theta,\sin\theta)$, let
$p_{t,r}(x,\theta)$ be the reduced two-point Palm probability that
the edge $xy$ is pivotal when toggled, with both centers restored.
Set
\[
 P_{t,r}(\theta)=\int p_{t,r}(x,\theta)\,dx,\qquad
 a_{t,r}(\theta)=1-e^{-4\pi r^2
                  (e^{2t}\cos^2\theta+e^{-2t}\sin^2\theta)}.
\]
Two-point Campbell and the contact Jacobian, divided by two for
unordered edges, give
\[
 \partial_r C=2r\int_0^{2\pi}a_{t,r}P_{t,r}\,d\theta,\qquad
 M=\frac{\int\cos(2\theta)a_{t,r}P_{t,r}\,d\theta}
          {\int a_{t,r}P_{t,r}\,d\theta}.
\]
At $t=0$, $P_{0,r}=P_0(r)$ is constant in angle. With $v=4\pi r^2$
and $D_t=\int a_{t,r}P_{t,r}\,d\theta$, continuity of the Palm
contact law yields, whenever $P_0(r)>0$,
\[
 M(t,r)=\frac{1-e^{-v}}{D_t}\int\cos(2\theta)
                   [P_{t,r}(\theta)-P_0(r)]\,d\theta
              +\frac{v}{e^v-1}t+o(t).
\]
The explicit pair-density term is positive; a desired negative
orientation response must overcome it through the Palm geometry.
That response cannot be omitted. Its kernel is exactly
$K^{!x,!y}=K-VG_{xy}^{-1}V^*$, with coherent columns
$V=[K(\cdot,x),K(\cdot,y)]$ and their Gram matrix $G_{xy}$.
Differentiation retains all three terms
$-V'G_{xy}^{-1}V^*-VG_{xy}^{-1}(V')^*
+VG_{xy}^{-1}G_{xy}'G_{xy}^{-1}V^*$.
Contact pairs stay separated on compact positive radius and
deformation intervals. Graph changes, including near-tie switches,
remain a further contribution. The source's independent audit
checks these formulas, not the missing sign.

\subsection{What these results do and do not establish}

Theorem~\ref{ell:large-area} supplies an actual strict comparison
of fixed-point percolation probabilities for equal-area disk
and fixed noncircular ellipse grains at sufficiently large
area.  Corollary~\ref{ell:ginibre-limit} supplies an actual
infinite-volume critical-threshold limit for the anisotropic
processes.  Proposition~\ref{ell:joint-hole} separates two
large-grain regimes, and Proposition~\ref{ell:vertical-coefficient}
identifies the finite-window first correction at the Poisson
end without assigning it an unproved sign. The recovered refinements
give a quantitative endpoint error, the exact integrated orientation
identity, and further signed Palm calculations. The switching examples
explain why positive edge curvature cannot settle their sign.

These are separate assertions.  They do not establish
$r_c(G)<r_c(\Pi_1)$.  In particular, neither monotonicity of
$t\mapsto r_c(\Phi_t)$ nor a one-sided comparison with its
limit follows from \eqref{ell:main-rc-limit}.  A logically
sufficient additional conclusion would be a proved radius
$r<r_c(\Pi_1)$ at which $\Phi_0=G$ percolates.  This section
does not supply that conclusion.  The large-area comparison
and the Poisson limit have been connected at the level of
well-defined models and valid finite-size criteria, but an
additional near-critical comparison is still absent.

% ===== END ellipse =====

% ===== BEGIN small_beta =====
\section{The recovered small-\texorpdfstring{$\beta$}{beta} argument}
\label{sb:section}

\subsection{Status and the result recorded in the archive}

The small-$\beta$ work has not been lost. The principal proof note is
\path{research/small_beta_critical_radius_asymptotic.md}; its recorded
audits are \path{small_beta_asymptotic_audit.md},
\path{pure_small_beta_theorem_audit.md}, and
\path{direct_pivotal_void_audit.md}. Those documents state that their
arguments passed the earlier checks. That historical statement is not
being substituted here for a complete new verification of every
geometric construction used by the argument.

The precise distinction is as follows. The finite-window differentiation
identities and the deterministic implication from a uniform pivotal
estimate to the asserted critical-radius asymptotic are written out
below. The necessary uniform geometric estimate is recorded with its
actual hypotheses and a substantial reconstruction of its proposed
proof. A bounded fresh review found no decisive counterexample, but did
not independently reconstruct every certificate selection, boundary
placement, and conditional normalization in that chain. Accordingly,
the geometric estimate, and hence the asymptotic which depends on it,
are designated here as \emph{archival claims awaiting complete fresh
verification}. This is neither a disproof nor a report of missing files.

Write $r_P=r_c(\Pi_1)$ and set
\begin{align}
 H_\beta(z,w)
 &=\frac1\beta\exp\left\{\frac{\pi z\overline w}{\beta}
       -\frac{\pi(|z|^2+|w|^2)}{2\beta}\right\},
 &G_\beta&=\operatorname{DPP}(\beta H_\beta),
 \label{sb:kernel}\\
 L_\beta&=\frac{-\log(1-\beta)}\beta,
 &&0<\beta<1.\label{sb:L}
\end{align}
The kernel $H_\beta$ is a projection of intensity $\rho=\beta^{-1}$,
whereas $G_\beta$ has intensity one. Thus $G_\beta$ is ordinary
unit-intensity Ginibre independently thinned with retention $\beta$ and
then contracted spatially by $\sqrt\beta$. In particular $G_1=G$.

The archived conclusion is the squeeze
\begin{equation}
 \frac{r_P}{\sqrt{L_\beta}}
 \le r_c(G_\beta)
 \le \frac{r_P}{\sqrt{(1-\epsilon_\beta)L_\beta}},
 \qquad \epsilon_\beta=o(\beta),
 \label{sb:squeeze}
\end{equation}
for sufficiently small positive $\beta$. Its lower bound is the
independently substantiated consequence of the exact Poisson upper
domination proved in the core part of this report. Its upper bound is
the part requiring the additional geometry. The screenshot's formula
was
\begin{equation}
 r_c(G_\beta)=r_P(1-\beta/4)+O(\beta^{6/5}).
 \label{sb:old-asymptotic}
\end{equation}
The later note sharpens the claimed remainder, using the same proposed
estimate with a different cutoff, to
\begin{equation}
 r_c(G_\beta)=r_P(1-\beta/4)
    +O\bigl(\beta^{3/2}[\log(1/\beta)]^{3/2}\bigr).
 \label{sb:new-asymptotic}
\end{equation}
Neither statement would, by itself, settle the target at $\beta=1$.

\subsection{Finite events, conditioning, and an exact transport identity}

All grains in this section are full closed disks of radius $r$. For
$m<L$, let $F_{m,L,r}$ be the indicator that an occupied component meets
both $B(0,m)$ and $\{z:|z|\ge L\}$. Centers are allowed everywhere;
the graph is not clipped to an annulus. A center attaches to the inner
terminal exactly when $|z|\le m+r$, and to the outer terminal exactly
when $|z|\ge L-r$. A finite witnessing path can be stopped when it
first attaches to the outer terminal, so the event is determined by
centers in a bounded disk, which may be enlarged whenever a later local
construction requires it. The closed convention also specifies the
events at artificially inserted roots on a terminal boundary.

For a fixed $\beta$, let
\[
 X_{s,\nu}=\operatorname{DPP}(sH_\beta)\cup\Pi_\nu,
 \qquad 0\le s<1,\quad \nu\ge0,
\]
with independent, retained colors. Let $\mu_s$ be the DPP law and
$\mu_s^{!x}$ its reduced Palm law at $x$; the root is omitted from this
Palm configuration. For $s>0$, the rank-one mixture identity gives the
positive measure
\begin{equation}
 \Lambda_x=\mu_s-(1-s)\mu_s^{!x},
 \qquad 0\le\Lambda_x\le\mu_s.
 \label{sb:weighted-measure}
\end{equation}
Write $T_x=d\Lambda_x/d\mu_s$, so $0\le T_x\le1$. This is a measure
definition, not an identification with the projection onto an
infinite random span of coherent vectors. Restriction to an observation
window gives the corresponding bounded local density and the same
integrals for events observed there.

For any bounded measurable local $F$, define
\[
 C_F(s,\nu)=\mathbb EF(X_{s,\nu}),\qquad
 D_xF(\eta)=F(\eta\cup\{x\})-F(\eta).
\]
The exact formulas are
\begin{align}
 \partial_s C_F(s,\nu)
 &=\frac1\beta\int
     \mathbb E_{\mu_s^{!x},\Pi_\nu}D_xF\,dx,
 \label{sb:s-derivative}\\
 \partial_\nu C_F(s,\nu)
 &=\int\mathbb E_{\mu_s,\Pi_\nu}D_xF\,dx.
 \label{sb:nu-derivative}
\end{align}
Indeed, realize the DPP by thinning one process of kernel $H_\beta$.
Conditional on that process, only its finitely many points in the
active window $W$ affect $F$. The expectation is a polynomial in $s$;
its derivative is the sum of the add-one differences over those
points, with all other retention marks independently sampled. Its
absolute value is at most $\operatorname{osc}(F)N(W)$, whose expectation
is $\operatorname{osc}(F)|W|/\beta$. Dominated convergence and reduced
Campbell give \eqref{sb:s-derivative}. The analogous Poisson add-one
calculation gives \eqref{sb:nu-derivative}. The same couplings give
continuity of these derivatives, including the appropriate one-sided
derivatives at $s=0$ and $\nu=0$. This is a finite-window assertion; its
dominating constant is not claimed to be independent of $W$.

For increasing $F$, put
\[
 I_F=\int\mathbb E[D_xF]\,dx,
 \qquad J_F=\int\mathbb E[T_xD_xF]\,dx.
\]
The expectations include the independent Poisson color. Substituting
\eqref{sb:weighted-measure} into \eqref{sb:s-derivative} gives the exact
identity
\begin{equation}
 \left(\partial_s-\frac{1}{\beta(1-s)}\partial_\nu\right)C_F
       =-\frac{J_F}{\beta(1-s)}.
 \label{sb:transport}
\end{equation}
Both $I_F$ and $J_F$ are nonnegative. Without any geometric estimate,
\eqref{sb:transport} has only the sign yielding the Poisson upper
domination. A gain in the opposite direction for connection events
requires a useful upper bound on $J_F/I_F$.

\subsection{The exact additional estimate and its parameter regime}

The principal input asserted in \path{direct_pivotal_void_split.md} is
the following. Fix
\begin{equation}
 0<r_-\le r\le r_+,\qquad
 m\ge200000r_+,\qquad L-m>240r_+.
 \label{sb:geometry-regime}
\end{equation}
Fix also $c_*>0$. For sufficiently small $\beta_0<1/2$, the proposed
estimate is uniform for
\begin{equation}
 0<s\le\beta\le\beta_0,\quad 0\le\nu\le2,
 \quad \max\{s/\beta,\nu\}\ge c_*,
 \quad 0<d<r_-/100.
 \label{sb:parameter-regime}
\end{equation}
For $F=F_{m,L,r}$ it states
\begin{equation}
 J_F\le\epsilon_\beta(d)I_F,
 \qquad
 \epsilon_\beta(d)=C d^3+C\exp\left\{-\frac{\pi d^2}{8\beta}\right\}.
 \tag{G}\label{sb:geometric-claim}
\end{equation}
The constant $C$ is asserted to depend only on the fixed radius and
intensity bounds, not on $m,L,d,s,\nu,\beta$ in the displayed ranges.
In particular, a proof with a constant growing in $L$, or one requiring
$\nu\ge\nu_0>0$ throughout, does not establish (G). Independence of
$C$ from $d$ is needed for the optimized cutoff yielding
\eqref{sb:new-asymptotic}.

The next subsections reconstruct the proposed proof of (G). They retain
its conditioning and repairs, but do not certify that its entire
geometric chain has been freshly proved. The later transport and limit
argument is a rigorous implication \emph{if (G) holds as stated}.

\subsection{The direct small-ball split}

Write $f(x)=\mathbb P(D_xF=1)$ and $I=\int f(x)\,dx$. Set
\[
 B_x=B(x,d),\qquad
 S_x(d)=\{w:2r-d<|w-x|\le2r\}.
\]
Suppose $x$ is insertion-pivotal and an original center $z$ lies in
$B_x$. Take a shortest positive terminal path after inserting $x$.
It uses $x$. If both neighbors of $x$ on that path are old centers,
at least one of them, say $w$, is not adjacent to $z$: otherwise
replacing $x$ by $z$ would produce a crossing in the old configuration.
Consequently $w\in S_x(d)$. If a neighbor is a terminal, the same
replacement works unless $x$ is attached to that terminal and $z$ is
not. Such a membership change forces $x$ into a deterministic strip of
width $d$ around the corresponding center-attachment circle.

Let $\mathcal T_d$ be the union of these two strips, and let $N_D$ and
$N_X$ count the DPP color and the full union. The deterministic
implication needed by the argument is therefore
\begin{equation}
 \{D_xF=1, N_D(B_x)>0\}
 \subset
 \{N_D(B_x)N_X(S_x(d))\ge1\}
       \cup\{x\in\mathcal T_d, N_D(B_x)>0\}.
 \label{sb:near-implication}
\end{equation}
The two random count regions are disjoint, with
\begin{equation}
 |B_x|=\pi d^2,
 \qquad |S_x(d)|=\pi(4rd-d^2)\le4\pi rd.
 \label{sb:areas}
\end{equation}

The current-law reduced Campbell bounds have coefficients
\[
 M_D=\frac{s}{\beta(1-s)},\qquad M=M_D+\nu.
\]
They are bounded uniformly in \eqref{sb:parameter-regime}. If $H_x$ is
an exterior event whose observed region excludes both $B_x$ and
$S_x(d)$, deleting the counted points does not change $H_x$. Applying
the colored Campbell upper bounds successively gives
\begin{align}
 \mathbb E[\mathbf1_{H_x}N_D(B_x)N_X(S_x(d))]
   &\le M_D M |B_x||S_x(d)|\mathbb P(H_x),
 \label{sb:two-count}\\
 \mathbb E[\mathbf1_{H_x}N_D(B_x)]
   &\le M_D|B_x|\mathbb P(H_x).
 \label{sb:one-count}
\end{align}
These estimates retain the exterior event. Replacing it by an
independent event would be unjustified for a DPP.

In the bulk, the proposed local separation lemma supplies an event
$H_x$, observed outside a fixed multiple of $r$, with
\begin{equation}
 \{D_xF=1\}\subset H_x,
 \qquad \mathbb P(H_x)\le C f(x).
 \label{sb:bulk-comparison}
\end{equation}
Then \eqref{sb:two-count} costs $O(d^3)f(x)$. At boundary roots,
the claimed estimates are an integrated replacement of
\eqref{sb:bulk-comparison} and the strip estimate
\begin{equation}
 \int_{\mathcal T_d}\mathbb P(H_x)\,dx\le Cd I.
 \label{sb:strip}
\end{equation}
Multiplying \eqref{sb:strip} by the $O(d^2)$ factor in
\eqref{sb:one-count} recovers the third power of $d$. With the
relocation described below, this yields the proposed conclusion
\begin{equation}
 \int\mathbb P(D_xF=1, N_D(B_x)>0)\,dx\le Cd^3 I.
 \label{sb:near-integrated}
\end{equation}
It is the geometric comparisons, not the elementary count calculation,
that require the substantial uniform argument.

\subsection{Canonical certificates and the empty-ball estimate}

The complementary event is
$E_x=\{D_xF=1, N_D(B_x)=0\}$. Choose a minimum-cardinality positive
path through the artificial root, with ties broken by an order
depending only on the candidate finite subset and its colors. Remove
the artificial root from that certificate. The remaining positive
requirements are paths to the two fixed terminals; the negative
requirements prohibit an old crossing and a preceding certificate.
Once a certificate is forced, these remainder constraints are
decreasing under deletion. This would fail for an order depending on
unretained random data, or for changing component labels whose terminal
connections were not themselves retained.

For the DPP part $A$ of a forced colored certificate, put
\begin{align*}
 P_A&=H_\beta(\cdot,A)H_\beta(A,A)^{-1}H_\beta(A,\cdot),
 &Q_A&=H_\beta-P_A,\\
 p_A(x)&=P_A(x,x)/\rho,
 &q_A(x)&=1-p_A(x).
\end{align*}
The reduced Campbell disintegration of \eqref{sb:weighted-measure}
is the measure identity
\begin{equation}
 s^k\det H_\beta(A,A)
 \left[p_A(x)\mu_A+q_A(x)\Lambda_{A,x}\right],
 \label{sb:campbell}
\end{equation}
where $\mu_A=\operatorname{DPP}(sQ_A)$ and
$\Lambda_{A,x}=\mu_A-(1-s)\mu_A^{!x}$. It follows from the determinant
ratio
$\rho_k^{!x}(A)/\rho_k(A)=q_A(x)$ and commutation of finite reduced
Palm operations. It is not an evaluation of an arbitrary almost
everywhere version of $T_x$ on a fixed Palm manifold. For $k$ DPP and
$j$ Poisson certificate points, the unordered coefficient is
\[
 \frac{s^k\det H_\beta(A,A)\nu^j}{k!j!}.
\]
Each configuration contributes its one selected colored subset, so the
sum has nonnegative integrands and Tonelli applies. At $\nu=0$ only
$j=0$ remains. Poisson certificate vertices may lie inside $B_x$; only
the DPP part must avoid it, and only that part enters $P_A$.

The void improvement can be calculated without any geometric surgery.
For a projection $Q$, let $\mu=\operatorname{DPP}(sQ)$, fix a root
with $Q(x,x)>0$, and write
\[
 u=\frac{Q(\cdot,x)}{\sqrt{Q(x,x)}},\quad
 K=sQ,\quad D=I-K_B,\quad t=\|\mathbf1_{B^c}u\|_2^2,
 \quad a_0=1+s\langle u_B,D^{-1}u_B\rangle.
\]
The determinant lemma gives
\[
 \frac{\mu^{!x}(N(B)=0)}{\mu(N(B)=0)}=a_0.
\]
The two void-conditioned exterior kernels satisfy
\begin{align*}
 K^{\rm void}&=K_{B^c}+K_{B^cB}D^{-1}K_{BB^c},\\
 (K^{!x})^{\rm void}&=K^{\rm void}-\frac{s}{a_0}vv^*,
 &v&=u_{B^c}+K_{B^cB}D^{-1}u_B.
\end{align*}
The second formula is Sherman--Morrison applied to the first. Kernel
order gives stochastic order of the conditional exterior DPPs. Since
$Q_B\ge u_Bu_B^*$, inverse monotonicity also gives
\[
 a_0\ge\frac1{1-s+st}.
\]
Thus for a decreasing exterior event $E_{B^c}$ and
$E=E_{B^c}\cap\{N(B)=0\}$,
\begin{equation}
 \Lambda_x(E)\le\frac{st}{1-s+st}\mu(E)
                  \le\frac{st}{1-s}\mu(E).
 \label{sb:void-improvement}
\end{equation}
The dependence between the exterior event and the void is retained
through their conditional laws.

For the remainder $Q_A$, its normalized row tail obeys
\begin{equation}
 q_A(x)t_A(x)\le2\bigl(h_B(x)+p_A(x)\bigr),
 \qquad
 h_B(x)=\rho^{-1}\int_{B^c}|H_\beta(x,y)|^2dy.
 \label{sb:row-tail}
\end{equation}
This follows by writing $Q_A=H_\beta-P_A$ and using
$\int|P_A(x,y)|^2dy=P_A(x,x)$. For $B=B(x,d)$,
$h_B(x)=e^{-\pi d^2/\beta}$. Combining
\eqref{sb:campbell}--\eqref{sb:row-tail} bounds each normalized
certificate contribution by
\begin{equation}
 p_A(x)+\frac{2s}{1-s}\bigl(e^{-\pi d^2/\beta}+p_A(x)\bigr).
 \label{sb:certificate-contribution}
\end{equation}

The analytic input recorded in \path{sparse_fock_projection.md} is
\begin{equation}
 p_A(x)\le C_k\exp\left\{-\frac{\pi\operatorname{dist}(x,A)^2}
                                      {8\beta}\right\},
 \label{sb:sparse}
\end{equation}
uniformly over finite sets with a fixed local clique or grid quota,
including colliding clusters of distinct points and arbitrarily large
total cardinality. The mechanism deserves emphasis. A quota forces
small physical clusters to have uniformly bounded cardinality.
Within each such cluster, confluent divided differences provide a
collision-uniform coherent-state basis. A compactness argument is used
only for this bounded number of nodes. Each cluster is then
orthonormalized; Gaussian decay makes the off-diagonal blocks of the
global Gram matrix have row sums below $1/2$ for small $\beta$.
The inverse is consequently bounded independently of the number of
clusters. Gaussian packing sums give \eqref{sb:sparse}. Applying a
Gram inverse bound to the original unnormalized near-collision matrix
would be invalid; the within-cluster normalization is essential.

A shortest two-path certificate has at most four points in each
radius-$r$ clique, so its DPP subset satisfies the needed quota. Thus
the claimed application of \eqref{sb:sparse} to
\eqref{sb:certificate-contribution} gives
\begin{equation}
 \mathbb E[T_x\mathbf1_{E_x}]
   \le C e^{-\pi d^2/(8\beta)}f(x).
 \label{sb:void-result}
\end{equation}
The source notes for this part are
\path{canonical_certificate_audit.md} and
\path{canonical_void_clearance.md}. Combining
\eqref{sb:near-integrated} and \eqref{sb:void-result}, using
$0\le T_x\le1$, is precisely the proposed proof of (G).

\subsection{How the proposed construction reaches zero bath intensity}

The positive-bath argument cannot simply be continued to $\nu=0$.
The archive uses two different filling constructions to cover the
parameter path.

\paragraph{A Poisson bath bounded away from zero.}
For a bounded active window, the conditional DPP $L$-kernel inside a
cleared region $B$, given the exact colored exterior in that window,
is a Schur complement. It satisfies
\[
 0\le L_B^\xi\le L_{BB},\qquad L_B^\xi(z,z)\le M_D.
\]
Hence its conditional void probability is at least
$e^{-M_D|B|}$, and its conditional factorial densities are bounded by
$M_D^n$. For any event $E$ of the interior union and the fixed exterior,
including a nonmonotone arm event, the bounds imply
\begin{equation}
 e^{-M_D|B|}\mathbb P_{\Pi_\nu,B}(E)
 \le\mathbb P_X(E\mid\text{colored exterior})
 \le e^{M_D|B|}\mathbb P_{\Pi_{\nu+M_D},B}(E).
 \label{sb:poisson-transfer}
\end{equation}
If $E$ has an interior certificate of at most $K$ vertices such that
arbitrary deletion of other vertices preserves it, independent thinning
gives
\[
 \mathbb P_{\Pi_v,B}(E)\le(v/\nu)^K
                     \mathbb P_{\Pi_\nu,B}(E),\qquad v\ge\nu>0.
\]
Shortest paths in a fixed bounded annulus give such a $K$ by a clique
cover. The recorded safe quotas are $K_0=4\cdot116^2$ in the bulk and
$K_1=2\cdot200^2$ at a boundary. Clearing the DPP and using separated
Poisson paths therefore transfers the ordinary Poisson separation
construction to the mixed process. The constants contain
$(\nu/(\nu+M_D))^{K_i}$ and Poisson filling probabilities; they
degenerate at $\nu=0$. The published bulk Poisson separation input is
Franceschetti--Penrose--Rosoman, Lemma 4.1, used at graph range one on
$C_{29}\setminus C_{20}$ with arbitrary fixed exterior vertex groups.
Its archived application and explicit inner filling are in
\path{cooperative_pair_bulk.md}; the mixed transfer is in
\path{cooperative_pair_dpp_bath.md}. This report does not elevate that
whole application to a new external certification.

\paragraph{A fixed-region determinantal filling.}
For $s/\beta$ bounded below, the archive replaces the Poisson filling by
a DPP-only filling on a favorable set of ordinary exterior
configurations. The construction in
\path{finite_window_leverage_operator.md} explicitly avoids a
conditional law given an arbitrary exact infinite exterior.

For a bounded active window $W$, the rank-one mixture and Janossy
identities construct a positive operator $0\le B_W(\xi)\le I$ on the
Bargmann range. Its coherent diagonal is the simultaneous continuous
version
\[
 T_W(x,\xi)=\langle f_x,B_W(\xi)f_x\rangle,
 \qquad f_x=H_\beta(\cdot,x)/\sqrt\rho.
\]
Positivity is first obtained simultaneously on a countable dense set
of Fock vectors, and then extended by continuity, so that one
configuration null set suffices. Magnetic translation and the
analytic submean inequality give
\begin{equation}
 \sup_{x\in K}T_W(x,\xi)
 \le\frac{e^\pi}{\pi\beta}
       \int_{K+B(0,\sqrt\beta)}T_W(w,\xi)\,dw.
 \label{sb:submean}
\end{equation}
Here $K$ is a fixed compact target region inside the cleared bounded
Borel region $B$.

Evaluate $B_W$ at the exterior configuration with the interior empty.
For $a=s/(1-s)$, the exact conditional interior kernel is
\begin{equation}
 L_{\rm cond}(x,y)=a\bigl[H_\beta(x,y)-B_W[H_\beta](x,y)\bigr].
 \label{sb:conditional-L}
\end{equation}
In particular its diagonal is $a\rho(1-T_W(x))$, while positivity
bounds the correction off the diagonal by
$a\rho\sqrt{T_W(x)T_W(y)}$. Independently,
$L_{\rm cond}\le L_W\le aH_{\beta,W}$ supplies the uniform combined
void lower bound $e^{-(a\rho+\nu)|B|}$. This lower bound permits transfer
from the distribution with an empty interior to the ordinary exterior
distribution without conditioning on a null set.

Suppose $H_{\rm ext}$ has intrinsic sparse certificates outside $B$
and a decreasing remainder. If
$b=\operatorname{dist}(K,B^c)>0$, the certificate void bound and
\eqref{sb:submean} yield the recorded estimate
\begin{equation}
 \frac{\mathbb P(H_{\rm ext},\ \sup_K T_W>\delta)}
      {\mathbb P(H_{\rm ext})}
 \le e^{(2+\nu_1)|B|}\delta^{-1}
       \frac{e^\pi |B|C'}{\pi\beta}
           e^{-\pi b^2/(32\beta)}.
 \label{sb:good-exterior}
\end{equation}
The coefficient tends to zero uniformly in the active window. Choose
$\delta=1/(8n_0)$ and then $\beta$ small. If at most $n_0$ target sets
$V_i\subset K$ have mutual distance at least $b_0>0$, the Gaussian
off-diagonal part is also at most $1/(8n_0)$. Gershgorin applied to
\eqref{sb:conditional-L} gives the conditional determinant lower bound
$(a\rho/2)^n$. Thus, on a good exterior, exactly one DPP point in each
$V_i$, no other DPP point in $B$, and no Poisson point in $B$ has
conditional probability at least
\begin{equation}
 e^{-(2+\nu_1)|B|}(c_*/2)^n\prod_{i=1}^n|V_i|,
 \qquad s/\beta\ge c_*.
 \label{sb:scaffold-probability}
\end{equation}
The $n!$ labelings cancel the unordered Janossy factor. Target sets may
depend measurably on the exterior because the good-exterior bound is a
supremum over the entire fixed compact $K$. A deterministic lower
bound on target areas is still necessary to turn this into a fixed
positive multiple of $\mathbb P(H_{\rm ext})$.

The region may have finitely many deterministic holes, provided its
area and the target depth are controlled. An arbitrary random stopping
region is \emph{not} authorized by this theorem. Section 8 of
\path{finite_window_leverage_operator.md} explicitly leaves that
extension unproved. The later geometry seeks to avoid needing it by a
finite family of deterministic punctured regions.

\subsection{The fixed-hole bulk geometry and its recorded repairs}

The bulk construction of \path{zero_bath_bulk_separation.md} translates
the root to zero and divides lengths by $2r$, so the graph range is
one. Both terminals lie outside $C_{29}$. Let $H$ be the event that
the two terminal groups have disjoint connections to
$C_{25}\setminus C_{24}$ using only centers outside $C_{24}$.
Pivotality implies $H$. The desired conclusion is
$\mathbb P(D_xF=1)\ge\kappa\mathbb P(H)$, uniformly in the outer
window and at zero Poisson bath.

A fixed grid of side $10^{-5}$ covers the compact annulus through
$C_{29}$. A pattern retains finitely many cells with assigned terminal
labels. Its cleared region is $C_{29}$ minus those retained cells,
\emph{clipped to $C_{24}^c$}. Every selected outer-path cell is required
to have a terminal-connected anchor outside $C_{25}$ whose connection
uses only vertices outside $C_{25}$. Entry cells separately retain
their prescribed connected entry points. These are actual corrections
recorded in \path{zero_bath_bulk_audit.md}:
\begin{itemize}
 \item Without the outside-$C_{25}$ path condition, an anchor connection
 could run through an inner entry, invalidating the radial gap used
 later.
 \item Without clipping at $C_{24}$, a retained cell could reveal extra
 points inside the disk excluded from the definition of $H$, so its
 exterior would no longer necessarily be a subconfiguration of the
 nonconnecting graph defining $H$.
\end{itemize}
With these corrections, each retained exterior is an ordinary fixed
colored exterior. There are only finitely many patterns. Their
positive requirements have shortest-path witnesses with a fixed local
quota, and their remaining nonconnection requirements are decreasing.
The target filling therefore seeks to apply
\eqref{sb:scaffold-probability} to a deterministic punctured region,
not to a configuration-dependent Markov domain.

The retained entry parents $y,z$ have radii between $24$ and $25$.
Label them so that $|y|\ge|z|$, and let $R_C$ be the supremum radius
of the fixed cell containing $y$. The split is:
\begin{itemize}
 \item In Case F, the opposite terminal has an outside-$C_{25}$
 connected point $z_0$ with $|y|\le|z_0|\le R_C+0.05$.
 Retain its exterior path in place of the inner opposite entry.
 \item In Case G, no such exterior connection reaches the fixed collar
 $C_{R_C+0.05}\setminus C_{25}$. This is a decreasing requirement
 after the parent cell has been fixed. Opposite outer-hole points then
 have radius at least $|y|+0.049$, while opposite inner-entry points
 lie within $0.001$ of $z$.
\end{itemize}

Here are the principal strict margins recorded in the construction.
In Case F the two parent radii differ by at most $0.051$. Parent
nonadjacency gives
\[
 \sin(\theta/2)>0.998/50.2.
\]
Move each branch $0.6$ inward and $0.1$ tangentially away from the
other, then continue to radial parameter $23$. Their vertical
separation is greater than
\[
 46(0.998)/50.2+0.14>1.05.
\]
For an opposite hole vertex with radius deficit at most $0.061$, the
radial-distance identity gives squared clearance greater than
$1.261$ before the tangential shift, hence clearance greater than
$\sqrt{1.261}-0.1>1.02$ afterward. Diverging circular arcs and
opposite radii then lead to endpoints at distance $0.75$ from the
root.

In Case G, parents more than three apart admit simple inward radial
paths, with cross-branch separation exceeding $2.6$. For closer
parents, rotate so that $y=(R,0)$, and choose the exterior predecessor
of $z$ on the upper side. The radial gap and the edge to that
predecessor force $z_2>0.5$. The first lower-branch target is constructed
by intersecting the unit circle around $y$ with the circle of radius
$R+0.049$, and then taking a point at distance $0.999$ from $y$ and
$1.01$ from the lower intersection. In the source notation it is
\begin{align*}
 U&=0.049+\frac{0.049^2-1}{2R},\quad V=\sqrt{1-U^2},\\
 t&=(0.999^2+1-1.01^2)/2,\quad q=\sqrt{0.999^2-t^2},\\
 b&=y+(tU-qV,-tV-qU).
\end{align*}
The listed bounds are
$b_1-R<-0.855$ and $-0.516<b_2<-0.513$.
The first upper-branch target is $c=z-0.999(1,0)$, with
$c_2>0.5$. A two-circle distance calculation is used to keep $b$
more than $1.01$ from the opposite outer holes; the vertical separation
keeps it more than $1.013$ from $z$. The corresponding calculation for
$c$ gives clearance exceeding $1.05$ from the opposite outer holes.
Horizontal corridors to coordinate $20$, followed by vertical
separation and diverging arcs, preserve these margins.

The final recorded discretization uses at most $1000$ targets. A
maximal separated subcollection retains connectivity of each branch;
disks of radius $10^{-4}$ around its centers have mutual separation
at least $0.0048$ and fixed depth at least $0.3$ in the punctured
region. Every connection or nonconnection survives displacement
inside these disks. The positive lower target area and the finite
number of patterns turn \eqref{sb:scaffold-probability} into the
claimed bulk comparison. These explicit margins explain the proposed
uniformity; complete fresh checking of all cases and event reductions
remains part of the unresolved verification of (G).

\subsection{Boundary construction, strips, and deep-root relocation}

At boundary roots, ordinary insertion pivotality can involve direct
terminal attachment. The archive therefore does not extend the bulk
two-arm lemma to the boundary by assertion. In normalized range-one
coordinates the nearby attachment circle has radius at least
$100000$. The construction in
\path{zero_bath_boundary_separation.md} uses the fixed radii
\[
 R_j=49+j/100,\qquad j=0,\ldots,99.
\]
Its event $H_j$ observes only the graph outside $C_{R_j}$ and asks for
a distant-terminal connection to the shell
$[R_j,R_j+0.01]$, together with terminal nonconnection. A witnessing
path first entering $C_{50}$ gives at least one such event. Each $H_j$
has a shortest-path certificate outside its deterministic cleared disk.

If $z$ is the distant-connected shell parent, $r_0=|z|$, and
$e=z/r_0$, then every nearby-terminal exterior vertex $v$ obeys
$|v|\ge r_0-0.01$ and $|z-v|>1$. The exact radial identity gives
\[
 |z-0.5e-v|^2
  =(1-0.5/r_0)|z-v|^2+0.25
       +(0.5/r_0)(|v|^2-r_0^2)>1.229.
\]
An allowed-side displacement of $0.05$ corrects the small penetration
of the large attachment circle. The recorded first-target clearance
is then greater than $\sqrt{1.229}-0.05>1.05$. A far branch follows
an allowed circular arc and a positive radius to $+0.75$; a separate
near branch ends at $-0.75$ and reaches the nearby terminal. For
normalized signed distance in $[-0.01,30]$, the construction makes
every insertion within $0.1$ of the original root pivotal. In physical
coordinates this supplies comparison to roots within $0.2r$.

The discretization keeps at most $1000$ targets, with target radius
$0.001$ and mutual separation at least $0.008$ in normalized units.
They have fixed positive depth in the cleared disk. The fixed-region
DPP filling applies to each of the $100$ bins; summing costs only a
fixed factor. On terminal strips, this robust comparison allows
averaging over nearby roots. Circle coarea, rather than a stationarity
assertion about a conditional exterior, gives \eqref{sb:strip}.

There is a further correction. A pivotal root can have signed
normalized distance as low as $-1$, whereas the unshifted robust
construction was justified only down to $-0.01$. For those deeper
roots, \path{direct_pivotal_void_split.md} moves the artificial root
to $+2$, replaces the branch endpoints by $+1.25$ and $+2.75$, and
extends the near branch to $-(a+0.25)$, where $a\ge-1$ is the signed
distance. The original root need not itself be made robustly pivotal.
Instead the claimed comparison is
\begin{equation}
 \mathbb P(H_x)\le C f\bigl(x+4r n(x)\bigr),
 \label{sb:relocation}
\end{equation}
where $n(x)$ points into the unattached side of the nearby terminal.
The recorded correction of the radial penetration bound is from
$0.002$ to less than $0.025$, still covered by the same $0.05$
allowed-side shift. The two branches retain separation $1.5$; both
endpoints are $0.75$ from the relocated root.

On either circular belt, the map in \eqref{sb:relocation} changes the
radial coordinate by $\pm4r$ and leaves the angle fixed. Its Jacobian
is $(|x|\pm4r)/|x|$, with a uniformly bounded inverse under
\eqref{sb:geometry-regime}; the two belt images have bounded
multiplicity. Thus integrating \eqref{sb:two-count} against
\eqref{sb:relocation} still costs $Cd^3I$. On the actual width-$d$
membership-change strips, the original unshifted comparison is in
its valid range and supplies the extra factor $d$. The explicit
recorded check is \path{zero_bath_boundary_audit.md}.

\subsection{The complete implication from (G) to an asymptotic}

The remainder of the argument has no geometric freedom once (G) is
available with its asserted uniform constants. For this application fix
$c_*=1/3$ in \eqref{sb:parameter-regime}; the characteristic below is
then covered for all sufficiently small $\beta$. Choose a cutoff $d=d_\beta$
so that $\epsilon_\beta=\epsilon_\beta(d_\beta)=o(\beta)$, and reduce
$\beta_0$ until $0\le\epsilon_\beta\le1/2$. Combining (G) with
\eqref{sb:transport} gives
\begin{equation}
 \left(\partial_s-
  \frac{1-\epsilon_\beta}{\beta(1-s)}\partial_\nu\right)C_F
   =\frac{\epsilon_\beta I_F-J_F}{\beta(1-s)}\ge0.
 \label{sb:reverse-transport}
\end{equation}
For $0\le s\le\beta$, define
\begin{equation}
 \nu_s=\frac{1-\epsilon_\beta}{\beta}
           \log\frac{1-s}{1-\beta},\qquad
 \lambda_\beta=\nu_0=(1-\epsilon_\beta)L_\beta.
 \label{sb:characteristic}
\end{equation}
This path has nonnegative bath, ends at $\nu_\beta=0$, and has
$\nu_s'=-(1-\epsilon_\beta)/[\beta(1-s)]$. It stays in the claimed
parameter regime: for $s\le\beta/2$,
\[
 \nu_s\ge\frac{1-\epsilon_\beta}{\beta}
                  \int_{\beta/2}^{\beta}\frac{dt}{1-t}
       \ge\frac{1-\epsilon_\beta}{2},
\]
which is at least $1/3$ for sufficiently small $\beta$; for
$s\ge\beta/2$, the DPP intensity $s/\beta$ is at least $1/2$.
Also $\nu_s\le L_\beta\le2$ for $\beta\le1/2$.

At each fixed finite window, the continuously differentiable formulas
above permit the ordinary chain rule on every characteristic, not
merely on almost every curve. One-sided continuity supplies $s=0$
and $\nu=0$. Alternatively, an almost-everywhere formulation can be
straightened and mollified before taking the prescribed curve; the
source note \path{positive_bath_pde_comparison.md} explains that
procedure. No derivative is taken after an infinite-volume limit.
Integration of \eqref{sb:reverse-transport} therefore gives
\begin{equation}
 \mathbb P(F_{m,L,r}(G_\beta)=1)
 \ge\mathbb P(F_{m,L,r}(\Pi_{\lambda_\beta})=1)
 \label{sb:annulus-comparison}
\end{equation}
for every admissible $m,L,r$. It is a comparison of these connection
events, not lower stochastic domination for all increasing observables.

Fix the compact radius interval to contain $r_P$, for example
$[r_P/2,2r_P]$. For small $\beta$, the Poisson threshold
$r_P/\sqrt{\lambda_\beta}$ lies in this interval. Take
$r>r_P/\sqrt{\lambda_\beta}$ and one admissibly large finite seed
$B(0,m)$. Poisson spatial scaling makes the process on the right of
\eqref{sb:annulus-comparison} supercritical at this radius. The seed
has positive probability of meeting an unbounded occupied component:
otherwise stationarity and a countable cover by translates of that
seed would give zero probability of any unbounded component.

Now let $L\to\infty$ while holding $m,\beta,r$ fixed. The events
$F_{m,L,r}$ decrease. Local finiteness gives only finitely many grains,
and hence finitely many occupied components, meeting the seed. Thus
their decreasing intersection is exactly the event that the seed meets
an unbounded component; arbitrarily distant connections cannot be
provided by infinitely many different seed components. Continuity of
probability from above passes \eqref{sb:annulus-comparison} to this
event. It follows that $G_\beta$ percolates with positive probability.
Its stationarity and mixing identify this with the usual almost-sure
percolation conclusion. Letting $r$ decrease to the Poisson threshold
gives the upper bound in \eqref{sb:squeeze}. No assertion about
percolation at the critical Poisson radius is used.

For the two error choices, first take $d_\beta=\beta^{2/5}$. For small
$\beta$ it is in the allowed range, and
\[
 \epsilon_\beta\le C\beta^{6/5}
       +C\exp\{-(\pi/8)\beta^{-1/5}\}=O(\beta^{6/5}).
\]
The exponential is smaller than every fixed positive power of
$\beta$, as follows by taking logarithms. Next take
\[
 d_\beta^2=(16/\pi)\beta\log(1/\beta).
\]
Then the exponential term is exactly $\beta^2$, and
\begin{equation}
 \epsilon_\beta
  =O\bigl(\beta^{3/2}[\log(1/\beta)]^{3/2}\bigr).
 \label{sb:optimized-error}
\end{equation}
Both choices give $\epsilon_\beta=o(\beta)$.

Finally,
\[
 L_\beta=1+\frac\beta2+\frac{\beta^2}{3}+O(\beta^3),
 \qquad
 L_\beta^{-1/2}=1-\frac\beta4-\frac{7\beta^2}{96}+O(\beta^3).
\]
For $0\le\epsilon\le1/2$,
$(1-\epsilon)^{-1/2}-1\le2\epsilon$. Since $L_\beta\ge1$, the
width of the squeeze \eqref{sb:squeeze} is at most
$2r_P\epsilon_\beta$. Hence that squeeze rigorously implies
\[
 \left|r_c(G_\beta)-r_P(1-\beta/4)\right|
    \le C r_P(\beta^2+\epsilon_\beta).
\]
This proves the algebraic implications
\eqref{sb:old-asymptotic} and \eqref{sb:new-asymptotic} from the
respective proposed error estimates. The asserted right derivative
$-r_P/4$ would follow by this deterministic squeeze, with no
differentiation of an infinite-volume percolation event.

Strictness in this implication comes from a positive intensity gain,
not from an unquantified limiting comparison. Since
$\epsilon_\beta=o(\beta)$ and $L_\beta=1+\beta/2+O(\beta^2)$,
$\lambda_\beta>1$ for all sufficiently small positive $\beta$.
For an explicit conservative version, take
$\epsilon_\beta\le\beta/8$ and $\beta\le1/2$. The inequalities
$L_\beta\ge1+\beta/2$ and $L_\beta\le2$ give
$\lambda_\beta\ge1+\beta/4$, and therefore
\[
 r_c(G_\beta)\le\frac{r_P}{\sqrt{1+\beta/4}}<r_P,
 \qquad r_P-r_c(G_\beta)\ge r_P\beta/16.
\]
This last display, like the upper squeeze bound, remains dependent on
the complete verification of (G).

\subsection{A separate strict comparison with a positive Poisson bath}
\label{sb:mixed-gap}

The recovered result of \path{positive_bath_pde_comparison.md} concerns
a different unit-intensity endpoint. Fix $a\in(0,1)$, put $p=1-a$,
$b=\beta p$, and
\[
 Y_{\beta,a}=\operatorname{DPP}(bH_\beta)\cup\Pi_a.
\]
Its hypothesis is the uniform distributional differential estimate
\begin{equation}\label{sb:mixed-pde}
 \left(\partial_s-\frac1{\beta(1-s)}\partial_\nu\right)C
 \ge -K_r\frac{s}{\sqrt\beta}\partial_r C
       -K_\nu s\partial_\nu C.
\end{equation}
Here $C=C_{m,L}(s,\nu,r)$, the constants are independent of the
connection windows and small beta, and the estimate is required for
$s\le\beta$, $\nu\in[a/2,2]$, $r\in[r_P/2,2r_P]$ and the seed
regime \eqref{sb:geometry-regime}. Its geometric dependencies are the
positive-bath cooperative and current-law cofactor estimates; this
subsection records their consequence without certifying that entire
chain anew.

The characteristic ending at bath $a$ is
\[
 \nu_s=a+\int_s^b\left[\frac1{\beta(1-t)}-K_\nu t\right]dt,
 \qquad r_s=r_0+\frac{K_rs^2}{2\sqrt\beta}.
\]
Thus its initial intensity and radius cost are
\[
 \lambda_{\rm eff}=a-\beta^{-1}\log(1-\beta p)
                    -\tfrac12K_\nu\beta^2p^2,\qquad
 \Delta_r=\tfrac12K_rp^2\beta^{3/2}.
\]
Finite-window probabilities are locally Lipschitz in all parameters.
Straightening the characteristic and mollifying the transverse
variables justifies integration of the distributional inequality on
every prescribed admissible curve. For small beta the entire path
stays in the allowed domain, and
\[
 C_{m,L}(b,a,r_0+\Delta_r)
       \ge \mathbb P(F_{m,L,r_0}(\Pi_{\lambda_{\rm eff}})=1).
\]
Keep one admissible seed fixed and send $L\to\infty$.
The decreasing-event argument already used above gives
$r_c(Y_{\beta,a})\le r_P/\sqrt{\lambda_{\rm eff}}+\Delta_r$.
If $K_\nu\beta\le1/2$, then
$\lambda_{\rm eff}\ge1+p^2\beta/4$; if also
$\sqrt\beta\le r_P/(16K_r)$, with zero-constant constraints omitted,
the recovered quantitative conclusion is
\begin{equation}\label{sb:mixed-strict}
 r_c(Y_{\beta,a})\le r_P-\frac{r_P(1-a)^2\beta}{32}<r_P.
\end{equation}
This conditional implication is a positive-bath theorem at each fixed
$a$, not a uniform removal of the bath and not a comparison at
$\beta=1$. The recorded checks are \path{weak_hybrid_threshold_audit.md}
and \path{mixed_cooperative_audit.md}.

\subsection{What this recovered argument does and does not settle}

The archive contains the proof chain, its repairs, and its claimed
audits. It is inaccurate to describe the small-$\beta$ argument as
lost or disproved. It is also inaccurate to treat its audit labels as
a substitute for checking the uniform geometric statement (G) in full.
The exact unresolved verification is whether all fixed-hole pattern
events, boundary targets, and conditional filling comparisons provide
the constants in (G) uniformly in the growing connection windows and
through the zero-bath endpoint. The transport calculation and its
infinite-volume consequence have been separated so that this question
is not hidden in a derivative of $r_c$.

Even a completely verified small-$\beta$ result is a separate statement
from the original full-Ginibre target. Ordinary thinning gives only
\[
 r_c(G)\le\frac{r_c(G_\beta)}{\sqrt\beta},
\]
which does not turn a small improvement of $r_c(G_\beta)$ over $r_P$
into an improvement for $G$. No monotonicity of $r_c(G_\beta)$ in
$\beta$ has been proved in these notes. Thus continuing from small
$\beta$ to $\beta=1$ remains an additional mathematical task, even if
the archival geometric claim is confirmed.

\subsection{Sharp obstructions to coupling the small-beta endpoint to Ginibre}
\label{sb:coupling-barriers}

The recovered calculations in
\path{research/ginibre_coupling_bridge_obstructions.md} rule out
several stronger ways to transfer the small-beta gain. Write
$G_{\beta,\lambda}$ for intensity-$\lambda$ beta-Ginibre, with kernel
\[
 K_{\beta,\lambda}(z,w)=\lambda
  \exp\left\{\frac{\pi\lambda}{\beta}z\bar w
     -\frac{\pi\lambda}{2\beta}(|z|^2+|w|^2)\right\}.
\]
If a coupling gives $G_{\beta,\lambda}\subset G$, products of counts
in disjoint sets order the continuous factorial densities. At distinct
fixed $a_1,\ldots,a_n$, the least-degree Vandermonde term gives
\[
 \lim_{\epsilon\downarrow0}
 \frac{\rho_n^{\beta,\lambda}(\epsilon a_1,\ldots,\epsilon a_n)}
      {\rho_n^{1,1}(\epsilon a_1,\ldots,\epsilon a_n)}
 =\lambda^{n(n+1)/2}\beta^{-n(n-1)/2}.
\]
This must be at most one for every $n$. First taking the coalescence
limit at fixed $n$, then $n\to\infty$, proves the necessary condition
\begin{equation}\label{sb:inclusion-barrier}
 \lambda\le\beta.
\end{equation}
Equality is realized by ordinary independent beta-thinning of $G$.
Dependent thinning with an exact beta-Ginibre output cannot improve
this intensity loss.

The same restriction holds for an injection with any fixed finite
displacement bound $D$. The count in a disk for a kernel $pH_\gamma$
is a sum of independent Bernoulli variables with parameters
$p\mu_j(u)$, where
$u=\pi R^2/\gamma$ and
$\mu_j(u)=\mathbb P(\operatorname{Poisson}(u)\ge j+1)$.
This is independence of count summands, not spatial points.
The bounds
$e^{-u}u^{j+1}/(j+1)!\le\mu_j(u)\le u^{j+1}/(j+1)!$,
applied to the first $n$ successes and to $\mathbb E\binom Nn$,
give the fixed-parameter expansion
\begin{equation}\label{sb:overcrowding}
 \log\mathbb P(N\ge n)
 =-\tfrac12n^2\log n+
          (\tfrac34+\tfrac12\log u)n^2+O_{u,p}(n\log n).
\end{equation}
For the upper bound, writing indices as $j-1+\ell_j$ bounds the
remaining sum by a product
$\prod_{j=1}^n\sum_{\ell\ge0}u^\ell j!/(j+\ell)!
\le C_u n^{2u}$; summing Stirling's formula supplies the expansion.
An injection would order the source count in $B_R$ by the Ginibre
count in $B_{R+D}$. If $\lambda>\beta$, choose a fixed $R$ such that
$\lambda R^2/\beta>(R+D)^2$. Equation~\eqref{sb:overcrowding} then
contradicts this order as $n\to\infty$.

Allowing collisions does not permit an arbitrarily small-displacement
covering at nearly unit intensity. Assume only every source point is
within $D$ of a Ginibre point, with $0<\lambda\le1$.
Translation averaging makes an admissible coupling stationary.
Assign each source point to a nearest Ginibre point and retain one
assignment per used target. The discarded source density is at most
$2\pi\lambda^2D^2$, by the source pair-density bound $\rho_2\le\lambda^2$;
mass transport bounds unused target density by
$1-\lambda+2\pi\lambda^2D^2$.
Comparing holes in a disk of radius $1/\sqrt\pi$, and paying for its
boundary strip, gives, for $\beta\le1/10$,
\begin{equation}\label{sb:cover-barrier}
 \delta_0:=e^{-10/9}-2e^{-2}
       \le1-\lambda+2\pi D^2+2\sqrt\pi D.
\end{equation}
The source hole is at least $e^{-\lambda/(1-\beta)}$; the first two
Ginibre disk eigenvalues bound its hole above by $2e^{-2}$.
Thus $\delta_0>0$, and at $\lambda=1$,
$D\ge(\sqrt{1+2\delta_0}-1)/(2\sqrt\pi)$.
When a small-beta radius improvement is of order beta, a useful
transfer with $\lambda\le1$ would require $\lambda\to1$ and
$D=O(\beta)$, contradicting \eqref{sb:cover-barrier}.

Two related decomposition shortcuts also fail. An independent union
of positive-intensity components has a nonzero cross-color pair density
at coincidence, unlike Ginibre's vanishing pair density. Orthogonal
spectral subspaces therefore do not in general yield independent
spatial subprocesses. The known independent-block identity after a
finite Ginibre power map concerns the powered process; its unbounded
spatial derivative does not transfer bounded disk grains. None of
these obstructions disproves the critical-radius conjecture; they
close the specified exact-output coupling routes.

% ===== END small_beta =====

% ===== BEGIN branches =====
\section{Attempted routes, exact obstructions, and recovered calculations}
\label{br:investigation}

This section records how the investigation developed beyond the principal
results proved elsewhere in this report.  It includes the earlier research
history supplied by the user, not just the later working notes.  Its purpose
is to preserve the mathematics and the reasons for changing direction.
It is not a second list of established theorems.

Three distinctions govern the presentation.  An elementary identity or
counterexample whose argument is reproduced below can be checked from that
argument.  A \emph{recovered calculation} is an assertion with a substantial
derivation in the preserved notes; its historical ``audited'' label is not
being adopted as a fresh certification of every premise.  An explicitly
\emph{unproved comparison} remains a missing estimate, even when several
inputs to it have been established.  None of the negative examples below
disproves the original Ginibre--Poisson critical-radius conjecture.

Throughout, $r$ is the grain radius and $R=2r$ is the center connection
range.  Write $H(z,w)=\exp\{\pi z\overline w-
\pi(|z|^2+|w|^2)/2\}$ for the unit-intensity Ginibre projection.
For a finite tuple $A$ of distinct points,
\[
 P_A=H(\cdot,A)H(A,A)^{-1}H(A,\cdot),\qquad Q_A=H-P_A.
\]
For an increasing Boolean event $F$, $D_xF(\eta)=F(\eta\cup\{x\})-F(\eta)$.
The words ``root'', ``parent'', and ``mark'' always refer to point
locations; a reduced Palm root is restored only when this is expressly
stated.

\subsection{What kind of advance would reach the target?}
\label{br:sufficient}

An exact local improvement is useful only if its normalization and its
scale range permit an infinite-volume conclusion.  For example, consider
\[
 C_{m,L}(s,\nu,r)=
 \mathbb P\bigl(\operatorname{DPP}(sH)\cup\Pi_\nu
 \text{ connects }B_m\text{ to }\partial B_L
 \text{ with radius }r\bigr).
\]
If one could prove, uniformly in appropriate $m,L$, a differential
comparison
\begin{equation}
 \partial_s C_{m,L}-a(s)\partial_\nu C_{m,L}\geq 0
 \label{br:characteristic}
\end{equation}
along the actual path
\[
 \nu(s)=\int_s^1 a(t)\\,dt\geq0,
 \qquad \lambda_*:=\int_0^1 a(t)\\,dt>1,
\]
then the endpoint comparison would be
$C_{m,L}(1,0,r)\geq C_{m,L}(0,\lambda_*,r)$.
After valid endpoint and infinite-volume passages, Poisson scaling would
give
\[
 r_c(G)\leq \frac{r_c(\Pi_1)}{\sqrt{\lambda_*}}
 <r_c(\Pi_1).
\]
Here the strictness is the definite excess $\lambda_*-1$, not merely a
strict finite-window probability inequality.  A version allowing a
radius increase must pay that increase from the same excess.

The exact local derivative can be written using the reduced Palm
likelihood ratio $R_x^{(s)}$:
\[
 \partial_s C=\int\mathbb E[R_x^{(s)}D_xF]\,dx,
 \qquad \partial_\nu C=\int\mathbb E[D_xF]\,dx.
\]
For $s<1$, let $T_x=1-(1-s)R_x^{(s)}$.  If the pivotal integral is
positive and
\[
 \ell(s)=
 \frac{\int\mathbb E[T_xD_xF]\,dx}
      {\int\mathbb E[D_xF]\,dx},
\]
the relevant normalized coefficient is $(1-\ell(s))/(1-s)$.  The desired
gain is therefore related to the improper integral
\begin{equation}
 \int_0^1\frac{s-\ell(s)}{1-s}\\,ds>0.
 \label{br:averaged-gap}
\end{equation}
This display does not assert convergence or positivity.  The dependence
of $\ell$ on the current bath, radius, and window is part of the problem.
Several attempted proofs failed precisely by suppressing one of these
dependencies.

\subsection{Voids, negative association, and the unknown vacant path}
\label{br:void}

The most immediate determinantal advantage is the smaller probability of
a prescribed hole.  For a bounded Borel set $A$, the Fredholm expansion
gives
\begin{equation}
 -\log\mathbb P(G(A)=0)
 \geq |A|+\frac12\iint_{A\times A}e^{-\pi|x-y|^2}\\,dx\\,dy.
 \label{br:hole-second}
\end{equation}
The strict positive second term suggested comparing vacant crossings.
A vacant curve $\gamma$ requires an empty radius-$r$ neighborhood
$\gamma+B_r$.  Thus \eqref{br:hole-second} does improve the cost of every
\emph{fixed} candidate curve.

The missing operation is taking the union over all possible curves.
Comparisons of each event, or of their finite intersections, do not
compare their union because inclusion--exclusion has alternating signs.
Nor is the length of a curve a lower bound for the area of its tube:
repeated traversals and closely spaced turns invalidate that substitution.
The earlier Voronoi route made this problem concrete.  A Voronoi vertex
has degree three in general position, but very short edges allow many
successive vertices in a small region.  A path count in a realized
Voronoi graph does not bound the integral over its random nuclei.

For forced nuclei $X=\{x_1,\ldots,x_k\}$, the correct weight is
\[
 \det H(X,X)\,
 \mathbb P^{!X}(G(A)=0).
\]
Rank-$k$ eigenvalue interlacing bounds the second factor after discarding
the first $k$ eigenvalues of $H_A$.  It does not permit discarding the
factor $\det H(X,X)$ when integrating nucleus positions.  The archive's
straight-strip calculation supplies the transverse spectral values
\[
 \lambda_w(a)=\sqrt2\int_{-w/2}^{w/2}e^{-2\pi(x-a)^2}\\,dx
\]
and a negative log-hole rate per longitudinal unit length
$\int_{\mathbb R}-\log(1-\lambda_w(a))\,da$.  Its use for arbitrary
curved, overlapping tubes would require geometric and spectral estimates
not supplied by that calculation.

Conditioning only on an empty region $A$ is favorable for further
emptiness estimates: the exterior conditional kernel is
\[
 H^A=H_{A^c,A^c}+H_{A^c,A}(I-H_A)^{-1}H_{A,A^c}
 \succeq H_{A^c,A^c}.
\]
An exploration also finds occupied points.  Palm conditioning subtracts
a coherent direction and can reverse the desired hole comparison.  The
supplied history records a boundary-root disk example where
$\mathbb P^{!x}(G(B_s)=0)>e^{-\pi s^2}$ for
$1\leq\pi s^2\leq11/10$.  This is a recovered counterexample from that
history; its original detailed note is not being replaced by an assertion
of uniform exploratory domination here.

Negative association has a valid but narrower role.  It bounds products
of increasing functions, or products of decreasing functions, on
\emph{disjoint} spatial supports.  It can supply a block criterion once a
single good-block probability is sufficiently large.  It does not order
Ginibre and Poisson crossing probabilities.  Even the two-site projection
DPP which chooses exactly one of two sites has a larger probability of
``at least one'' and a smaller probability of ``both'' than independent
Bernoulli-$1/2$ occupancy.  This also explains why directionally convex
count orders, factorial-moment bounds, Laplace-functional orders, and
negative association cannot be used as crossing orders without an
additional argument.

The history contains successive finite combinatorial refinements of a
block criterion, culminating in the recorded sufficient condition
\[
 g_{L,2L}^{\rm in}(r)+2g_{L,L}^{\rm in}(r)>23/8.
\]
Its use at $r=r_c(\Pi_1)$ would establish strictness by continuity of the
finite-window probabilities.  No required Ginibre crossing probability
was proved to satisfy it.  The improved finite criterion and its original
finite checks are historical inputs, not fresh certifications in this
section.

\subsection{Affine deformations: the correct limit and the wrong sign claim}
\label{br:affine}

The ellipse calculations and the justified Poisson-limit argument are
treated separately in this report.  Two additional obstructions belong
to the general investigation.

First, isotropy does not imply that disks minimize a critical scale over
equal-area ellipses.  The archive gives a direct counterexample using a
Poisson process of parents, each carrying an independently rotated finite
$\delta$-net of the unit disk.  Let
\[
 \lambda=(100\pi)^{-1},\quad a=6400\pi,
 \quad b=a^{-1},\quad\delta=b/10,
 \quad E=\{x_1^2/a^2+x_2^2/b^2\leq1\}.
\]
Every radius-two offspring disk is contained in its parent's radius-three
disk.  The parent connection path bound is
$36\pi\lambda=9/25<1$, proving nonpercolation of these radius-two
offspring disks.  Conversely, the ellipses centered at a cloud's net
points cover the rectangle
$[-a/2,a/2]\times[-1/2,1/2]$ about its parent.  Indeed, choose a cloud
point within $\delta$ of $(0,x_2)$; the ellipse quadratic form is at most
$(3/4)^2+(1/4)^2<1$.

An area-preserving linear map turns this parent rectangle into a square
of side $\sqrt a$.  Squares of side $\sqrt a/2$ contain independent parent
points with vacancy probability $e^{-16}$, and nearest-neighbor occupied
squares have overlapping grains.  An elementary closed-circuit sum is
less than one, hence these ellipse grains percolate.  After deterministic
rescaling to intensity one the strict shape comparison persists.  The
process is stationary, isotropic and mixing because the clouds have
bounded radius.  It is clustered; this is not a Ginibre counterexample.

Second, an anisotropic approach to Poisson does not improve all crossings
in a common direction.  Put $\varepsilon=e^{-t}$ and
$\Phi_t=S_t^{-1}G$ with $S_t=\operatorname{diag}(e^t,e^{-t})$.
For a bounded observable $F$ supported in a fixed window $W$, the
recovered determinantal difference expansion is
\begin{equation}
 \mathbb EF(\Phi_t)=\mathbb EF(\Pi_1)
 -\frac12\iint |K_t(x,y)|^2\mathbb ED_xD_yF(\Pi_1)\,dx\,dy
 +O_W(\varepsilon^2).
 \label{br:anisotropic-expansion}
\end{equation}
Only one coordinate collapses.  For disk crossings its leading term is
\[
 -\frac{\varepsilon}{2}
 \int_{\mathbb R^3}
 \mathbb ED_{(u,v)}D_{(u,w)}F(\Pi_1)\,du,dv\,dw+o(\varepsilon).
\]
Replacing the two insertions by a repeated insertion at one point is
incorrect: $v$ and $w$ remain distinct integration variables.

The sign mechanism is explicit.  For complete coverage of the horizontal
segment $[0,3r]\times\{0\}$, disks centered on a common vertical line
produce nested intervals on the segment.  The mixed difference is
nonpositive for every background and strictly negative on a positive
measure set of backgrounds and insertions.  For the vertical segment of
length $3r$, an empty background can instead require both insertions.
The notes bound its collapsed integral below by
\[
 r^3\left[\frac1{100}e^{-(6+\pi)r^2}
       -50\bigl(1-e^{-(6+\pi)r^2}\bigr)\right]>0
\]
when $r^2<\log(1+1/5000)/(6+\pi)$.  By decreasing thin-rectangle
limits the signs persist for nondegenerate rectangles.  These are actual
Ginibre examples.  They disprove the orientation-independent crossing
claim, while leaving the scalar critical-radius comparison open.  The
remainder in \eqref{br:anisotropic-expansion} is fixed-window; its
constants do not justify an unrestricted near-critical window limit.

\subsection{Structure factors, clusters, and exposed contacts}
\label{br:local-observables}

The structure-factor route used
\[
 S_G(k)=1-e^{-\pi|k|^2},\qquad S_{\Pi_1}(k)=1,
 \qquad
 \operatorname{Var}G(A)
 =\iint_{A\times A^c}e^{-\pi|x-y|^2}\\,dx\,dy.
\]
The supplied history records the perimeter bound
$\operatorname{Var}G(A)\leq\operatorname{Per}(A)/(2\pi)$ and stronger
spectral hole estimates involving both $|A|$ and this variance.  These
improve prescribed-hole costs but do not solve the unknown-curve union.
Second-order information alone does not determine connectivity.

The cluster route tried to compare the Palm susceptibility, namely the
expected size of the component containing a restored root.  The proposed
global inequality in the favorable direction is already false at small
range.  With $V=\pi R^2\leq1$, the recorded path estimate is
\[
 \chi_G(R)\leq
 \frac{1}{1-(V-1+e^{-V})}<1+V\leq\chi_{\Pi_1}(R).
\]
The strict middle inequality follows from $e^{-V}<1/(1+V)$ for $V>0$.
This is consistent with Ginibre's repulsion suppressing immediate Palm
neighbors.  It does not rule out a different comparison near criticality.
The connectedness Ornstein--Zernike identity, when its integrability
assumptions hold, contains the full pair connectedness density, not just
the ordinary structure factor; replacing one by the other was invalid.

The contact route compared the rate of exposed contacts as the radius
increases.  The history records a favorable exact local ratio over a
specified radius interval, but an exposed contact may either merge two
components or close a cycle within one component.  Euler characteristic
counts components minus holes and therefore does not resolve that
distinction.  A macroscopic comparison needs the contact weighted by
the exterior connections that make it pivotal.

For finite Ginibre, angularly averaging the conditional law of a pair
$m\pm se^{i\theta}$ gives a half-separation density proportional to
\[
 s^3e^{-2\pi s^2}\sum_k|c_k|^2s^{4k}.
\]
On a sufficiently short truncated interval this is more spread out than
the corresponding Poisson density proportional to $s$.  Its angular law
also changes, and spreading a pair can break its exterior attachments.
No connection-preserving transport to the Ginibre law was obtained.

\subsection{Coulomb dynamics and finite spherical systems}
\label{br:dynamics}

These calculations occur in Sections 7--8 of the user's research
history.  They are retained here as recovered finite-system work; this
report does not certify a Poisson-start infinite Ginibre diffusion.

An unrenormalized Poisson force cannot be used for such a construction.
The annular increment
\[
 Z_R=\sum_{z\in\Pi_1\,:,R<|z|<2R}\frac{z}{|z|^2}
\]
has mean zero and covariance $\pi\log2$ times the identity.  Its summands
are bounded by $1/R$, and the Poisson characteristic-function expansion
therefore yields a nondegenerate Gaussian limit.  In particular the
radial cutoff force is not Cauchy in probability.  Force differences at
two fixed roots have square-integrable far tails, so this rules out the
unrenormalized construction, not all relative-coordinate constructions.
Equilibrium stochastic differential equations do not by themselves
construct the intended process from Poisson initial data.

Two proposed instantaneous conclusions were also too strong.  If a
coupling gives a Poisson displacement process $Y_t$ and
$\sup_i|X_i(t)-Y_i(t)|\leq a(t)$, then
\[
 \mathcal O_{r-a(t)}(Y_t)\subseteq\mathcal O_r(X_t)
 \subseteq\mathcal O_{r+a(t)}(Y_t),
 \qquad |r_c(X_t)-r_c(\Pi_1)|\leq a(t).
\]
If $a(t)\to0$, there is no instantaneous jump of fixed size.  Likewise,
local convergence to Poisson plus sufficiently uniform spatial
decorrelation can imply $\liminf_{t\downarrow0}r_c(X_t)\geq r_c(\Pi_1)$.
Neither hypothesis was proved for the proposed unscreened dynamics.

On a sphere of area $A$, let $\mu_n$ be $n$ independent uniform points
and let $p_n(r)=\mathbb E_{\mu_n}F_r$.  For the repulsive diffusion with
generator
\[
 L_n^+=\sum_i\Delta_i+
 2\sum_i\sum_{j\ne i}
 \nabla_i\log|x_i-x_j|\cdot\nabla_i,
\]
the archived initial-response calculation gives
\begin{equation}
 \left.\frac{d}{dt}\mathbb E_{\mu_n}F_r(X_t)\right|_{t=0+}
 =\frac{4\pi n(n-1)}{A}\,[p_n(r)-p_{n-1}(r)]\geq0.
 \label{br:sphere-initial}
\end{equation}
Distances in the logarithm are chord distances; gradients and Laplacians
are spherical.  The invariant density is proportional to
$\prod_{i<j}|x_i-x_j|^2$.  The local equilibrium limit at $A=n\to\infty$
is the intended Ginibre law.  The limit of the coefficients in
\eqref{br:sphere-initial}, for a fixed physical observation window, is
$4\pi\partial_\lambda\mathbb EF_r(\Pi_\lambda)|_{\lambda=1}$.
This is a limit of initial derivatives, not a derivative of a constructed
infinite-volume evolution.  No common positive-time interval uniform in
particle number and crossing scale was obtained.

The attractive process killed at collisions supplies an exact finite
duality in the archive.  With $c_n=4\pi n(n-1)/A$, the adjoint repulsive
semigroup is $e^{c_nt}$ times the attractive killed semigroup.  Thus
\[
 \mathbb E^+_{\mu_n}F(X_t)
 =\mathbb E^-_{\mu_n}[F(Y_0)\mid\tau>t].
\]
A useful proof would compare crossing with attractive survival.  The
identity alone gives no such correlation sign.

Merging particles at collisions yields a finite coalescent whose
positions, conditional on its count, remain independent uniform at
deterministic times.  Its count decreases from $n$ at rate $c_n$, and the
recorded density limit is $(1+4\pi t)^{-1}$.  Reversing it creates births
at existing points.  A renewal identity yields the positive linear term
in \eqref{br:sphere-initial} and, under specified contact-surface
regularity hypotheses, a negative correction of order $t^{3/2}$.
Those remainders were fixed-system.  Removing a newborn changes later
forces on all particles, so a birth is not an independent local
enhancement.  Movement, feedback, and rare-pivotal normalization were
not controlled.

Finally, the conjectured explicit identification of this evolution with
the $\beta$-Ginibre family fails.  The pair equation agrees with
$\beta'=8\pi(1-\beta)$, but at three collinear points of adjacent squared
distance $\beta\log2/\pi$ the recorded three-point discrepancy is
\[
 \frac{\pi(1-\beta)}{2\beta}(5-6\log2)>0,
 \qquad0<\beta<1.
\]
Pair-correlation agreement therefore cannot identify the law.  This
countercalculation must be retained if the dynamics route is resumed.

\subsection{Static interaction-strength interpolation}
\label{br:energy}

On a sphere of area $N$ and radius $R_N$, the separate static family in
the history is
\[
 d\mu_{N,\theta}=Z_{N,\theta}^{-1}e^{-\theta H_N}d\mu_N,
 \quad
 H_N=-\sum_{i<j}\log\frac{|x_i-x_j|}{2R_N},
 \quad0\leq\theta\leq2.
\]
For a crossing-failure event $B$ of positive probability, differentiation
of the normalized finite integral gives exactly
\begin{equation}
 \log\frac{\mu_{N,2}(B)}{\mu_{N,0}(B)}
 =-\int_0^2
 \bigl(\mathbb E_{N,\theta}[H_N\mid B]
                  -\mathbb E_{N,\theta}H_N\bigr)\,d\theta.
 \label{br:energy-identity}
\end{equation}
Thus a uniform positive integrated energy excess for macroscopic failure
would be useful.  That is the unproved principal quantity, not a
consequence of logarithmic repulsion alone.

The recorded Poisson-endpoint calculation is singular in ambient volume:
for a fixed local observable $F$,
\[
 \left.\partial_\theta\mathbb E_{N,\theta}F\right|_{\theta=0}
 =-\frac{\log N}{4}
   \left.\partial_\lambda^2\mathbb EF(\Pi_\lambda)
   \right|_{\lambda=1}+O_F(1).
\]
Some small rectangular crossings have positive intensity curvature,
which rules out an unrestricted positive crossing derivative in this
interpolation.  Conditioning on exactly $n$ points in a fixed planar
region $W$ removes that divergent count contribution.  The limiting
derivative then involves
\[
 -\operatorname{Cov}_{\mu_{n,W}}(F,J_W),\qquad
 J_W=-\sum_{i<j}\log|x_i-x_j|
       +\sum_i\int_W\log|x_i-y|\\,dy.
\]
A favorable conditional example in the history does not control the
unconditioned count distribution or the entire integral in
\eqref{br:energy-identity}.

\subsection{Common-realization square replacement and its genuine endpoint gain}
\label{br:squares}

Partition the plane into half-open squares of side $\ell$.  Independently
mark each square with probability $q$; use the restriction of one common
Ginibre realization in marked squares and an independent Poisson
realization in the others.  Write $\Sigma_{q,\ell}$ for the result.
Its endpoints are exactly Poisson and Ginibre, but the intermediate law
is generally not determinantal.  For a local event,
\[
 \partial_q\mathbb EF(\Sigma_{q,\ell})
 =\sum_Q\mathbb E\bigl[F(\Sigma^{Q,1}_{q,\ell})
                       -F(\Sigma^{Q,0}_{q,\ell})\bigr].
\]

At $q=0$, the exterior of the replaced square is independent Poisson.
When the square is a clique, write $A,B\subset Q$ for locations
attaching an inserted disk to the two relevant exterior components.  Put
$h(S)=e^{-|S|}-\mathbb P(G(S)=0)$.  The exact replacement gain,
conditional on this exterior, is
\begin{equation}
 h(A)+h(B)-h(A\cup B).
 \label{br:square-gain}
\end{equation}
Even though $h\geq0$, \eqref{br:square-gain} need not be nonnegative.
The positive endpoint argument instead compares robust favorable
geometries, such as $A=B=Q$, with adverse thin-shell geometries.

The supplied history records, using its Poisson arm inputs and local
modifications, a strictly positive derivative
\[
 \partial_q C_{L,\ell}(0,r_c(\Pi_1))
 \geq c\ell^2 I_L,
 \qquad
 I_L=\int\mathbb ED_xF_{L,r_c(\Pi_1)}(\Pi_1)\,dx.
\]
This is a substantial historical endpoint result.  The detailed earlier
arm proofs are not reproduced and freshly certified here.  It was not
lost when the later report excluded an unjustified positive-density
continuation.

The exact obstruction to a pointwise continuation is the dependence of
the common Ginibre realization across squares.  The relevant void law
becomes conditional on observed Ginibre exterior points.  The history
records essential conditional hole-probability extrema zero and one for
every $q>0$, using rigidity and finite fully marked annuli.  Hence uniform
improvement for every exterior cannot replace an averaged argument.
Moreover, when the Goldman extra point lands in a square, its mark forces
the entire square's color; it is not an independent mark on one point.

For fixed $L$ the recorded mesh expansion, with $a=\ell^2$, is
\[
 \partial_q C_{L,\ell}(0,r_c(\Pi_1))
       =\frac a2I_L+o_L(a),
 \qquad
 \partial_q^2 C_{L,\ell}(0,r_c(\Pi_1))\longrightarrow-J_L,
\]
where
\[
 J_L=\iint e^{-\pi|x-y|^2}
                \mathbb ED_xD_yF_{L,r_c(\Pi_1)}(\Pi_1)\,dx\,dy.
\]
The history's controlled positive-density regime is
$L\leq c\ell^{-1/3}$ and $q\leq c\ell^2$.  It is not a macroscopic
critical-radius comparison.  Indeed, in the natural coupling,
\[
 d_{\rm TV}\bigl(\Sigma_{q,\ell}|_W,\Pi_1|_W\bigr)
 \leq2q|W|.
\]
For its growing windows this is $O(L^{-4})$, so the total crossing gain
tends to zero.  Also a lower derivative bound by a fixed positive
multiple of $I_L\to\infty$ cannot hold throughout a fixed interval in
$q$: integrating would make a probability increase by more than one.
The necessary repair is a comparison normalized by influences of the
\emph{current} law.

The subsequent fine-mesh calculation identifies the local limit
\[
 \Sigma_{q,\ell}\longrightarrow
 \operatorname{thin}_q(G)\cup\Pi_{1-q}
\]
in total variation, uniformly in $q$ on a fixed window.  Couple square
marks to point marks unless two source points share a square; the union
of the two source processes has second factorial density at most four,
giving an error at most $2\ell^2|W|$.  In this limit the first derivative
at zero vanishes and the second derivative is $-J_L$.  A real order-$a$
endpoint effect therefore disappears with the mesh; it cannot be
silently retained as a positive limiting derivative.

\subsection{Fixed-kernel Poisson expansion and a positive long-arm family}
\label{br:fixed-kernel}

The recovered source
\path{research/fixed_kernel_poisson_second_variation.md}, with
\path{fixed_kernel_poisson_second_variation_audit.md}, goes beyond the
second derivative in the preceding subsection. For
$X_s=\operatorname{DPP}(sH_1)\cup\Pi_{1-s}$ and a bounded functional
$F$ supported in $W$, put $m_W=|W|$ and let $B_n$ have zero diagonal
and off-diagonal entries $H_1(x_i,x_j)$. The exact centered factorial
density is $s^n\det B_n$, so
\begin{equation}\label{br:entire-expansion}
 \mathbb EF(X_s)=\sum_{n\ge0}\frac{s^n}{n!}
  \int_{W^n}\det B_n\,
       \mathbb E_\Pi[D_{x_1}\cdots D_{x_n}F]\,dA^n.
\end{equation}
Hadamard's bound $|\det B_n|\le(n-1)^{n/2}$ and
$|D^nF|\le2^n\|F\|_\infty$ make this an entire series in $s$.
Writing $C(s)=\mathbb EF(X_s)$, a cycle majorant gives, for
$|s|\le1/8$,
\[
 |C(s)-C(0)-\tfrac12s^2C''(0)|
 \le\|F\|_\infty
 \left[\tfrac{64}{3}m_W|s|^3+
       \tfrac{25}{2}m_W^2s^4e^{5m_Ws^2}\right].
\]
This is explicitly a fixed-window error.

For a deterministic background with two terminal-attached parents
at distance $d$, where $R<d<2R$ and $R=2r$, let $V$ be their
radius-$R$ lens and $A,B$ the two exclusive crescents. With
$q(h)=e^{-\pi|h|^2}$, the joint pair contribution is
\[
 \mathrm{Red}=\int_{V^2}q(y-x)\,dx\,dy,\qquad
 \mathrm{Coop}=2\int_{A\times B}q(y-x)
                        \mathbf1_{\{|x-y|\le R\}}\,dx\,dy.
\]
Gaussian translation and bounded variation imply
$\mathrm{Red}\ge|V|-\operatorname{Per}(V)/(2\pi)$.
The two exclusive sets lie on opposite sides of the perpendicular
bisector; integrating the Gaussian over these halfplanes and their
vertical span gives $\mathrm{Coop}\le2R/\pi$.
Consequently
\[
 \liminf_{d\downarrow R}(\mathrm{Red}-\mathrm{Coop})
 \ge R\left[(2\pi/3-\sqrt3/2)R-\tfrac23-\tfrac2\pi\right].
\]
It is strictly positive if
\begin{equation}\label{br:lens-positive-range}
 R>\frac{2/3+2/\pi}{2\pi/3-\sqrt3/2},
\end{equation}
including the source's rationally checked value $\pi R^2=9/2$.
The limit is one-sided: at $d=R$ the old parent edge opens.
For comparison, at $d=5R/2$ there is no singleton bridge but an
open set of cooperative pairs, hence a strictly negative contribution.
This proves neither a universal conditional sign nor a Poisson average.

The positive estimate extends to arbitrarily long arms and their
neutral Poisson background. Put the tips at $0$ and $d$ on the
horizontal axis. Keep all other left-arm centers at first coordinate
at most $-\ell_0$, and all other right-arm centers at least
$d+\ell_0$, with $2R-d<\ell_0<R$; all centers lie in a strip
$|y|\le b_0$. Let $U_S,U_T$ be the radius-$R$ neighborhoods of
the two connected arms. Their intersection is exactly the tip lens.
For every neutral configuration outside $U_S\cup U_T$,
\begin{equation}\label{br:neutral-long-arms}
 \mathrm{Red}-\mathrm{Coop}
 \ge |V|-\frac{\operatorname{Per}(V)}{2\pi}
                         -\frac{2(R+b_0)}{\pi}.
\end{equation}
The bound permits cooperative paths through any number of neutral
vertices. It is positive for $d$ sufficiently close to $R$ and
$b_0$ small under \eqref{br:lens-positive-range}.
Conditional on the arms and no additional points in $U_S\cup U_T$,
the entire remaining Poisson process can therefore be integrated.
The actual void factor $e^{-|U_S\cup U_T|}$ is retained.

For the true seed-to-outer event, take collinear arms from the
deterministic terminal attachment sets to tips $u_0,v_0=u_0+d$.
With positive margins
$D_{\rm seed}=v_0-R-(m+r)$ and
$D_{\rm outer}=L-r-(u_0+R)$, restrict the other centers to the free
part of the active window, outside both attachment sets and arm
neighborhoods. The bound becomes
\[
 \mathrm{Red}-\mathrm{Coop}\ge
 |V|-\frac{\operatorname{Per}(V)}{2\pi}-\frac{2R}{\pi}
 -2\pi(m+r)^2e^{-\pi D_{\rm seed}^2}
 -2|U_S|e^{-\pi D_{\rm outer}^2}.
\]
At $u_0=(L+m)/2$ with fixed $m,d$, the errors vanish as $L\to\infty$,
since $|U_S|=O(L)$. Strict edge and terminal slack allow open
perturbations of the arm positions. This is a favorable
positive-probability family for the actual event, with all free neutral
points integrated. Its probability may be very small, and no uniform
fraction of pivotal weight or verification of
\eqref{br:lens-positive-range} at $r_P$ is asserted.

\subsection{Bottleneck thresholds, cloning, and contact surgery}
\label{br:cloning}

For a finite terminal graph the first crossing radius is a minimax:
\[
 \mathcal R(\eta)=\min_{\gamma}\max_{e\in\gamma}w_e,
 \qquad w_{xy}=|x-y|/2.
\]
Terminal edge weights are distances to the relevant terminal.  With
$T=\max(a,\min(b,\mathcal R))$, the deterministic layer-cake identity is
\[
 T=b-\int_a^bF_r,dr.
\]
This avoids differentiating a nonsmooth threshold at every radius.

A stronger deterministic comparison of a two-point gain with integrated
one-point gains is false.  In $[0,6]\times[0,1]$, take centers
\[
 (1/2,1/2),(1,1/2),(3,1/2),(5,1/2),(11/2,1/2).
\]
The left--right threshold is one and remains one after any single
insertion.  Insert instead $(3-h,1/2)$ and $(3+h,1/2)$, where
$0<h<2/3$.  The threshold becomes $1-h/2$.  The insertions can be
arbitrarily close.  Perturbing the central old point removes the tied
bottleneck while making the integral of one-point gains arbitrarily
small and retaining a positive two-point gain.  Probability weighting,
rather than a deterministic inequality, is essential.

For the whole-plane seed-to-outer event, adding a clone at displacement
$h$ decreases $T$ by at most $|h|$.  The parent integration must include
the whole plane: a parent outside a nominal center window may have a
clone entering the active region.  Dropping this term was one of the
corrected cofactor-boundary errors.

For Poisson intensity $\lambda$, define the Gaussian clone gain $I_\varepsilon$
and the mean absolute move gain $M_\varepsilon$ using displacement
$\varepsilon Z$, with density $e^{-\pi|z|^2}$ for $Z$.  The exact
decomposition is
\begin{equation}
 I_\varepsilon=\frac{M_\varepsilon}{2\lambda}+J_\varepsilon,
 \qquad J_\varepsilon\geq0,
 \label{br:clone-decomp}
\end{equation}
where $J_\varepsilon$ integrates
\[
 \min\{T(\eta+x),T(\eta+y)\}-T(\eta+x+y).
\]
It counts connections requiring both insertions.  The movement part
satisfies
\[
 M_\varepsilon\leq\frac{\varepsilon}{\pi}
          \mathbb P(a<\mathcal R<b).
\]
The coefficient follows because a generic unique center--center
bottleneck has two endpoint gradients of norm $1/2$, a terminal
bottleneck has one gradient of norm one, and
$\int e^{-\pi|z|^2}|z\cdot v|\,dz=|v|/\pi$.
Mecke's formula, followed by translation of the moved point, makes this
bound uniform in the window.  In contrast, the dominated convergence
giving
$I_\varepsilon/\varepsilon\to
\mathbb P(a<\mathcal R<b)/(2\pi\lambda)$ is fixed-window.

The cooperative-pair notes develop a more ambitious bound of order
$|h|^2$ times the ordinary pivotal integral.  Their construction retains
two occupied arms, excludes bypasses, and treats circular terminal
strips separately.  The positive-bath version fills a resampling region
with Poisson points and its constants deteriorate as the bath vanishes.
The zero-bath version uses sparse Fock projection estimates and
conditional Janossy densities.  These lengthy geometric chains are
\emph{recovered calculations}; their exact relationship to the claimed
small-$\beta$ theorem and their outstanding verification are set out in
the separate small-$\beta$ section.

Contact surgery isolated a further precise issue.  At a regular
artificial-grain contact $y$--$w$, $|y-w|=2r$, the open lens
$B(y,2r)\cap B(w,2r)$ contains no background point.  A predecessor $z$
may nevertheless approach $y$ from the rear.  Its normalized coherent
contribution is $e^{-\pi|z-y|^2/\beta}$ and tends to one.  The
canonical contact partition gives the archived reduction
\[
 N\leq sB+(1-s)\int p_A(y),d\mu_{\rm contact}(A,y),
\]
where $B$ is total contact mass and $p_A$ is the forced-certificate
projection diagonal.  Thus the random remainder is controlled, while
the contact-weighted predecessor term still needs geometry.

Two shortcuts were explicitly ruled out.  First, two terminal clusters
with gap $4r+d$ have zero one-grain contact mass at radius $r$, but two
nearby insertions can bridge them.  Second, at fixed $\beta$ and zero
bath a local connection requiring $n$ points has leading contact mass
$b_{n-1}s^{n-1}$, and its leverage-weighted mass has a positive leading
coefficient of the same order.  Their ratio therefore does not vanish
like $s$ as $s\downarrow0$.  This does not refute a bound restricted to
$s\geq\beta/2$, or a uniform constant below one at small $\beta$.

\subsection{Recovered half-Gaussian and twin-contact identities}
\label{br:contact-refinements}

Sections 5--6 of \path{research/current_law_contact_obstructions.md}
calculate the portion of contact surgery that a half-space estimate
actually controls. At a regular contact $|y-w|=2r$, put
$n=(w-y)/(2r)$ and $z=y+h$. Absence of a bypass for $|h|<2r$
requires $h\cdot n\le |h|^2/(4r)$. For
$q_\beta(h)=\beta^{-1}e^{-\pi|h|^2/\beta}$ the full Gaussian mass
of the allowed side obeys
\[
 p_\beta(r)\le\frac12+\frac{\sqrt\beta}{6\pi r}
                         +e^{-\pi r^2/\beta}.
\]
Indeed, on $|h|\le r$, its positive normal coordinate is at most
the square of the tangential coordinate divided by $3r$; the latter
coordinate has variance $\beta/(2\pi)$.

This factor bounds only contacts which remain pivotal when $z$ is
deleted. In the contact measure
$\delta(|y-w|-2r)C(\eta+z;y,w)\,dy$, the complementary nonnegative
part has $C(\eta;y,w)=0$: the new point creates the contact's
connectivity requirement. Denote its integrated mass by $Q_\beta$.
The original contact mass can vanish while $Q_\beta>0$, as in the
gap example above. Swapping $y,z$ moves its matching contact to a
different artificial radius and does not cancel it at radius $r$.
Equivalently, with the $z$-grain fixed at radius $r$ and $y$-grain
radius $u$, the difference
$J(u)=F(\eta+B(y,u)+B(z,r))-F(\eta+B(y,u))$
is an interval indicator when the background and $z$ alone fail.
The contact difference is the signed measure $dJ/du$.

There is also an exact symmetric perimeter identity. For fixed
$0<|h|<2r$, define the successful dimer midpoint set
\[
 U_h(\eta)=\{c:F_r(\eta+\{c-h/2,c+h/2\})=1\}.
\]
At a regular boundary point exactly one endpoint supplies the
critical contact. Coarea, with both contact and midpoint translation
speed one, gives
\[
 \int dc\,[A_{c-h/2}^r(\eta+c+h/2)
               +A_{c+h/2}^r(\eta+c-h/2)]
       =\operatorname{Per}(U_h(\eta)).
\]
Here $A_y^r$ denotes the contact measure for increasing only the
artificial grain at $y$; identities may be read against radius test
functions. Gaussian symmetry then yields
\[
 \int_{|h|<2r}q_\beta(h)\int A_y^r(\eta+y+h)\,dy\,dh
 =\frac12\int_{|h|<2r}q_\beta(h)
                           \operatorname{Per}(U_h(\eta))\,dh.
\]
The comparison with $\operatorname{Per}(U_0)$ remains unproved;
the separate cofactor count term has additional weights. Thus neither
identity gives a half-factor for the full leverage integral.

\subsection{Deletion certificates and conditional count tails}
\label{br:counts}

The certificate method produced estimates normalized by the actual
event probability.  Suppose a point process $\eta$ obeys reduced
Campbell domination
\[
 \mathbb E\sum_{z\in\eta}g(z,\eta-z)
 \leq M\int\mathbb Eg(z,\eta)\,dz
\]
for nonnegative $g$.  If an event $E$ has a witness using at most $k$
protected points in a cell $Q$, deletion of every other point preserves
$E$, whence
\[
 (N(Q)-k)\mathbf1_E\leq
 \sum_{z\in\eta\cap Q}\mathbf1_E(\eta-z).
\]
Applying Campbell gives
\begin{equation}
 \mathbb E[N(Q)\mathbf1_E]\leq(M|Q|+k)\mathbb P(E).
 \label{br:count-mean}
\end{equation}
A shortest path uses at most two vertices of a clique cell; if the
artificial pivotal root lies in that clique, at most one removable
witness vertex is needed there.

The exponential version is equally explicit.  Multiplication by
$e^{tN(Q)}$ and the same Campbell argument yield
\[
 F'(t)\leq(k+M|Q|e^t)F(t),
 \quad F(t)=\mathbb E[e^{tN(Q)}\mathbf1_E],
\]
and hence
\begin{equation}
 \mathbb E[e^{tN(Q)}\mid E]
 \leq\exp\{kt+M|Q|(e^t-1)\},\qquad t\geq0.
 \label{br:count-mgf}
\end{equation}
Factorial-moment domination supplies the needed exponential
integrability, or one may first use finite generating-series
truncations.  For
$\operatorname{DPP}(sH_\beta)\cup\Pi_\nu$, $s<1$, the applicable
constant is $M=\nu+s/[\beta(1-s)]$.
These are actual event-weighted estimates, not independence claims.

They control counts but not a nearly singular coherent Gram inverse.
At $s\uparrow1$ the displayed Campbell constant also diverges.  An
abstract count-only attempt to absorb this divergence fails for a
rank-one projection supported in a region $Q$: the event
$\{N(Q)\geq2\}$ has a two-point deletion certificate but probability
zero, whereas one artificial insertion can make it occur whenever the
single retained point is present.  The residual leverage is then one.
This is not a Ginibre pivotal counterexample; it disproves the proposed
abstract count-only implication.

The canonical-certificate ordering is important.  Choose candidates by
an intrinsic ordering depending only on the candidate's own vertices
and deterministic geometry.  After forcing a chosen tuple $A$, absence
of an earlier valid witness and absence of a bypass are decreasing
conditions on the remainder.  An ordering whose validity changes when
unselected points are deleted would not justify the same partition.

For a normalized root row $f$ of a residual projection, let
$t=\int_{B^c}|f|^2$.  The empty-region rank-one resolvent calculation
gives, under an additional decreasing exterior event, the useful
remainder bound
\[
 \frac{st}{1-s+st}.
\]
The sparse-Fock notes then aim to make $t$ and the certificate
projection small at small correlation scale.  They use divided
differences for near-colliding points before block Gram estimates.
They do not assert well-conditioning of the original coherent Gram
matrix or a deterministic localization theorem at $\beta=1$.

\subsection{A protected close-pair bound uniform at the projection endpoint}
\label{br:close-pairs}

The recovered estimate in \path{pivotal_close_pair_estimate.md}
concerns an internal consecutive DPP pair $z,z+h$ on an induced
pivotal path, with $|h|\le\delta<r$. Both outside neighbors must be
old random centers, rather than the artificial root or a terminal.
Chordlessness places them in the two disjoint crescents
$B(z,2r)\setminus B(z+h,2r)$ and
$B(z+h,2r)\setminus B(z,2r)$, each of area at most $4r|h|$.

For $z$ in a fixed bounded set $Q$, let $K_{Q,\delta}$ count all
eligible ordered quadruples. Choose nested fixed disks containing
the quadruples, and let $E_{\rm ext}$ be measurable outside the
larger disk. Its weighted factorial bound is preserved under deletion
of the four points. The buffered exterior Bergman estimate supplies
constants $b_0,b_1$ on the smaller disk with diagonal at most $b_0$
and pair determinant at most $b_0b_1^2|h|^2$.
For $\operatorname{DPP}(sH_1)\cup\Pi_\nu$, summing the neighbor colors
and integrating the pair separation gives
\begin{equation}\label{br:delta-six}
 \mathbb E[\mathbf1_{E_{\rm ext}}K_{Q,\delta}]
 \le \frac{16\pi}{3}s^2b_0b_1^2(sb_0+\nu)^2
           r^2|Q|\delta^6\,\mathbb P(E_{\rm ext}).
\end{equation}
The sixth power consists of two separation dimensions, two crescent
widths and the quadratic determinantal zero. The earlier Campbell
bound had constant
$(16\pi^2/3)[s/(1-s)]^2[s/(1-s)+\nu]^2$ and diverged at $s=1$.
Equation~\eqref{br:delta-six} does not.

The note records submission for independent audit; the estimate is
restored here as a recovered local calculation. A separate
arm-preserving comparison
$\mathbb P(E_{\rm ext})\le C\mathbb P(\mathrm{Piv}_x)$ is still
needed. Root and terminal neighbors require additional strip
arguments, and a shrinking close-pair cutoff cannot exclude the
fixed-separation saturation family below.

\subsection{Fock saturation, negative thin lenses, and positive ray examples}
\label{br:saturation}

The exact bridge multiplier and its variational characterization are
proved elsewhere in this report.  Their significance here is that they
separate two effects: conditioning on occupied arms removes coherent
directions, while the necessarily empty bridge region increases the
resolvent.

For two parents $A=\{-a,a\}$ with $R/2<a<R$, let
$B_a=B(-a,R)\cap B(a,R)$.  Direct inversion of the two-by-two Gram
matrix gives
\[
 P_A(0,0)=\frac{2e^{-\pi a^2}}{1+e^{-2\pi a^2}}
          =\operatorname{sech}(\pi a^2)>0.
\]
As $a\uparrow R$, the lens area tends to zero.  For the
Lebesgue-integrated trace surplus
$\int_{B_a}(M_{A,B_a,s}(x)-1)\\,dx$, the loss of coherent directions
is bounded below by a positive constant times $|B_a|$, whereas the
positive resolvent trace correction is $O(|B_a|^2)$.  These are
integrated estimates: the pointwise deficit at the origin does not
vanish with the lens area.  This proves the failure of a
certificate-by-certificate nonnegative integrated surplus, including
at $s=1$ once $|B_a|<1$.

The later saturation construction is stronger and is preserved as a
recovered geometric calculation.  Two induced paths traverse paired
grid points with macronode spacing
$D=R+\epsilon/2$ and fixed small within-pair displacement.  If $D^2<2$,
their eventual doubled-lattice subset has zero density $2/D^2>1$.
Jensen's formula and the Fock evaluation bound
$|f(z)|\leq\|f\|e^{\pi|z|^2/2}$ imply completeness of those coherent
directions.  The forced root deficits $q_N=1-P_{A_N}(0,0)$ therefore
tend to zero.  Strict geometric inequalities retain two disconnected
arms and an actual pivotal insertion, including circular terminals.

Their \emph{entire} bridge region is one fixed thin lens $B$, independent
of path length, with $|B|<1$.  The scalar estimate is elementary:
\[
 \|(Q_{A_N})_B\|\leq|B|,
 \qquad
 \int_B|Q_{A_N}(z,0)|^2\\,dz\leq q_N|B|,
\]
so its multiplier satisfies
\begin{equation}
 0\leq M_{A_N,B,s}(0)
 \leq\frac{q_N}{1-s|B|}
 \leq\frac{q_N}{1-|B|}\longrightarrow0
 \label{br:saturated-multiplier}
\end{equation}
uniformly for $0\leq s\leq1$.  The full geometric verification is in
the source notes.  This construction shows why deterministic certificate
positivity cannot be the missing estimate; it does not assess the
probability of those certificates under actual pivotal conditioning.

The archive also contains a concrete \emph{positive} family.  Put
\[
 \pi a^2=11/5,\qquad\pi t^2=2/5,
 \qquad R=a+t,\quad\ell=99R/100,
 \quad A_N=\{\pm(a+n\ell):0\leq n\leq N\}.
\]
Let $B_t=B(0,t)$, a disk contained in the two-parent bridge lens.
The recovered rational bounds for a void in this disk give
$M_{A_N,B_t,s}(0)>1+3/1000$, uniformly in $N$ and
$24/25\leq s\leq1$.  The proof retains the two-parent Schur complement
before controlling all remaining ray vertices; it is not an estimate
obtained by ignoring the arms.  But no uniform positive fraction of
pivotal configurations was shown to have that geometry, and the
available analytic threshold bounds did not place its radius below the
Poisson critical radius.

A second local construction moved two tight triples and increased their
probability by a holomorphic circle average.  This did not make them
positive bridge targets.  For its particular six-point family with
$\pi R^2=9/2$, the recorded analytic calculation gives
\[
 P_A(0,0)>27/100,\qquad |B|<15/56,
 \qquad M_{A,B,s}(0)<1022/1025<1
\]
uniformly in the phase and $s\leq1$.  Probability amplification and
positive pivotal multiplier are separate requirements.

\subsection{Count-preserving motions and exterior holomorphic fields}
\label{br:motions}

The conditional Ginibre density in a disk, given its full exterior and
its rigid inside count, has a Vandermonde factor, a Gaussian factor, and
the squared modulus of a nonvanishing holomorphic exterior factor.
Consequently a fixed count-preserving motion must be compared under
that exact density, not an independent Ginibre density with the
exterior discarded.

For two parents moved from $\pm a$ to $\pm b$, the full Palm likelihood
contains
\[
 \prod_{w}\left|\frac{w^2-b^2}{w^2-a^2}\right|^2
\]
with the prescribed radial product and normalizing constant.  The
motion notes establish the exact inverse-moment range $0<q<1$ for
the unrestricted likelihood.  The obstruction at $q\geq1$ is already
local: a target Palm zero at $b$ has density vanishing quadratically,
whereas the source Palm density there is positive.  Integrating the
inverse likelihood therefore includes
$\int_{|w-b|<\delta}|w-b|^{-2q}\\,dw$, which diverges exactly at
$q\geq1$.  The same local obstruction persists under every nonzero
thinning without an uncolored Poisson bath.  Large-order exterior
moment estimates under a genuine latent void cannot be substituted
for these unrestricted moments.

\paragraph{Fixed-count support and averaged motion.}
The source \path{bridge_motion_continuation.md} gives a separate
geometric obstruction. Put two exterior ports at
$\pm(a+R-\delta)e_1$ and retain exactly two movable inside parents.
Any replacement preserving their separate terminal attachments has
parents $u,v$ satisfying $|u-v|\ge2a-2\delta$.
A proposed widening requiring $|u-v|\le2b$ is therefore impossible
when $b<a-\delta$, even before probabilities are compared.
The ports and paths admit open perturbations.

Under a genuine latent void of radius $D>a>b>0$, the motion notes
give an averaged quantitative repair. Put
\[
 \Delta=a^2-b^2,\quad B_0=\frac{2\Delta}{D^2-a^2},\quad
 V_0=\frac{4\pi\Delta^2}{D^2-a^2},\quad
 \kappa=e^{2\pi\Delta}\frac{b^2}{a^2}
                    \frac{1-e^{-4\pi a^2}}{1-e^{-4\pi b^2}}.
\]
For the same residual event $E$ under the reduced Palm laws
$\mu_a,\mu_b$, with $E\subset\{X(B_D)=0\}$ and $p=\mu_a(E)>0$,
\begin{equation}\label{br:motion-relative}
 \frac{\mu_b(E)}{\mu_a(E)}
 \ge\kappa e^{-2B_0}p^{1/q}
             \exp[-\tfrac12q e^{qB_0}V_0]\qquad(q>0).
\end{equation}
The bounded far log-product has variance proxy $V_0$ and centered
log-MGF at most $\tfrac12t^2e^{|t|B_0}V_0$; Holder gives this bound.
Optimizing it, as checked in
\path{ginibre_motion_independent_audit.md}, gives a relative factor
$\exp[-(B_0+o(1))\log(1/p)/\log\log(1/p)]$ as $p\downarrow0$.
It still vanishes for rare events. This optimization concerns the
void-restricted product and does not extend the unrestricted
inverse-moment range $q<1$.

More generally, for $L=d\mu_b/d\mu_a$ and independent thinning $T_s$,
the exact pushed-forward likelihood is
\[
 \frac{d(T_s\mu_b)}{d(T_s\mu_a)}(\eta)
       =\mathbb E_{\mu_a}[L(X)\mid T_sX=\eta].
\]
Adding a bath and forgetting colors is another Markov kernel with
the same conditional-expectation rule. Finite full-law inverse moments
transfer by conditional Jensen, but a thinned observed void does not
imply the latent void needed in \eqref{br:motion-relative}.
Changing the parent-dependent definition of $E$ also requires a
geometric inclusion or a justified configuration map.

A holomorphic circle does not uniformly close a lone two-parent gap.
If $z_1(t),z_2(t)$ are holomorphic in the unit disk and continuous on
its boundary, the maximum modulus principle prevents
$|z_1(t)-z_2(t)|$ from being strictly smaller at every boundary point
than at $t=0$.  More generally a fixed witness path cannot have all
its edge lengths uniformly improved on the whole parameter circle
unless the central path was already improved.

Redundancy allows the active witness to change with phase.  The archive
constructs four points within $\eta$ of $c$, moving with velocities
$\rho(1,i,-1,-i)$.  If, for example, $\rho=r/4$ and
$\eta\leq r/100$, their radius-$r$ occupied union contains the old
occupied union at every boundary phase.  Jensen's formula for the six
Vandermonde factors gives the recovered density amplification
\[
 e^{-4\pi\rho^2}
 \left(\frac{\rho}{\sqrt2\eta}\right)^{12}.
\]
The corresponding two-triple construction has factor
\[
 \frac1{20}e^{-6\pi\rho^2}
 \left(\frac{\sqrt3\rho}{2\eta}\right)^{12}.
\]
The factor $1/20$ is a labeling allocation, not a free multiplicity
gain.  These are comparisons for \emph{preexisting} redundant clusters.
No uniform supply of such clusters under rare two-arm conditioning was
proved.

Moving whole arms with tapered displacements preserves their slack
edges, but integrating over the moving root lens introduces an
additional obstruction.  Its area, for parent separation $d$, is
\[
 A(d)=2R^2\arccos(d/(2R))-\frac d2\sqrt{4R^2-d^2}.
\]
Although $\Delta A(d)>0$ for thin lenses, one has
$\Delta\log A(d)<0$ for every $0<d<2R$.  Indeed
$A'=-\sqrt{4R^2-d^2}$, $A''=d/\sqrt{4R^2-d^2}$, and
$A(d)<(2R-d)\sqrt{4R^2-d^2}$ imply $AA''-(A')^2<0$;
the radial term $A'/(dA)$ is also negative.

Under a holomorphic motion $Z+tv$, the exact density-weighted lens
mass has the form
\[
 W(t)=e^{-\kappa|t|^2}A(|2a-2\delta_0t|)|f_O(t)|^2,
 \quad\kappa=\pi\sum_j|v_j|^2,
\]
where $f_O$ is nonvanishing holomorphic and depends on the exterior.
Writing $u=\nabla\log A$ and $v_O=\nabla\log|f_O|^2$ at zero gives
\[
 \frac{\Delta W(0)}{W(0)}
 =4\delta_0^2\Delta\log A(2a)-4\kappa+|u+v_O|^2.
\]
An exterior field that cancels $u$ makes this strictly negative.
The later compact-factor support construction was needed to realize
such fields by actual infinite-Ginibre exteriors; a finite polynomial
ensemble example alone would not have supplied that conclusion.
The fixed-Palm exterior-ratio result proved elsewhere in this report
records the completed form of this obstruction.  It still concerns
arbitrary exteriors, not their weights in a pivotal average.

\subsection{Density deficits, random zero sets, and the rare-event loss}
\label{br:density}

A pivotal two-arm certificate can be chosen as two induced paths in
\emph{different} background components.  Both paths cannot meet the
same clique cell.  Thus, for every such certificate $A$ in the same
realization,
\begin{equation}
 |A\cap Q|\leq\min\{G(Q),2\}.
 \label{br:cap-two}
\end{equation}
This is cap two, not the cap four for an arbitrary union of two paths.
Let $a=|Q|$ and
$\Delta=a-\mathbb E\min\{G(Q),2\}
=\mathbb E(G(Q)-2)_+>0$.
Strict positivity follows by integrating the strictly positive
three-point determinant over three separated target disks in $Q$.
Negative association of disjoint cells and the bounded range $[0,2]$
give, for a union $U$ of $n$ cells and suitable $\delta>0$,
\[
 \mathbb P\left(\sum_{Q\subset U}\min\{G(Q),2\}
          \geq(1-\delta/2)|U|\right)
 \leq\exp\{-\delta^2a|U|/8\}.
\]
This bounds every certificate simultaneously, without a union bound
over possible paths.

\paragraph{A recovered conditional density gap.}
\label{br:conditional-density}
Sections 5--6 of \path{research/induced_certificate_density_gap.md}
strengthen the preceding unconditional estimate under a specific
exterior conditioning. For concentric radii $M<R$, $M\le R/10$,
and $\pi R^2\ge5$, the weighted Bergman series gives
\[
 0\le B_R(z,z)-1\le\varepsilon_R
       :=6e^{-c_0\pi R^2},\qquad |z|\le M,\quad
 c_0=(1-\log2)/2.
\]
To obtain the bound, split the monomial series at
$\lfloor\pi R^2/2\rfloor$. A Poisson Chernoff bound controls the
low modes and the first term of the Poisson upper tail controls
the high modes. Thus the constants are uniform on the smaller disk.

Let $U\subset B_M$ be a union of clique cells of area $a$, let
$S_U=\sum_{Q\subset U}\min\{G(Q),2\}$, and choose $R$ large enough
that $\varepsilon_R\le\delta/4$, where the ordinary capped mean is
at most $a(1-\delta)$. Conditional domination by the fixed Bergman
DPP, followed by its negative association, yields
\begin{equation}\label{br:conditional-cap-tail}
 \mathbb P\left(S_U\ge(1-\delta/2)|U|
                      \,\middle|\,G|_{B_R^c}\right)
       \le e^{-\delta^2a|U|/32}.
\end{equation}
Indeed monotone coupling with $G|_{B_M}$ increases each capped mean
by at most $a\varepsilon_R$, because the cap is one-Lipschitz in
the count. The bounded-variable exponential estimate then applies.
For every event $E$ measurable outside $B_R$, multiplying
\eqref{br:conditional-cap-tail} by $\mathbf1_E$ retains the factor
$\mathbb P(E)$, however small it is. The buffer is unobserved;
arbitrary pivotal conditioning is not covered. Negative association
is used for the dominating Bergman process, not assumed for every
conditioned Ginibre law.

The advanced analytic continuation of this argument is preserved as
recovered work.  It builds an almost surely finite Lipschitz field of
smoothing radii from these count tails, smears certificate atoms below
density one, and solves a singular weighted $\overline\partial$ problem.
The recorded conclusion is a simultaneous lower deficit
\begin{equation}
 1-P_A(x,x)\geq c_\delta h^{2+8\pi}e^{-C R(x)^2},
 \qquad h=\min\{\operatorname{dist}(x,G)/2,1\},
 \label{br:random-deficit}
\end{equation}
with a Gaussian one-point tail for $R(x)$.  A simpler product argument
in the archive says that deleting any three actual Ginibre points
leaves the exact zero set of a nonzero Fock function.  It uses the
Bufetov--Qiu Palm-product first moment and Tonelli; it does not assert
that deletion of one or two points suffices.  These claims concern
analytic nonuniqueness and unconditional moments, not a near-critical
pivotal comparison.

The probabilistic loss is exact even assuming
\eqref{br:random-deficit}.  If
$\mathbb P(R(x)^2>u)\leq Ke^{-cu}$ and $\mathbb P(E)=p>0$, then
\[
 \mathbb P(R(x)^2>u\mid E)
 \leq\min\{1,Ke^{-cu}/p\}.
\]
Integrating gives
\[
 \mathbb E[R(x)^2\mid E]\leq c^{-1}(\log(K/p)+1),
\]
and for $0<\eta<c$,
\[
 \mathbb E[e^{\eta R(x)^2}\mid E]
 \leq\frac c{c-\eta}(K/p)^{\eta/c}.
\]
This loss is sharp from a tail bound alone: take an exponential random
variable and condition on its upper tail.  Consequently unconditional
inverse moments of a deficit do not give a relative lower bound by
a fixed constant times the rare pivotal probability.

The same issue appears in the variance-sensitive potential estimates.
For a projection DPP statistic $S_f=\sum f(x)$, with oscillation $M$
and variance $V$, the recovered exponential-tilt calculation gives
\[
 \log\mathbb E e^{t(S_f-\mathbb ES_f)}
 \leq\frac V{M^2}(e^{tM}-1-tM),\qquad t\geq0.
\]
For $f_L(z)=\min\{L^{-4},|z|^{-4}\}$, the Ginibre mean is
$2\pi/L^2$ and the variance is at most $L^{-8}$.  Hence
\[
 \mathbb E[S_{f_L}\mid E]
 \leq\frac{2\pi}{L^2}
  +L^{-4}\bigl(\sqrt{2e\log(1/p)}+2\log(1/p)\bigr).
\]
Optimizing the Bennett expression improves the logarithmic profile
but does not eliminate its dependence on $p$.  The profile cannot
be replaced by a window-independent constant without geometric
information about $E$.

\paragraph{Thinning and a Poisson bath change the variance term.}
\label{br:mixed-bennett}
The recovered extension in \path{thinned_ginibre_potential_bennett.md}
applies to $X=\operatorname{DPP}(sH_1)\cup\Pi_\nu$, for a bounded
nonnegative integrable $f$ with $f\in L^2$. Put $M=\sup f$ and
\[
 V_*=s^2\Var_G\Bigl(\sum f\Bigr)+[s(1-s)+\nu]\int f^2.
\]
The same Bennett expression holds with mean $(s+\nu)\int f$ and
variance parameter $V_*$. A finite-rank proof embeds $sP_N$ as the
physical block of the projection
\[
 \begin{pmatrix}
 sP_N&\sqrt{s(1-s)}\,U_N\\
 \sqrt{s(1-s)}\,U_N^*&(1-s)I_N
 \end{pmatrix},
\]
where $U_N$ is the basis isometry onto $\operatorname{ran}P_N$.
The statistic is zero on the auxiliary atoms. Trace-norm exhaustion
and the independent Poisson MGF give the plane limit.
For $f_L=\min(L^{-4},|z|^{-4})$ one may use
\[
 V_*\le s^2L^{-8}+\frac{4\pi}{3}[s(1-s)+\nu]L^{-6}.
\]
Consequently an event of probability $p>0$ has conditional potential
mean at most
\[
 (s+\nu)\frac{2\pi}{L^2}
       +\sqrt{2eV_*\log(1/p)}+2L^{-4}\log(1/p).
\]
For root-indexed events in a deterministic region of area $A$ with
integrated probability $I>0$, entropy averaging gives the same
bracket with $\log(A/I)$, multiplied by $I$.
The $L^{-6}$ term is required away from the full projection without
a bath. This extension preserves the rare-event loss.

The corridor proposal aimed to supply that missing normalization.
If an exterior-support event $H$ contains the actual pivotal event
$E$, conditional domination may bound a local bad event by
\[
 \mathbb P(E\cap\mathrm{Bad})\leq\tau_D\mathbb P(H).
\]
The required next factor is $\mathbb P(H)\leq C_D\mathbb P(E)$.
Emptying a large disk costs on the quartic scale, too much to absorb
an area-exponential bad-density tail by that particular surgery.
A fixed-width corridor has compressed Ginibre norm
$\operatorname{erf}(w\sqrt{\pi/2})<1$. The buffered conditional
void theorem below makes the exponential-in-length cost precise.
But creating a pivotal-producing filling
with comparable conditional cost remained missing.  Upper domination
and hole bounds supply no insertion lower bound.

\subsection{Recovered conditional corridor and annular clearing laws}
\label{br:clearing}

The sources \path{research/rare_pivotal_weighting_and_corridors.md}
and \path{rare_pivotal_corridor_spectral_audit.md} establish a
conditional hole estimate under explicit buffer conditions.
If $D\subset B_M$ lies in a width-$w$ strip, write
$\kappa_w=\operatorname{erf}(w\sqrt{\pi/2})<1$.
Choose $R\ge10M$, $\pi R^2\ge5$, so that
$\varepsilon_R\le1$ and
$\varepsilon_R|D|\le(1-\kappa_w)/2$.
The Bergman approximation then gives
$\|B_R|_D\|\le(1+\kappa_w)/2$ and $\Tr(B_R|_D)\le2|D|$.
Using $\log(1-u)\ge-u/(1-\|B_R|_D\|)$ on its spectrum,
buffered conditional domination yields
\begin{equation}\label{br:corridor-hole}
 \mathbb P(N_K(D)=0\mid X|_{B_R^c})
       \ge e^{-4|D|/(1-\kappa_w)}
       \qquad(0\le K\le H_1).
\end{equation}
Fixed reduced-Palm remainders are included. A colored independent
bath contributes $e^{-\nu|D|}$. A finite collection of fixed-width
straight corridor pieces also has an exponential area cost under
the buffered spectral conditions in the source. These statements
give hole lower bounds, not occupied insertion lower bounds.

The count-preserving construction in
\path{count_preserving_corridor_clearing.md} is a different positive
result. Let $U$ be a rectangle of length $\ell$ and width $w$,
let $U_b$ expand all four sides by $b$, and put $V=U_{2\delta}$.
Move each unprotected point of $U$ to the exit position of planar
Brownian motion from $U_b$, with $b$ uniform in $[\delta,2\delta]$;
points in $V\setminus U$ stay fixed. The resulting Markov kernel
$T$ preserves count and lands outside $U$.
Its adjoint Lebesgue density is at most
$1+(w+4\delta)/(2\delta)$, by the rectangle torsion bound.

Given the Ginibre exterior and inside count, the positional density
is Gaussian times a squared holomorphic modulus. The exact gauge
\[
 e^{-\pi|z|^2}|F(z)|^2
       =e^{-2\pi(\Im z)^2}|e^{-\pi z^2/2}F(z)|^2
\]
and subharmonic stopping give a density cost per moved coordinate
of at most $A=\exp[2\pi(w/2+2\delta)^2]$.
Let
$C=A[1+(w+4\delta)/(2\delta)]$.
If almost every source configuration in $E$ is sent into
$F_{\rm tar}$ with probability one under the product move, then
\begin{align}
 \mathbb P(F_{\rm tar}\mid\mathrm{outside}\ V)
   &\ge C^{-n}\mathbb P(E\mid\mathrm{outside}\ V),\label{br:clearing-count}\\
 \mathbb P(F_{\rm tar})
   &\ge\mathbb E[C^{-N(V)}\mathbf1_E].\label{br:clearing-integrated}
\end{align}
Here $n$ is the conditioned count. The all-output geometric condition
is essential: preserving an event under deletion does not ensure its
preservation when the deleted points land in a surrounding layer.
Protected fixed Palm points enter the holomorphic factor; subsequent
canonical-certificate reintegration still needs multiplicity control.

The same construction applies to marked latent full Ginibre plus a
Poisson bath, moving all unprotected latent points and carrying their
retention marks. The count is the latent count plus bath count.
If $p=\mathbb P(E)$ and $M=(1+\nu)|V|$, its count MGF and Holder give
\[
 \mathbb P(F_{\rm tar})
       \ge p^{1+1/q}\exp[-M(C^q-1)/q]\qquad(q>0).
\]
For $|V|=O(t)$ and $\log(1/p)=O(t^2)$, a suitable $q$ proportional
to $\log t$ gives relative cost $\exp[-O(t^2/\log t)]$.
There is a stronger normalization for an event $H$ observed wholly
outside $B_t$ when $V$ lies at depth at least $d>0$ inside $B_t$:
the conditional intensity is at most
$b_d=(1-e^{-\pi d^2})^{-1}$. If clearing every configuration in $H$
has the target property, conditional Jensen gives
\[
 \mathbb P(F_{\rm tar}\cap H)
       \ge e^{-(b_d+\nu)|V|\log C}\mathbb P(H).
\]
No lower bound on $\mathbb P(H)$ is needed in this separated
conditioning regime.

For an annulus of middle radius $a\ge W:=w+4\delta$ and width $w$,
\path{count_preserving_annular_clearing.md} gives the same law with
\begin{equation}\label{br:annular-constant}
 C_{\rm ann}=e^{3\pi W^2/4}(1+2W/\delta).
\end{equation}
The harmonic gauge $2\pi a^2\log(|z|/a)$ leaves a Gaussian potential
of oscillation at most $3\pi W^2/4$, and annular torsion bounds the
adjoint exit density by $2W/\delta$. These constants do not depend
on $a$, while the landing region has area $2\pi aW$.
Its two landing layers are disconnected. If $R<w<2R$, all
attachments and the protected arm of one terminal lie inside the
ring, and those of the other lie outside, then every output keeps
them separate. An inserted root within range of both protected arms
is pivotal for every output. This is an eligible source geometry.

For a generic local bulk root both terminals may lie outside the
annulus; clearing a full ring then disconnects the root from both.
Ports, sufficient movable source vertices, and canonical multiplicity
remain separate tasks. The independent checks in
\path{count_preserving_clearing_audit.md} retain these restrictions.
Thus the clearing laws resolve a density-comparison step without
claiming a general pivotal-producing filling.

\subsection{Omitted points, adaptive choices, and multi-pair optimization}
\label{br:multipair}

The omitted-point route retained rather than discarded the event in a
Palm norm identity.  For distinct $q,p,t$, a normalized entire product
$F_Y(q)=1$ under the reduced $(q,p,t)$-Palm law gives a norm $N_q(Y)$
and the recovered exact formula
\begin{equation}
 \mathbb E^{qpt}[N_q\mathbf1_E]
 =\int W_{qpt}(z)\mathbb P^{zpt}(E)\,dz,
 \quad
 W_{qpt}(z)=
 \frac{|q-p|^2|q-t|^2\det H(z,p,t)}
      {\det H(q,p,t)|z-p|^2|z-t|^2}.
 \label{br:weighted-norm}
\end{equation}
The same function $E$ of the remainder appears on both sides; it is
not recomputed after moving its geometric root.  Determinantal zeros
cancel the apparent poles and the far tail is quartic.  Vandermonde
factors also cancel in the unconditioned norm mean as a fixed compact
triple coalesces.  These cancellations repair several earlier apparent
singular constants.

Decreasingness alone still does not compare
$\mathbb P^{zpt}(E)$ with $\mathbb P^{qpt}(E)$: their kernels differ
by one positive and one negative rank-one direction.  The notes also
give large-disk voids for which the relative norm mean can grow with
the disk's displacement.  Thus retaining the event is necessary but
not sufficient.

\paragraph{Steiner certificates and exact local capture.}
\label{br:steiner}
The deterministic result in
\path{research/steiner_certificate_clique_excess.md} explains how
local augmentation can retain every nearby pivotal alternative.
In a finite graph, let an inclusion-minimal connected induced
subgraph contain $k$ prescribed vertices. Except for the trivial
one-vertex case, deleting at most $k-2$ exceptional vertices makes
it triangle-free. To see this, every nonprescribed vertex is a
cutvertex. If $n_0$ is the number of noncutvertices, $c$ the number
of cutvertices and $b$ the number of blocks, the block-cut tree gives
\[
 \sum_{\rm blocks}(|B|-2)=n_0+c-b-1\le n_0-2\le k-2.
\]
Delete all but two vertices in each block of size at least three.
This deletion is for the counting bound; it need not preserve
connectivity or the prescribed vertices.

For two disconnected terminal components, prescribe every center in
$B(q,L_0)$ on either side and connect each prescribed set to its
virtual terminal by a minimal induced subgraph. Their union $A$ has
an exceptional set $E$ with
$|E|\le N_\eta(B(q,L_0)\cap W)$, and $A\setminus E$ has cap two
in every clique cell. Keeping the full $A$, and writing
$R=2r$, the exact capture identity is
\begin{equation}\label{br:local-capture}
 B_A\cap B(q,L_0-R)
       =\operatorname{Piv}(\eta)\cap B(q,L_0-R).
\end{equation}
A pivotal insertion in the smaller disk has all its old neighbors
in the prescribed local sets, with their paths retained; direct
terminal attachments are retained as well. Conversely $A\subset\eta$
and the old background fails, so a bridge for $A$ is an actual
pivotal insertion.
Canonical disintegration first fixes the local prescribed sets,
then orders intrinsic candidate subgraphs. A global order whose
validity changes when a new local point appears would not justify
the needed decreasing remainder. The quantitative Fock treatment of
the finite exceptional set is in
\path{augmented_certificate_fock_calculation.md} and its audit.

Adaptive crowded-cell arguments were developed to obtain actual omitted
marks.  A cap-two certificate omits a pair from every cell containing at
least four points.  If the certificate is augmented by at most $N$
exceptional local vertices, among its first $2N+1$ crowded cells at least
$N+1$ still omit a pair.  Averaging trial norms over those cells, then
using Campbell and the exact Palm product, cancels a whole Gram
determinant.  The recovered result is a common inverse-first-moment
bound polynomial in the local augmentation area, uniformly over the
admissible augmented certificates.  It is under the root Palm law
\emph{without} a rare pivotal conditioning.  Restoring a forced mark
without restoring it as a zero of the trial function would invalidate
this argument.

For $k\geq2$ actual omitted marks $P$, the exact relocation kernel is
\[
 W_{q,P}(z)=\frac{\det H(z,P)}{\det H(q,P)}
        \prod_{p\in P}\frac{|q-p|^2}{|z-p|^2}.
\]
It has tail $|z|^{-2k}$.  Campbell selection weights must be evaluated
on the restored configuration; for a uniform ordered choice from a
pool of size $m$ the weight is $(m)_k^{-1}$.  After reversing Campbell,
the pool and certificate are evaluated \emph{after deleting the
relocated point}, not under an unrelated ordinary configuration.
This dependence is part of the transport kernel.

Several pair-deletion trials for one common omitted set span only
\[
 F_Y\{\text{polynomials of degree at most }k-2\},
\]
a space of dimension $k-1$, not $\binom{k}{2}$.  With outside-norm
Gram matrix $G^B(Y)$, random optimizing coefficients are legitimate
only if retained inside the transported expectation.  They cannot be
replaced by an optimizer of the mean Gram matrix.  For example, if
$G$ is equally likely to be
$\operatorname{diag}(\epsilon,1)$ or
$\operatorname{diag}(1,\epsilon)$ and $c_1+c_2=1$, then
\[
 \mathbb E\min c^*Gc=\frac{\epsilon}{1+\epsilon},
 \qquad
 \min c^*(\mathbb EG)c=\frac{1+\epsilon}{4}.
\]
The difference is real, but its favorable size under pivotal weighting
is not supplied by this example.

The decisive variational limitation is deterministic.  Every admissible
trial vanishes on the same retained certificate $A$.  Its normalized
full Fock norm is at least $1/(1-P_A(q,q))\geq1$.  If $|B|<1$, its
outside norm is at least
\[
 \frac{1-|B|}{1-P_A(q,q)}.
\]
Thus no choice of many trial functions repairs the saturation examples
in \eqref{br:saturated-multiplier} pointwise.

The confluent benchmark in the later notes takes distinct marks to a
cluster with $k$ fixed, then lets $k\to\infty$, retaining fixed arm
zeros and an explicitly stated local event.  Its optimized limiting
cost is precisely $1/M_{A,B}(q)$, the original outside-norm optimum.
It can reach that multiplier, including its favorable examples, but
cannot turn an unfavorable certificate into a favorable one.  The
orders of the two limits and the full spatial-integral domination
are part of the recorded calculation.  Neither prescribed clusters
nor regular polygons can simply be assumed present with a useful
selection weight; sufficiently dense polygons may form an occupied
circuit and exclude bulk pivotality themselves.

\subsection{Explicit confluent means and a finite-Palm outside-norm gain}
\label{br:confluent-explicit}

The scalar calculations of \path{coalesced_multi_deletion_fock_mean.md}
and \path{multipair_fock_gram_optimization.md} give more than the
trial-space dimension. Fix a normalization root $q$ and let $k$
distinct omitted marks approach zero, first at fixed $k$.
For $u=\pi|q|^2$, define
\[
 Q_k(u)=1-e^{-u}\sum_{j=0}^{k-1}\frac{u^j}{j!},\qquad
 a_k(u)=\frac{u^k}{k!Q_k(u)},\qquad
 w_k(z)=k\int_0^1t^{k-1}e^{-\pi t|z|^2}\,dt .
\]
Removable values are taken continuously. The limiting scalar expected
density of the normalized common product is $a_k(u)w_k(z)$.
This is a limit of exact distinct-Palm expectations, not an assertion
of convergence of random products under a confluent law.
Tonelli gives the complete finite polynomial moments
\[
 \int |z|^{2j}w_k(z)\,dA
       =\frac{k\,j!}{\pi^j(k-j-1)},\qquad0\le j\le k-2;
\]
higher moments diverge. Put
$D_k(v)=\sum_{j=0}^{k-2}(k-1-j)v^j/j!$.
The deterministic optimal polynomial with value one at $q$ and
degree at most $k-2$ is
\[
 P_k(z)=\frac{D_k(\pi z\bar q)}{D_k(u)},\qquad
 \min_P a_k(u)\int|P|^2w_k=\frac{a_k(u)k}{D_k(u)}.
\]
These follow from the diagonal monomial Gram matrix. At $q=0$
the continuous value of the minimum is $k/(k-1)$.

As $k\to\infty$, the optimized densities converge pointwise to
$e^{-\pi|z-q|^2}$ and their total masses converge to one.
Scheffe's lemma upgrades this to $L^1$, so for every fixed Borel $B$,
\begin{equation}\label{br:confluent-gaussian}
 a_k(u)\int_{B^c}|P_k(z)|^2w_k(z)\,dA
        \longrightarrow\int_{B^c}e^{-\pi|z-q|^2}\,dA.
\end{equation}
For $B=B(q,b)$ this is $e^{-\pi b^2}<1$.
First select finite $k$ giving a strict gain, then choose sufficiently
close distinct marks; divided-kernel continuity and the integrable
$|z|^{-4}$ tail preserve it on an open set of tuples.
The omitted cluster can lie outside the desired bridge disk.
This is an unconditional finite-Palm improvement. High-order Palm
depletion empties every fixed local compact in the limit, so the
required arms are not automatically present. Retaining forced arms
and pivotal events leads back to the outside-norm optimum already
described, with its possible unfavorable sign.

\subsection{The complete outside-norm partition and its required coefficient}
\label{br:complete-partition}

The source \path{pivotal_outside_norm_campbell.md}, checked in
\path{neutral_janossy_outside_norm_audit.md}, makes a precise
distinction between complete certificate coverage and an eligible
pair in one selected cell. In this subsection set
\[
 I_F=\int\mathbb E_GD_qF\,dq,\qquad
 J_F=\int\mathbb E^{!q}D_qF\,dq;
\]
$J_F$ here is the full reduced-Palm pivotal mass, not the
leverage-weighted quantity of the small-beta section.
For the intrinsic canonical certificate $A$, let $E_A$ be its
selection and root-pivotal event in the reduced $A$-Palm remainder.
It is decreasing and contains its full bridge void. Conditional-void
Palm order gives
\[
 Q_A(q,q)\mathbb P^{!(A,q)}(E_A)
       \ge M_{A,B(A),1}(q)\mathbb P^{!A}(E_A).
\]
Divide by the positive multiplier and sum the full Campbell
partition. The determinant identity
$\det H(A\cup\{q\})=\det H(A)Q_A(q,q)$ and Tonelli yield
\begin{equation}\label{br:complete-outside}
 I_F\le S_F:=
 \int\mathbb E^{!q}\left[
   \frac{D_qF(X)}{M_{C_q(X),B(C_q(X)),1}(q)}\right]dq .
\end{equation}
No integrability of $S_F$ is presumed before a transport bound.
If one proved $S_F\le c_0I_F+c_1J_F$, with $c_1>0$ and
$c_0+c_1<1$, then
\[
 J_F-I_F\ge\frac{1-c_0-c_1}{c_1}I_F.
\]
This identifies a sufficient strict coefficient, not a proved
estimate. Uniform continuation and endpoint control would still
be needed for a critical-radius comparison.

The outside-norm transport removes redundant insertions in the same
certificate bridge. Its remaining redundancy is exactly
$z\in B_{\rm full}(Y)\setminus B(C_q(Y))$, with $q,z$ individually
successful in the same failing background. Ordinary contributions
still evaluate mark selection after deleting the relocated point.
Equation~\eqref{br:local-capture} removes this residual in a local ball,
but does not normalize the remaining rational potential to one.
The fixed-cell eligibility condition of at least two omitted points
is not decreasing and cannot be inserted into the order argument
for \eqref{br:complete-outside}. Adaptive complete selection must
retain its unbounded spatial scale. Moreover choosing marks beyond
an augmentation radius $L_0$ leaves the existing far bound with
coefficient $D^4/(L_0-R)^2=O(L_0^2)$ when $D$ is comparable to
$L_0$; enlargement alone does not make it small.

\subsection{Higher-mark angular gains and the annular obstruction}
\label{br:polygon}

The detailed refinements of
\path{multiple_omitted_marks_far_transport.md} and
\path{k_omitted_marks_radial_bennett_audit.md} retain geometry beyond
the tail exponent $2k$. For root zero and regular polygon marks
$p_j=D\omega^j$, $\omega=e^{2\pi i/k}$, the rational factor is
\[
 R(0,z;P)=\frac{D^{2k}}{|z^k-D^k|^2}.
\]
For $|z|\ge L>D$, it is at most
$M_{\rm poly}(L/|z|)^{2k}$, where
$M_{\rm poly}=(D/L)^{2k}/[1-(D/L)^k]^2$.
Extending the majorant inside by
$M_{\rm poly}(|z|/L)^{2k}$ gives mean
$2\pi kL^2M_{\rm poly}/(k^2-1)$, variance at most
$kM_{\rm poly}^2$, and far integral
$\pi L^2M_{\rm poly}/(k-1)$.
Marks within $\epsilon<L-D$ of their polygon slots incur only
\[
 C_{\rm slot}=(1+\epsilon/D)^{2k}
                       (1-\epsilon/(L-D))^{-2k}.
\]
Thus these far coefficients can decay exponentially in $k$ for
$D=cT$, $c<1$, and $L=T-R$. Actual admissible slot occupancy
and its selection weights remain necessary.

Conversely, if all marks lie outside $B_T$, any associated rational
function analytic inside $B_T$ and normalized to one at zero has
circle mean squared modulus at least one. Hence
\[
 \int_{T-R<|z|<T}|R_0(z)|^2\,dA
                      \ge\pi[T^2-(T-R)^2].
\]
Even retaining the Gram depletion does not remove the mass for
sparse polygons. For marks on radius $T\ge2R$ set
\[
 \alpha=e^{-\pi R^2/4},\quad
 \sigma=2\sum_{m\ge1}e^{-8\pi(T/k)^2m^2},\quad
 \epsilon_*=2\sum_{m\ge1}e^{-2\pi(T/k)^2(2m-1)^2}.
\]
If $\sigma,\epsilon_*\le(1-\alpha)/8$, on
$A=\{T-R<|z|<T-R/2\}$ the residual diagonal is at least
$c_R=(1-\alpha)/2$, the parent determinant is at least
$(1-\sigma)^k$, and exact angular averaging yields
\begin{equation}\label{br:polygon-annulus}
 \int_A\det H(z,P)R(0,z;P)\,dA
 \ge(1-\sigma)^kc_R\frac{\pi T(T-R)}{k}\log2 .
\end{equation}
For $k$ of order $\log T$ this has order $T^2/k$.
This is a kernel bound; the relocated pivotal expectation could
still suppress that annulus and must not be discarded.

The correlated polygon calculation is also explicit. Set
$a=\pi D^2$, $b=\pi r^2$,
$E_\ell(a)=\sum_{n\ge0}a^{\ell+nk}/(\ell+nk)!$, and
$c_{\ell n}(r)=(ab)^{(\ell+nk)/2}/(\ell+nk)!$.
For $r<D$ put $t=(r/D)^k$. Then
\begin{align}
 &\frac1{2\pi}\int_0^{2\pi}
       Q_P(re^{i\theta},re^{i\theta})R(0,re^{i\theta};P)\,d\theta
 \notag\\
 &\quad=\frac{1-e^{-b}\displaystyle\sum_{\ell=0}^{k-1}
         \frac{\sum_{n,m\ge0}c_{\ell n}c_{\ell m}t^{|n-m|}}
              {E_\ell(a)}}{1-t^2}.
 \label{br:polygon-correlated}
\end{align}
For $r>D$, use $t=(D/r)^k$ and multiply the right side by $t^2$.
At $r=D$ its removable value is
$\frac1{2k}\sum_\ell w_\ell\mathbb E|N_\ell-N_\ell'|$,
where $w_\ell=e^{-a}E_\ell(a)$ and $N_\ell,N_\ell'$ are independent
Poisson$(a)$ variables conditioned to be congruent to $\ell$ modulo
$k$. Residue-class coherent expansions and Fourier averaging prove
the formula; separate angular means cannot be multiplied.
For the normalized BQ density divide by
$Q_P(0,0)=1-1/E_0(a)$; for the Campbell-cancelled kernel multiply
by $\det H(P)$ instead.

Finally, very dense restored polygons may exclude the source event:
if $2D\sin(\pi/k)\le R=2r_0$ and $D\cos(\pi/k)>r_0$,
their occupied circuit surrounds the entire root grain. When both
terminal attachment sets lie outside this circuit, any two arms
reaching the root already meet the circuit, so the root is not
pivotal. The independent correlated-kernel audit preserves this
distinction between prescribed Palm marks and actual restored marks.

\subsection{What buffered Bergman domination does and does not supply}
\label{br:bergman-scope}

The conditional Bergman theorem proved elsewhere conditions outside a
larger disk $D$ while leaving the buffer between $D$ and the target $C$
unobserved.  It bounds the conditional process in $C$ above by the
weighted Bergman DPP, including prescribed Palm zeros.  Its unconditioned
disk compression has eigenvalues
$p_n(M)/p_n(R)$ for concentric radii $M<R$, and hence norm strictly
below one.  The limit to the full exterior is a conditional-expectation
martingale limit; it is not the assertion that an inverse on the
entire exterior remains a bounded operator.

Three limitations matter.  Observing immediately outside $C$ removes
the free buffer and permits rigidity to force its count.  Adding a
decreasing condition in the buffer is additional conditioning not
covered by the same theorem.  Finally, an upper conditional intensity
or a lower hole probability does not provide a lower probability of
inserting a prescribed scaffold.

Prescribed clustered zeros make local conditional occupancy very small:
the recorded bound is of the form $C(k+1)\rho^{2k}$ for a compact target
with relative radius $\rho<1$.  But comparing through a common exterior
law incurs a root normalization.  For a $k$-fold jet at the disk center,
the asymptotic factor recorded in the notes is
\[
 \frac{B_D^{\mathrm{jet}}(q,q)}{Q_k(q,q)}
 \sim
 \frac{e^{\pi R^2}(k+1)!}
      {(\pi R^2)^{k+1}(1-|q|^2/R^2)}.
\]
Its factorial growth can overwhelm the local hole gain.  Direct
fixed-Palm exterior comparison was therefore investigated separately;
the resulting unbounded likelihood ratios are a confirmed obstruction
to uniform arbitrary-exterior normalization, not a proof that the
actual pivotal average has the wrong sign.

The zero/one-neutral-point Janossy calculation illustrates the same
discipline.  With residual projection $Q$, a bounded retained region
$V$, and $L=Q_V(I-Q_V)^{-1}$, the zero-point root factor is $L(q,q)$;
with one neutral point $u$ it is
\[
 L(q,q)-\frac{|L(q,u)|^2}{L(u,u)}.
\]
The neutral set must exclude direct one-sided terminal contacts as well
as the bridge.  These factors can still vanish along saturated
certificates.  Few neutral points are not automatically a favorable
pivotal configuration.

\subsection{Reversible moves and genuine Goldman deletion weights}
\label{br:goldman}

A finite-rank reversible move was an intermediate repair for the
dependence problem.  Given a configuration $\eta$ and $y\in\eta$, let
$Q=H-P_{\eta-y}$.  Then
\[
 q_\eta^y(x)=\frac{|Q(x,y)|^2}{Q(y,y)}
\]
is a normalized move density because $Q$ is a projection.  Schur
factorization gives
\[
 \det H(\eta)q_\eta^y(x)
 =\det H(\eta-y)|Q(x,y)|^2,
\]
which is symmetric in $x,y$.  This establishes detailed balance of
the move for the thinned finite-rank DPP.  It does not establish
stationarity at intermediate positions on the straight segment from
$y$ to $x$, which the Poisson movement argument used.

The incoming density is a nonorthogonal dual-vector frame, not the
coherent projection diagonal.  For $n$ points coalescing near zero in
the rank-$n$ Fock space, that frame tends to
$n|\phi_{n-1}(x)|^2$, while their coherent span is the entire rank-$n$
space.  At the origin the ratio can vanish, and at distant points it
can grow with $n$.  The conditional squared jump distance tends to
$n/\pi$.  Thus a uniform conditional Gaussian jump tail or a
configuration-uniform comparison of frame and leverage is false.
Unweighted frame and leverage expectations agree, but a pivotal weight
introduces an uncontrolled covariance.  A process-level
infinite-Ginibre limit of this finite-rank move was not constructed.

The final representation uses the \emph{genuine} Goldman coupling,
rather than inventing an insertion rule from a convenient trial norm.
Disintegrating its deletion given $G$ gives weights $w_q(G,z)$ on actual
points with $\sum_{z\in G}w_q(G,z)=1$.  Its true conditional insertion
density from a reduced Palm remainder $X$ is
\[
 k_q(X,z)=w_q(X+z,z)\frac{d\mu^{!z}}{d\mu^{!q}}(X),
 \quad
 \int k_q(X,z)\\,dz=1,
 \quad
 \mathbb E^{!q}k_q(X,z)=e^{-\pi|z-q|^2}.
\]
The Gaussian formula is a marginal, not independence.  The exact
signed redundancy-minus-cooperation identity, and its thinning
extension, are proved in the principal-results section of this report.

Only an \emph{annealed} spatial column identity follows automatically:
\[
 \mathbb E^{!z}\int w_q(X+z,z)\,dq=1.
\]
It does not imply column mass one for every realized point.  The archive
contains a finite projection-DPP counterexample to that automatic
inference; it does not disprove the existence of a specially balanced
Ginibre coupling.  The rank-three failure of normalized Lagrange-square
weights, proved elsewhere in this report, rules out another convenient
candidate.  Thus the final unresolved sign is attached to actual
nonlinear deletion weights.  Removing those weights would change the
quantity being compared.

\subsection{Local Palm insertion and excluded analytic weight representations}
\label{br:palm-local}

Sections 6--8 of \path{ginibre_palm_relocation_markov.md}, with
\path{palm_relocation_local_independent_audit.md}, contain additional
recovered results. For $D=B(q,R)$, condition ordinary Ginibre on
its exterior $\xi$. Its inside law is a rank-$n$ polynomial-ensemble
projection $P_\xi$, with $n$ fixed by the exterior; its conditional
$q$-Palm law has rank $n-1$. The exterior density of the reduced
$q$-Palm law relative to the ordinary exterior is
\[
 t(\xi)=P_\xi(q,q),\qquad
 0\le t\le(1-e^{-\pi R^2})^{-1},\qquad \mathbb E t=1.
\]
Couple these two finite conditional projections by adding one point
inside $D$. Starting from the actual reduced $q$-Palm law reconstructs
the conditional ordinary inside law, but leaves its exterior tilted
by $t$. The resulting full law $\widetilde\mu$ therefore satisfies
\begin{equation}\label{br:local-insertion-tv}
 d_{\rm TV}(\widetilde\mu,\mu)
       =\mathbb E(t-1)_+\le e^{-\pi R^2}.
\end{equation}
The inequality follows from $(t-1)_+\le e^{-\pi R^2}t$.
The source audit supplies a measurable selection of the finite
couplings using their compact sections on configuration space.
This is an absolute local approximation, not relative control of
an arbitrarily rare pivotal event, and no limiting identification
with a particular Goldman coupling is claimed.

A different exact obstruction excludes holomorphic-square weights.
For an actual configuration $G$ with at least two points, suppose
entire functions $f_{G,z}$ vanish on $G\setminus\{z\}$ and obey
\[
 w_q(G,z)=e^{-\pi|q|^2}|f_{G,z}(q)|^2,\qquad
                   \sum_{z\in G}w_q(G,z)=1
\]
for almost every $q$. Submean bounds give local uniform convergence
of the covariance series. Analytic polarization forces
\[
 e^{\pi q\bar w}=\sum_{z\in G}
                     f_{G,z}(q)\overline{f_{G,z}(w)}.
\]
At two distinct actual nodes the right side vanishes termwise and
the left side is nonzero, a contradiction.

The nodal vanishing is itself forced for any genuine deletion family
continuous in $q$. The reduced-Palm intensity gives
\[
 \mathbb E^{!q}N(B(q,\epsilon))
       =\pi\epsilon^2-(1-e^{-\pi\epsilon^2})=O(\epsilon^4).
\]
Integrating over roots in an enlarged compact set, expressing the
Palm law by deletion, dividing by $\pi\epsilon^2$, and applying
Fatou to the nonnegative sum over actual pairs yields
$w_x(G,z)=0$ simultaneously for actual $x\ne z$.
This avoids differentiating at a random point.

It also excludes a positive-operator representation
$w_q(G,z)=\langle H(\cdot,q),A_{G,z}H(\cdot,q)\rangle$
with $A_{G,z}\ge0$ and $\sum_zA_{G,z}=I$ weakly.
The nodal zeros would imply
$A_{G,z}^{1/2}H(\cdot,x)=0$ for $x\ne z$; summing mixed inner
products would make two distinct coherent vectors orthogonal.
These impossibility statements concern the specified analytic and
positive-operator classes. They do not contradict general measurable
Goldman deletion weights. A root-dependent nonlinear normalization
leaves the analytic-square class and still needs both marginal laws
checked.

\subsection{Finite-rank reveal dynamics and total insertion influence}
\label{br:residual-dynamics}

The source \path{residual_projection_dynamics.md} gives finite-rank
identities separate from the reversible move. Let $H$ be a rank-$N$
projection with finite spatial moments. Reveal the points of
$\operatorname{DPP}(H)$ in independent uniform mark order, writing
$X_s$ for those already revealed. Conditional on $X_s=\eta$, the
unrevealed process has residual projection $Q_\eta=H-P_\eta$.
Its future birth intensity is $Q_\eta(z,z)/(1-s)$, and a birth at
$z$ changes the residual kernel by
\[
 Q_\eta^z(x,y)=Q_\eta(x,y)
       -\frac{Q_\eta(x,z)Q_\eta(z,y)}{Q_\eta(z,z)}.
\]
Zero-denominator locations have zero birth intensity.
With $F_s(x,y)=\mathbb E|Q_s(x,y)|^2$, the generator and projection
identity give
\[
 \partial_sF_s(x,y)=-\frac{2F_s(x,y)}{1-s}
 +\frac1{1-s}\int\mathbb E
       \frac{|Q_s(x,z)|^2|Q_s(z,y)|^2}{Q_s(z,z)}\,dz.
\]
Define $M_p(s)=\iint|x-y|^pF_s(x,y)\,dx\,dy$.
Subadditivity of $|\cdot|^p$ for $0<p\le1$ bounds the integrated
positive term by $2M_p(s)$; the squared triangle inequality bounds
it by $4M_2(s)$. Consequently
\begin{equation}\label{br:reveal-moments}
 M_p(s)\le M_p(0)\quad(0<p\le1),\qquad
 M_2(s)\le M_2(0)/(1-s)^2.
\end{equation}
For the rank-$N$ Bargmann projection at scale beta, the monomial
trace calculation gives $M_2(0)=2N\beta/\pi$ and
$M_1(0)\le N\sqrt{2\beta/\pi}$.

The thinned finite-rank law has Papangelou intensity
$c_s(x,\eta)=aQ_\eta(x,x)$, $a=s/(1-s)$. Adding $y$ changes it by
$a|Q_\eta(x,y)|^2/Q_\eta(y,y)$, whose integral in $x$ is exactly
$a$ whenever the denominator is positive.
A common-birth-proposal and common-death coupling for the associated
birth--death process therefore contracts expected Hamming
disagreements at rate at least $1-a$ for $s<1/2$.
This is a count estimate. Neither \eqref{br:reveal-moments} nor
the contraction permits multiplication by a rare pivotal indicator
without additional generator terms. No infinite-Ginibre process
limit is asserted by this recovered finite-rank calculation.

\subsection{Endpoint passages and separate analytic threshold bounds}
\label{br:endpoints}

Several local calculations survived after their intended large-scale
interpretation failed.  The replacement for an unproved rectangular
crossing step law is the full-plane seed-to-outer event with uncut
grains.  At fixed seed $m$, the events of reaching $\partial B_L$
decrease as $L\to\infty$.  Only finitely many components meet the seed,
so their intersection is exactly the event that an unbounded component
meets it.  Letting $m\to\infty$ and using ergodicity gives the
percolation zero--one law away from the critical radius.  A deterministic
slow diagonal and bounded capped thresholds then recover $r_c$ in
probability and mean.  No assertion at criticality is needed.  The
required comparison must nevertheless hold uniformly along that
diagonal; the diagonal does not create uniformity in a fixed-window
Taylor remainder.

Similarly, finite local thinning converges to the full projection
endpoint, and independent baths of intensity tending to zero disappear
locally.  These facts justify bounded-window endpoint limits when the
comparison itself has controlled coefficients.  They do not justify
integrating $(1-s)^{-1}$ errors or exchanging the $s\uparrow1$ and
$L\to\infty$ limits without a bound.  Number rigidity is one reason
that an arbitrary conditional argument need not extend uniformly to
$s=1$.

The separate-bounds route was also checked.  The published general
sub-Poisson upper bound used in the later audit is
$r_c(G)\leq\sqrt{2\log7}$.  A fully elementary induced-path exclusion
gives
\[
 r_c(\Pi_1)\geq
 \frac{1}{2\sqrt{\pi/3+\sqrt3/2}}.
\]
The intervals do not overlap in the useful direction.  The history
records sharper contour work, but it likewise does not establish an
upper Ginibre radius below a proved Poisson lower radius.

The threshold audit excluded Monte Carlo confidence intervals as
deterministic theorem inputs.  It also found an unjustified substitution
of an unconditional clustering coefficient for a history-selected
overlap in one proposed published bound.  For a parent at normalized
distance $h$, the actual overlap fraction is
\[
 c(h)=\frac{2\arccos(h/2)-(h/2)\sqrt{4-h^2}}{\pi},
\]
and at $h=1$ it differs by
$1/3-\sqrt3/(4\pi)>0$ from the asserted averaged value.
This rejects that supplied proof step, not necessarily the truth of
the numerical inequality it claimed.  No numerical experiment was
performed for the present investigation or report; the earlier
history's numerical explorations remain historical material only.

\subsection{Result-level map of the restored source discoveries}
\label{br:coverage-map}

The following table maps every finding of the coverage audit to its
mathematical treatment in this amended report. The associated source
filenames and scope are given in those passages; the complete source
index remains a separate archival inventory. The restored arguments
retain the verification distinctions stated at the start of the report.

\begin{longtable}{@{}L{0.08\textwidth}L{0.56\textwidth}L{0.27\textwidth}@{}}
\toprule
Audit & Restored result or refinement & Location\\
\midrule
\endhead
1 & Quantitative stretching-limit rate and its extra Poisson input
  & \ref{ell:recovered-rate}\\[3pt]
2 & Integrated critical-radius orientation, Riesz coefficient, and two-Palm response
  & \ref{ell:orientation}, \ref{ell:riesz-palm}\\[3pt]
3 & Switching atoms, rotated cusps, and smooth negative curvature
  & \ref{ell:switching}\\[3pt]
4 & Inclusion, displacement, overcrowding, and covering barriers
  & \ref{sb:coupling-barriers}\\[3pt]
5 & Strict positive-bath mixed-model threshold implication
  & \ref{sb:mixed-gap}\\[3pt]
6 & Buffered corridor holes, Brownian-exit clearing, and annular repair
  & \ref{br:clearing}\\[3pt]
7 & Entire fixed-kernel expansion and positive arms with neutral points retained
  & \ref{br:fixed-kernel}\\[3pt]
8 & Residual reveal dynamics and total Papangelou influence
  & \ref{br:residual-dynamics}\\[3pt]
9 & Protected internal close-pair estimate through the projection endpoint
  & \ref{br:close-pairs}\\[3pt]
10 & Local Palm reconstruction and excluded analytic/positive-operator weights
  & \ref{br:palm-local}\\[3pt]
11 & Steiner clique excess and exact local capture
  & \ref{br:steiner}, \eqref{br:local-capture}\\[3pt]
12 & Buffered conditional cap-density tail
  & \ref{br:conditional-density}, \eqref{br:conditional-cap-tail}\\[3pt]
Q1 & Complete outside-norm partition and sufficient strict coefficient
  & \ref{br:complete-partition}\\[3pt]
Q2 & Confluent polynomial moments, optimizer, and finite-Palm outside-norm gain
  & \ref{br:confluent-explicit}\\[3pt]
Q3 & Half-Gaussian contact split and dimer-perimeter identity
  & \ref{br:contact-refinements}\\[3pt]
Q4 & Balanced-mark tails and the correlated polygon/annular obstruction
  & \ref{br:polygon}\\[3pt]
Q5 & Thinning-plus-bath Bennett variance and entropy bound
  & \ref{br:mixed-bennett}\\[3pt]
Q6 & Fixed-count motion support, pushed-forward likelihood, and optimized rare-event cost
  & \ref{br:motions}, \eqref{br:motion-relative}\\
\bottomrule
\end{longtable}

\subsection{Source map and precise remaining boundary of the investigation}
\label{br:sourcemap}

The following map identifies the detailed working notes underlying this
section.  Their text is preserved in the source supplement, including
historical confidence language; that language is subject to the status
conventions at the beginning of this section.  A note ending in
\texttt{audit} records a collaborator's check at the time, not a proof
that all later uses of its conclusion were valid.

\begin{description}
\item[Earlier history.]
The attached research-history PDF, Sections 3--9, contains the initial
void/Voronoi, structure-factor, cluster, contact, dynamics, static-energy,
and square-replacement routes.  The supplied ellipse note is preserved
separately.  The extracted text is a reading aid, not a substitute for
the original mathematical typesetting.
\item[Affine and fixed-window signs.]
\path{isotropic_affine_counterexample.md};
\path{research/anisotropic_variation.md};
\path{ellipse_critical_variation.md};
\path{ellipse_minmax_switching.md};
\path{research/fixed_kernel_poisson_second_variation.md}.
\item[Cloning and contacts.]
\path{research/bottleneck_cloning.md};
\path{poisson_clone_decomposition.md};
\path{poisson_uniform_cloning_bound.md};
\path{mixed_clone_and_compensation.md};
\path{research/contact_certificate_disintegration.md};
\path{research/contact_scaffold_absorption.md};
\path{research/current_law_contact_obstructions.md}.
\item[Certificates and clearance.]
\path{replacement_count_lemma.md};
\path{conditional_count_tails.md};
\path{canonical_void_clearance.md};
\path{current_law_clearance.md};
\path{sparse_fock_projection.md};
\path{finite_window_leverage_operator.md};
the bulk, boundary and zero-bath separation notes and their audits.
\item[Bridge signs and saturation.]
\path{bridge_lens_resolvent.md};
\path{full_ginibre_transport_and_certificate_obstruction.md};
\path{two_parent_wide_lens_resolvent.md};
\path{straight_arm_wide_lens_resolvent.md};
\path{ginibre_triple_resolvent_certificate.md};
\path{pivotal_zero_one_neutral_janossy.md}.
\item[Motions and exterior support.]
\path{ginibre_motion_inverse_moments.md};
\path{ginibre_full_circle_cluster_motion.md};
\path{ginibre_triangle_cluster_motion.md};
\path{tapered_arm_area_motion_obstruction.md};
\path{count_preserving_corridor_clearing.md};
\path{corridor_filling_source_obstructions.md};
\path{ginibre_infinite_compact_factor_support.md};
\path{fixed_palm_infinite_exterior_ratio_obstruction.md}.
\item[Random deficits and rare-event normalization.]
\path{research/induced_certificate_density_gap.md};
\path{random_certificate_fock_nonuniqueness.md};
\path{ginibre_finite_deletion_fock_zero_sets.md};
\path{research/rare_pivotal_weighting_and_corridors.md};
\path{projection_dpp_variance_bennett.md};
\path{thinned_ginibre_potential_bennett.md}.
\item[Omitted points and Gram optimization.]
\path{three_palm_weighted_fock_identity.md};
\path{adaptive_crowded_cell_palm_fock.md};
\path{augmented_crowded_cell_palm_fock.md};
\path{pivotal_omitted_pair_campbell.md};
\path{multipair_fock_cross_moments.md};
\path{multiple_omitted_marks_far_transport.md};
\path{forced_arm_confluent_scalar_limit.md};
\path{retained_arms_multipair_variational_limit.md};
\path{research/multipair_projection_and_marking_obstructions.md}.
\item[Conditional kernels and genuine coupling.]
\path{exterior_palm_bergman_bounds.md};
\path{buffered_forced_marks_bergman_hole.md};
\path{reversible_projection_move.md};
\path{ginibre_palm_relocation_markov.md};
\path{thinned_goldman_relocation_identity.md};
\path{finite_ginibre_deletion_weights.md};
\path{research/goldman_column_mass_audit.md}.
\item[Limits, constants, and chronology.]
\path{annulus_threshold_convergence.md};
\path{analytic_critical_radius_bounds_audit.md};
\path{kong_yeh_clustering_audit.md};
\path{checkpoint_front_version2.md};
\path{continuation_front.md}.  The latter two record successive
states of the investigation and must not be read as simultaneous
certifications of every evolving claim.
\end{description}

The final mathematical boundary is specific.  The full Ginibre endpoint
has exact current-law interpolation and genuine deletion identities.
There are local bridge configurations of both signs, arbitrary-exterior
likelihood obstructions, and unconditional analytic estimates that lose
their constants under rare pivotal conditioning.  What was not obtained
is a signed, quantitatively positive comparison of the actual weighted
pivotal contributions, uniform over a sufficient interpolation and
growing windows, with controlled endpoints.  Consequently no strictly
positive radius gap for $G$ versus $\Pi_1$ follows from the investigation
as it stands.

% ===== END branches =====

% ===== BEGIN conclusion =====
\section{The remaining comparison and where strictness would enter}
\label{sec:remaining}

The original target is a strict inequality between two numbers at equal
intensity. It is weaker than monotonicity of every crossing probability,
every finite-cluster event, every affine deformation, or an entire
interpolation. Counterexamples to those stronger assertions eliminate
particular proofs, not the target itself.

\subsection{A finite-window route that would produce a radius gap}
The good-edge construction in the ellipse section supplies a universal
number \(b_0>0\): a stationary negatively associated process percolates
if, for one square scale \(L>2r\), both orientations of its local good-edge
event have failure probability less than \(b_0\). This is an
infinite-volume criterion proved by a dual-circuit sum, not an informal
inference from a high crossing probability.

Consequently, the following precise missing assertion is sufficient:
there is a finite scale \(L>2r_P\) such that
\[
 \max_{o\in\{\mathrm{horizontal},\mathrm{vertical}\}}
 \mathbb P_G(\text{the scale-}L\text{ edge of orientation }o
          \text{ is bad at }r_P)<b_0.
 \tag{S}
\]
If (S) were proved with a strict margin, it would persist at some
\(r_P-\delta>0\). To justify this last step, fix one slightly enlarged
bounded observation window containing every relevant center for radii
near \(r_P\). The point count there is almost surely finite. Conditional
on the count, its Janossy law is absolutely continuous. The finitely many
disk-intersection and side-contact incidences are stable in radius except
on tangency sets of Lebesgue measure zero. Equivalently one can describe
intersection of the finitely many clipped disks by finitely many
semialgebraic inequalities; the nontrivial equality cases have zero
probability. Thus each fixed block-event probability is continuous in
radius, by bounded convergence. The same fixed \(L\) and a small positive
\(\delta\) retain the criterion, giving
\[
                 r_c(G)\le r_P-\delta<r_P.
\]
Here the strict margin in (S), followed by finite-window continuity, is
the exact source of a positive radius gap. No proof of (S) is supplied by
this investigation.

\subsection{What the two recovered endpoint results do and do not say}
The verified limit
\[
          r_c(S_t^{-1}G)\longrightarrow r_P
\]
does not order the value at \(t=0\). This is a logical issue, not a lost
estimate: convergence alone allows approach from either side. The
large-area comparison holds for each fixed nonzero eccentricity beyond
a shape-dependent area threshold. It does not identify that threshold
with a neighborhood of a critical area, uniformly in eccentricity.
The joint hole regime makes the nonuniformity of the two limits explicit.

If the archived small-\(\beta\) geometric estimate is fully verified, its
squeeze proves
\[
 r_c(G_\beta)=r_P(1-\beta/4)+o(\beta)
\]
and hence a strict comparison for every sufficiently small positive
\(\beta\). The effective Poisson intensity exceeds one because
\(L_\beta=1+\beta/2+O(\beta^2)\), while its geometric loss is
\(o(\beta)\). That is a valid, explicit strictness mechanism at the
Poisson endpoint. It gives no continuation to \(\beta=1\) without
additional control. Even a right derivative at zero does not determine
a function's value at one. No monotonicity or uniform derivative estimate
on the whole interval has been established here.

\subsection{The event-weighted obstruction at full Ginibre}
The exact signed identity reduces the local interpolation derivative to
redundancy minus cooperation under a coupled reduced Palm law. Conditional
Fock and Bergman estimates retain substantial information, but the required
quantity weights that information by actual pivotality. An unconditional
inverse-moment bound, or a positive multiplier for a particular deterministic
certificate, does not control that weighted quantity. The report contains
explicit mechanisms preventing the replacement of the weighted estimate
by a uniform pointwise or arbitrary-exterior assertion.

The unresolved step is therefore not an omitted formal integration or a
Taylor expansion. It is a quantitative estimate retaining the pivotal
geometry and its probability, strong enough either to prove (S) or to
produce a positive intensity/radius gain uniformly over the growing
connection windows of a valid interpolation. No such full-Ginibre estimate
has been proved or disproved here.

\section{Published inputs and audit record}
\label{sec:inputs}

The following are external mathematical inputs, not results first proved
in this investigation. Normalizations and hypotheses are stated where
they enter the proofs.

\begin{enumerate}[leftmargin=*]
\item \textbf{A. Goldman, Ginibre and reduced-Palm coupling.}
\par
\href{https://arxiv.org/pdf/math/0610243}{Original paper, arXiv:math/0610243}.
Theorem 1 supplies the one-point coupling, in intensity-\(1/\pi\)
normalization. Scaling gives the kernels used here. The added point's
location is not assumed independent of the Palm remainder.

\item \textbf{R. Lyons, determinantal negative dependence and order.}
\href{https://arxiv.org/html/1406.2707v1}{Determinantal Probability:
Basic Properties and Conjectures, arXiv:1406.2707v1}.
Theorem 3.7 supplies negative association; Theorem 3.8 supplies stochastic
domination from operator order for locally trace-class positive
contractions on the same Radon-measure space. For a possibly irregular
observation set, its compressed kernels can be extended by zero to the
plane. Area-preserving changes of variables retain the hypotheses.

\item \textbf{D. Ahlberg, V. Tassion, A. Teixeira, planar Poisson crossings.}
\href{https://arxiv.org/html/1605.05926v1}{Sharpness of the phase
transition for continuum percolation in \(\mathbb R^2\),
arXiv:1605.05926v1}.
Theorem 1.1(i),(iii) supplies off-critical rectangle-crossing limits under
the radius-distribution second-moment hypothesis. Deterministic disks
satisfy that hypothesis. Their intensity formulation converts to a radius
formulation by spatial dilation, which also scales the rectangle.

\item \textbf{K. Adhikari, N. K. Reddy, Ginibre hole asymptotics.}
\href{https://arxiv.org/html/1604.08363v3}{Hole probabilities for finite
and infinite Ginibre ensembles, arXiv:1604.08363v3}.
Theorem 6 and Example 17 supply the ellipse coefficient
\((ab)^3/[2(a^2+b^2)]\) in their normalization. Ellipses satisfy the
required exterior-ball regularity. The report performs the intensity
and area changes explicitly.

\item \textbf{A. Bufetov, Y. Qiu, multiplicative functionals and Palm laws.}
\href{https://arxiv.org/html/1608.03736v2}{arXiv:1608.03736v2}.
The cited Proposition 1.1 and Corollary 1.6 concern equal-order Palm
equivalence and the radial conditional-density formula. The proof here
controls the radial analytic tail and handles fixed tuples explicitly;
it does not infer a uniform likelihood bound from equivalence.

\item \textbf{S. Ghosh, Y. Peres, Ginibre rigidity and tolerance.}
\href{https://arxiv.org/abs/1211.2381}{Rigidity and Tolerance in Point
Processes: Gaussian Zeros and Ginibre Eigenvalues,
arXiv:1211.2381}.
For Ginibre, the outside configuration determines the inside count,
and the conditional inside law is equivalent to Lebesgue measure on
that fixed-count configuration space. The argument does not impose the
additional sum rigidity of the distinct Gaussian-zero model.

\item \textbf{Strassen's closed-order coupling theorem.}
The local-to-global coupling extension uses the theorem for probability
measures on a Polish space with a closed partial order; see
\href{https://www.maths.tcd.ie/EMIS/journals/EJP-ECP/article/download/1005/1005-3375-1-PB.pdf}{Lindvall's account}
and its
\href{https://www.maths.tcd.ie/EMIS/journals/EJP-ECP/article/downloadSuppFile/1005/1005-8956-1-SP.pdf}{published correction}.
The critical-radius lower bound can also be obtained by the direct local
event exhaustion given in the proof, without this global coupling step.

\item \textbf{I. Manolescu, L. V. Santoro, near-critical Poisson crossings.}
\href{https://arxiv.org/html/2211.11661v1}{Widths of crossings in
Poisson Boolean percolation, arXiv:2211.11661v1}.
Theorem 2.4 and equation (2.16) are the additional inputs to the
recovered rate in Section~\ref{ell:recovered-rate}. The paper states
the near-critical result for fixed-radius disks and sketches its proof.
Only those inputs are used; radius conversion here uses
$\lambda=r^2$.
\end{enumerate}

\subsection{What was checked in preparing this report}
The six core statements were reread, with separate checks of the coupling
normalizations, conditional operator bounds, fixed-tuple exterior argument,
and finite-rank deletion counterexample. The ellipse critical-radius proof
received a separate reconstruction of its Gaussian estimates, support
colorings, planar gluing, source hypotheses, and order of limits. The
small-\(\beta\) archive was recovered and its dependency chain documented;
the entire geometric estimate was not newly certified by this review.
The additional branch survey distinguishes short reproduced calculations
from advanced historical derivations.

For the amendment, each coverage finding was matched to its original
derivation and audit where available; formulas, hypotheses, normalization,
and scope were restored in the relevant section. The additional
near-critical source statement was checked directly. This coverage and
consistency review does not promote all recovered arguments to the
six core statements' verification status.

These checks justify the statuses stated in this report. They do not turn
an unresolved geometric assertion into a theorem, and they do not establish
the original strict inequality. Compilation and typesetting checks of the
document are separate from these mathematical checks.

% ===== END conclusion =====

% ===== BEGIN source_index =====
\appendix
\section{Complete source index}
\label{app:source-index}
The companion archive contains the standalone report source and PDF, the two original research attachments, the consolidated historical checkpoint, and the individual notes indexed below. The originals are preserved byte for byte. Their titles and old audit labels are historical metadata, not fresh status assignments.


There are 130 individual mathematical Markdown notes. The file \path{manifest.json} records source paths, sizes, and SHA-256 digests. This is a preservation check, not a mathematical validation.
\small
\begin{longtable}{@{}p{0.48\textwidth}p{0.46\textwidth}@{}}
\toprule Source path & Original title\\\midrule\endhead
\path{active_fock_density_calculation.md} & Active Fock calculation after the capped-count theorem\\[4pt]
\path{adaptive_crowded_cell_independent_audit.md} & Independent audit of the adaptive crowded-cell Fock estimate\\[4pt]
\path{adaptive_crowded_cell_palm_fock.md} & Adaptive crowded cells and a uniform Palm first inverse Fock moment\\[4pt]
\path{analytic_critical_radius_bounds_audit.md} & Audit of explicit analytic critical-radius bounds\\[4pt]
\path{annulus_threshold_convergence.md} & A seed-to-boundary threshold that converges to the Boolean critical radius\\[4pt]
\path{audit_sources_findings.md} & Independent audit: primary inputs and finite-window counterexamples\\[4pt]
\path{augmented_certificate_fock_audit.md} & Audited quantitative Fock bound for augmented certificates\\[4pt]
\path{augmented_certificate_fock_calculation.md} & Augmented certificates: finite exceptions and local bridge capture\\[4pt]
\path{augmented_crowded_cell_palm_fock.md} & First inverse moments for augmented certificates under one-point Palm\\[4pt]
\path{bridge_lens_resolvent.md} & Bridge-lens resolvent: exact full-correlation calculation\\[4pt]
\path{bridge_lens_resolvent_audit.md} & Independent audit: bridge-lens resolvent\\[4pt]
\path{bridge_motion_continuation.md} & Fixed-count bridge motion: exact obstructions and averaged identities\\[4pt]
\path{buffered_bergman_domination_audit.md} & Independent audit of buffered Bergman domination\\[4pt]
\path{buffered_forced_bergman_independent_audit.md} & Independent audit of buffered holes with forced Palm marks\\[4pt]
\path{buffered_forced_marks_bergman_hole.md} & Buffered conditional holes retain prescribed Palm zeros\\[4pt]
\path{canonical_certificate_audit.md} & Canonical exterior certificates and the weighted Palm decomposition\\[4pt]
\path{canonical_void_clearance.md} & Canonical certificates, empty-ball conditioning, and residual leverage\\[4pt]
\path{checkpoint_front_version2.md} & Ginibre versus Poisson: continuation checkpoint\\[4pt]
\path{coalesced_multi_deletion_fock_mean.md} & Exact multi-deletion mean in a confluent Palm family\\[4pt]
\path{conditional_count_tails.md} & Exponential count bounds under deletion certificates\\[4pt]
\path{conditional_interior_audit.md} & Finite-rank conditional interior law and its planar probability limit\\[4pt]
\path{continuation_front.md} & Ginibre versus Poisson: active research checkpoint\\[4pt]
\path{cooperative_pair_boundary.md} & Boundary cooperative insertions: an explicit one-arm Poisson surgery\\[4pt]
\path{cooperative_pair_boundary_audit.md} & Independent audit of the boundary cooperative-insertion estimate\\[4pt]
\path{cooperative_pair_bulk.md} & A uniform bulk estimate for cooperative nearby Poisson insertions\\[4pt]
\path{cooperative_pair_dpp_bath.md} & Cooperative insertion estimates for a thinned projection DPP with a positive Poisson bath\\[4pt]
\path{corridor_filling_source_obstructions.md} & Source-count and exterior-weight obstructions to corridor filling\\[4pt]
\path{count_preserving_annular_clearing.md} & Count-preserving annular clearing and its exact geometric scope\\[4pt]
\path{count_preserving_clearing_audit.md} & Independent audit of Brownian-exit corridor clearing\\[4pt]
\path{count_preserving_corridor_clearing.md} & A count-preserving conditional Ginibre corridor-clearing comparison\\[4pt]
\path{current_law_clearance.md} & Averaged exterior clearance for the current interpolation law\\[4pt]
\path{current_law_clearance_proposal.md} & Proposed clearance estimate for current-law surgery --- unproved\\[4pt]
\path{direct_pivotal_comparison_route.md} & Direct insertion-pivotal route to the small-\(\beta\) pure endpoint\\[4pt]
\path{direct_pivotal_void_audit.md} & Independent audit: direct pivotal small-ball split\\[4pt]
\path{direct_pivotal_void_split.md} & Direct pivotal leverage bound by a small empty ball\\[4pt]
\path{ellipse_bridge.md} & Audited ellipse results and a critical-radius limit\\[4pt]
\path{ellipse_critical_variation.md} & Exact ellipse critical variation and the signed obstruction\\[4pt]
\path{ellipse_critical_variation_audit.md} & Independent audit: ellipse variation and quantitative Poisson limit\\[4pt]
\path{ellipse_minmax_switching.md} & Stretch derivatives of a finite bottleneck radius\\[4pt]
\path{exterior_palm_bergman_bounds.md} & Exterior Palm likelihood bounds and local Bergman domination\\[4pt]
\path{exterior_palm_bergman_independent_audit.md} & Independent audit: exterior Palm and Bergman bounds\\[4pt]
\path{finite_ginibre_deletion_weights.md} & Explicit finite Ginibre deletion weights and a rank-three test\\[4pt]
\path{finite_window_leverage_operator.md} & A positive operator version of local Palm leverage\\[4pt]
\path{fixed_kernel_poisson_second_variation_audit.md} & Independent audit of fixed-kernel second variation at Poisson\\[4pt]
\path{fixed_palm_exterior_transfer_independent_audit.md} & Independent audit of the fixed-Palm infinite-exterior transfer\\[4pt]
\path{fixed_palm_exterior_transfer_pending.md} & Fixed-Palm exterior transfer: superseded continuation record\\[4pt]
\path{fixed_palm_infinite_exterior_ratio_obstruction.md} & Fixed-Palm upgrade of the actual infinite-exterior ratio obstruction\\[4pt]
\path{fixed_region_zero_bath_audit.md} & Independent audit of the fixed-region zero-bath scaffold theorem\\[4pt]
\path{forced_arm_confluent_scalar_limit.md} & High-order deletion mean with fixed forced arms retained\\[4pt]
\path{full_ginibre_obstructions_audit.md} & Independent audits of two full-Ginibre obstructions\\[4pt]
\path{full_ginibre_transport_and_certificate_obstruction.md} & Full Ginibre transport and a saturated pivotal certificate\\[4pt]
\path{ginibre_compact_factor_support_audit.md} & Audit of infinite-Ginibre compact-factor realization\\[4pt]
\path{ginibre_finite_deletion_fock_zero_sets.md} & Finite deletion and Fock zero sets for an actual Ginibre configuration\\[4pt]
\path{ginibre_finite_deletion_fock_zero_sets_audit.md} & Independent audit of the three-deletion Fock zero-set argument\\[4pt]
\path{ginibre_full_circle_cluster_motion.md} & A full-circle Ginibre motion with a quantitative density gain\\[4pt]
\path{ginibre_full_circle_motion_independent_audit.md} & Independent audit: full-circle redundant-cluster motion\\[4pt]
\path{ginibre_infinite_compact_factor_support.md} & Actual infinite-Ginibre realization of a compact exterior factor\\[4pt]
\path{ginibre_motion_audit.md} & Ginibre two-parent motion: conditional ratios and the averaging gap\\[4pt]
\path{ginibre_motion_independent_audit.md} & Independent audit of Ginibre parent motion\\[4pt]
\path{ginibre_motion_inverse_moments.md} & Exact inverse-moment range for moving two Ginibre Palm parents\\[4pt]
\path{ginibre_palm_relocation_markov.md} & Genuine Ginibre Palm relocation density and its normalization\\[4pt]
\path{ginibre_quad_jensen_audit.md} & Independent audit: four-particle angular surgery\\[4pt]
\path{ginibre_triangle_cluster_motion.md} & Three-point version of the full-circle cluster repair\\[4pt]
\path{ginibre_triple_resolvent_certificate.md} & Exact uniform resolvent certificate for two triple clusters\\[4pt]
\path{induced_certificate_density_gap_audit.md} & Independent audit: induced-certificate density gap\\[4pt]
\path{isotropic_affine_counterexample.md} & Isotropy and ergodicity do not make disks minimize critical radius\\[4pt]
\path{k_omitted_marks_radial_bennett_audit.md} & k omitted marks: sharp radial tail and Bennett audit\\[4pt]
\path{kong_yeh_clustering_audit.md} & Independent audit of the Kong--Yeh clustering lower-bound argument\\[4pt]
\path{mixed_clone_and_compensation.md} & Current-law clone bound and differential compensation with a positive bath\\[4pt]
\path{mixed_cooperative_audit.md} & Independent audit of the positive-bath cooperative estimate\\[4pt]
\path{multipair_confluent_independent_audit.md} & Independent audit of multipair Gram identities and confluent scalar means\\[4pt]
\path{multipair_fock_cross_moments.md} & Polarized multi-pair Fock moments with exact Palm weights\\[4pt]
\path{multipair_fock_gram_optimization.md} & Multi-deletion trial spans and exact Gram optimization\\[4pt]
\path{multiple_omitted_marks_far_transport.md} & Multiple omitted Ginibre marks: exact transport and far-potential bounds\\[4pt]
\path{neutral_janossy_outside_norm_audit.md} & Independent audit: neutral Janossy states and the complete outside-norm partition\\[4pt]
\path{palm_continuation_checkpoint.md} & Palm branch continuation checkpoint\\[4pt]
\path{palm_current_law.md} & Current-law form of the Palm interpolation loss\\[4pt]
\path{palm_independent_audit.md} & Independent audit of current-law Palm and deletion estimates\\[4pt]
\path{palm_relocation_local_independent_audit.md} & Independent audit of local Palm relocation and measurable coupling\\[4pt]
\path{pivotal_close_pair_estimate.md} & A local protected close-pair estimate with its event weight retained\\[4pt]
\path{pivotal_omitted_pair_campbell.md} & Two omitted points: exact pivotal Campbell transport\\[4pt]
\path{pivotal_outside_norm_campbell.md} & Bridge-void Fock norm and exact alternative-certificate residual\\[4pt]
\path{pivotal_pair_bennett_independent_audit.md} & Independent audit: omitted-pair transport and projection Bennett bound\\[4pt]
\path{pivotal_zero_one_neutral_janossy.md} & Exact zero/one neutral-point law after forcing augmented arms\\[4pt]
\path{poisson_clone_decomposition.md} & Uniform Poisson movement bound and the remaining cloning term\\[4pt]
\path{poisson_cloning_bv_audit.md} & Independent audit: the Poisson cloning estimate at a fixed radius\\[4pt]
\path{poisson_uniform_cloning_bound.md} & Uniform Poisson cloning upper bounds\\[4pt]
\path{polygon_correlated_kernel_independent_audit.md} & Audit of the correlated polygon kernel and source geometry\\[4pt]
\path{positive_bath_pde_comparison.md} & From the positive-bath differential inequality to a strict mixed-model threshold gap\\[4pt]
\path{projection_dpp_variance_bennett.md} & A variance-sensitive exponential bound for projection DPP statistics\\[4pt]
\path{projection_dpp_variance_bennett_audit.md} & Independent audit of the projection-DPP Bennett estimate\\[4pt]
\path{pure_small_beta_theorem_audit.md} & Independent complete audit: strict small-beta threshold improvement\\[4pt]
\path{random_certificate_fock_nonuniqueness.md} & Fock nonuniqueness from a logarithmic-scale density gap\\[4pt]
\path{random_certificate_fock_nonuniqueness_audit.md} & Independent audit of the random-certificate Fock construction\\[4pt]
\path{rare_pivotal_corridor_spectral_audit.md} & Independent audit of strip and polygonal-corridor clearing estimates\\[4pt]
\path{replacement_count_lemma.md} & A deletion bound for crossing pivotals and the fine-mesh replacement limit\\[4pt]
\path{research/anisotropic_variation.md} & Anisotropic Ginibre expansion and opposite crossing signs\\[4pt]
\path{research/bottleneck_cloning.md} & Bottleneck thresholds: exact identities, counterexample, and Poisson cloning coefficient\\[4pt]
\path{research/contact_certificate_disintegration.md} & Canonical certificate disintegration of artificial-grain contact\\[4pt]
\path{research/contact_scaffold_absorption.md} & Contact-scaffold reduction of the zero-bath leverage term\\[4pt]
\path{research/current_law_contact_obstructions.md} & Artificial-grain contact: geometric facts and two ruled-out stronger bounds\\[4pt]
\path{research/fixed_kernel_poisson_second_variation.md} & Fixed-kernel expansion at Poisson and the combined two-parent geometry\\[4pt]
\path{research/ginibre_coupling_bridge_obstructions.md} & Exact obstructions to a thinning or bounded-displacement bridge\\[4pt]
\path{research/ginibre_exterior_analytic_product_audit.md} & The distant Ginibre analytic product tends to one\\[4pt]
\path{research/goldman_column_mass_audit.md} & Audit of genuine Palm insertion and the separate column-mass question\\[4pt]
\path{research/induced_certificate_density_gap.md} & A strict density gap for all induced two-arm certificates\\[4pt]
\path{research/multipair_projection_and_marking_obstructions.md} & Adversarial checks for optimized multi-pair deletion trials\\[4pt]
\path{research/rare_pivotal_weighting_and_corridors.md} & Rare pivotal weighting and a conditional corridor estimate\\[4pt]
\path{research/small_beta_critical_radius_asymptotic.md} & Exact Poisson domination and the first-order small-beta critical radius\\[4pt]
\path{research/steiner_certificate_clique_excess.md} & Minimal Steiner certificates: clique excess and capture of nearby relocations\\[4pt]
\path{residual_projection_dynamics.md} & Finite-rank residual projection dynamics\\[4pt]
\path{retained_arms_multipair_variational_limit.md} & Retained finite arms in the high-order multi-deletion limit\\[4pt]
\path{reversible_projection_move.md} & A reversible one-point move for a thinned finite-rank projection DPP\\[4pt]
\path{small_beta_asymptotic_audit.md} & Independent audit: exact domination and small-beta critical asymptotic\\[4pt]
\path{sparse_fock_projection.md} & Gaussian localization of a sparse coherent-state projection\\[4pt]
\path{sparse_fock_projection_audit.md} & Independent audit of sparse coherent-state localization\\[4pt]
\path{straight_arm_wide_lens_resolvent.md} & Positive wide-lens surplus with arbitrarily long straight forced arms\\[4pt]
\path{tapered_arm_area_motion_obstruction.md} & Tapered arm motion: the integrated lens-area obstruction\\[4pt]
\path{thinned_ginibre_potential_bennett.md} & Bennett control for a thinned Ginibre process with a Poisson bath\\[4pt]
\path{thinned_goldman_relocation_identity.md} & Thinned Goldman weights and the exact interpolation coefficient\\[4pt]
\path{thinned_relocation_independent_audit.md} & Independent audit of the thinned Goldman interpolation identity\\[4pt]
\path{three_palm_weighted_fock_identity.md} & Exact event-weighted three-Palm Fock identity and its limitation\\[4pt]
\path{two_parent_wide_lens_resolvent.md} & An explicit positive two-parent wide-lens multiplier\\[4pt]
\path{weak_hybrid_threshold_audit.md} & Independent audit of the strict weak-hybrid threshold comparison\\[4pt]
\path{wide_lens_independent_audit.md} & Independent audit of the positive wide-lens calculation\\[4pt]
\path{zero_bath_boundary_audit.md} & Independent audit: fixed-region boundary separation\\[4pt]
\path{zero_bath_boundary_separation.md} & Fixed-region boundary separation with no Poisson bath\\[4pt]
\path{zero_bath_bulk_audit.md} & Independent audit: fixed-region zero-bath bulk separation\\[4pt]
\path{zero_bath_bulk_separation.md} & Two-arm separation by deterministic holes and explicit disk geometry\\[4pt]
\path{zero_bath_certificate_route.md} & Canonical certificates and a proposed zero-bath clearance bound\\[4pt]
\bottomrule
\end{longtable}
\normalsize
\subsection*{Original attachments and consolidated record}
\begin{itemize}
\item \path{attachments/ginibre_poisson_attempts_summary.pdf}: the supplied research history, dated 5 September 2026.
\item \path{attachments/ginibre_ellipses_large_area.tex}: the supplied ellipse note, unchanged.
\item \path{history/ginibre_poisson_continuation_2026-09-15.md}: the consolidated historical checkpoint, including its full proof-note appendices.
\item \path{history/ginibre_proved_results.md}: the earlier selective account, retained to make the subsequent restoration of the ellipse theorem traceable.
\end{itemize}
The source archive intentionally retains assertions that the present report does not adopt as established. Consult the status discussion and the small-\(\beta\) section before using those historical claims.

% ===== END source_index =====

\end{document}
